\documentclass[reqno]{amsart}
\usepackage{amsmath,amssymb,amsthm,mathtools}
\usepackage[margin=0.92in]{geometry}
\usepackage{enumitem}
\setlist{nosep,leftmargin=*}
\usepackage[colorlinks=true,linkcolor=blue,citecolor=blue]{hyperref}

\theoremstyle{plain}
\newtheorem{theorem}{Theorem}[section]
\newtheorem{proposition}[theorem]{Proposition}
\newtheorem{lemma}[theorem]{Lemma}
\newtheorem{corollary}[theorem]{Corollary}
\theoremstyle{definition}
\newtheorem{definition}[theorem]{Definition}
\newtheorem{example}[theorem]{Example}
\newtheorem{remark}[theorem]{Remark}
\newtheorem{convention}[theorem]{Convention}

\DeclareMathOperator{\soc}{soc}
\DeclareMathOperator{\topp}{top}
\DeclareMathOperator{\Hom}{Hom}
\DeclareMathOperator{\Ext}{Ext}
\DeclareMathOperator{\im}{im}
\DeclareMathOperator{\End}{End}
\newcommand{\dimv}{\underline{\dim}\,}
\newcommand{\la}{\lambda}
\newcommand{\al}{\alpha}
\newcommand{\om}{\omega}
\newcommand{\ka}{\kappa}
\newcommand{\Pia}{\Pi}
\newcommand{\eps}{\varepsilon}
\newcommand{\hi}{\mathrm{hi}}
\newcommand{\lo}{\mathrm{lo}}
\newcommand{\ms}{\mathrm{m}}

\makeatletter
\def\@secnumfont{\bfseries}
\makeatother
\setlength{\topsep}{2pt plus 1pt minus 1pt}
\begin{document}

\title[Simple spectrum in type $D$]{Extremal torus fixed points in type $D$:\\
a classification of the simple-spectrum locus}

\author{Jaideep Sai Padhi}
\address{Department of Mathematics, Purdue University, West Lafayette, IN 47907}
\email{jpadhi@purdue.edu}

\begin{abstract}
Let $\Gamma=D_n$, let $\la=\om_k$ be a fundamental weight with $2\le k\le n-2$,
let $w$ be an element of the Weyl group and let $V$ be the unique $T$-fixed
representation on the Nakajima quiver variety $\mathfrak{M}(\la-w\la,\la)$.
Writing $\mu=w\la=\sum_{i\in S}\eta_ie_i$ in Bourbaki coordinates, put
$m=\#\{i\le n-2: i\in S,\ \eta_i=-1\}$ and $r=|S\cap\{n-1,n\}|$.  We prove,
over an arbitrary ground field, that $T$ acts on $V$ with simple
spectrum if and only if $r+2m\le2$.  For the two spin weights every $w$ is
$\la$-minuscule and the spectrum is always simple.  The proof identifies $V$
with an iterated socle of the injective $\Pia$-module $I_\ka$, decomposes
$I_\ka$ into two multiplicity-free ``strands'' under the leaf-exchanging
diagram automorphism, and reduces the statement to an explicit inequality
between two integer sequences at $9(k-1)$ rank-uniform weights.  The strand
decomposition requires $2$ to be
invertible; it is used only to construct one witness submodule in the generic
fibre, and a specialisation argument over $\mathbb{Z}_{(p)}$ carries the
conclusion to every characteristic, $p=2$ included.  We also record
an element of $D_4$ that is fully commutative but does not have simple
spectrum.  This answers, in type $D$ for non-spin fundamental $\la$, the open
problem of \cite[\S5.4.2]{DKK}; \S\ref{ss:remains} records what of that problem
remains.
\end{abstract}

\maketitle

\section{Introduction}\label{sec:intro}

\subsection{The problem}
Fix a simply-laced Dynkin diagram $\Gamma$ with vertex set $I$, Cartan matrix
$C=(a_{ij})$, simple roots $\al_i$, fundamental weights $\om_i$ and Weyl group
$W$.  Let $\mathfrak{M}(v,w)$ be the Nakajima quiver variety \cite{Nak} with dimension
vector $v$ and framing $w$, and write $\la=\sum_iw_i\om_i$ and
$\mu=\la-\sum_iv_i\al_i$.  When $\mu=w\la$ is extremal the variety is a single
point, and that point is a representation $V=\bigoplus_{i\in I}V_i$ of the
doubled quiver, acted on by a torus $T$.  Dinkins, Karpov and Krylov
\cite[\S5.4]{DKK} attach to such a $w$ a coloured graded poset, and their
formulas for the vertex function require the $T$-action on $V$ to have
\emph{simple spectrum}, i.e.\ pairwise distinct weights on each $V_i$.  They
ask \cite[\S5.4.2]{DKK} for a combinatorial characterisation of the $w$ for
which this holds.

The present paper answers that question in type $D_n$ for every non-spin
fundamental $\la$, in the sharpest possible form: the answer depends only on
the \emph{signs} of the Bourbaki coordinates of $\mu=w\la$, and not on $w$,
on any reduced word, or on the rank $n$.

\begin{theorem}[type $D$]\label{thm:typeD}
Let $\Gamma=D_n$ with the Bourbaki labelling $\al_i=e_i-e_{i+1}$ for
$i\le n-1$ and $\al_n=e_{n-1}+e_n$, let $t=n-2$ be the trivalent node, and let
$\la=\om_k$ with $2\le k\le n-2$.  Let the ground field $\mathbb{K}$ be
arbitrary.  Write
\[
   \mu=w\la=\sum_{i\in S}\eta_ie_i,\qquad |S|=k,\quad \eta_i=\pm1,
\]
and put
\[
   m=\#\{\,i\le n-2:\ i\in S,\ \eta_i=-1\,\},
   \qquad
   r=|S\cap\{n-1,n\}| .
\]
Then $T$ acts on $V$ with simple spectrum if and only if
\[
   \boxed{\,r+2m\le 2\,}
\]
--- equivalently, iff $m=0$, or $m=1$ and $r=0$.
\end{theorem}

\begin{theorem}[the spin weights]\label{thm:spin}
Let $\Gamma=D_n$ with $n\ge4$, let $\la=\om_\ka$ with $\ka\in\{1,n-1,n\}$, and
let the ground field $\mathbb{K}$ be arbitrary.  Then every $w$ of minimal
length in $wW_\la$ is $\la$-minuscule, and the module
$V=\soc_\sigma(I_\ka)$ of Theorem \ref{thm:main} satisfies
\[
   \dim_{\mathbb{K}}V_i[d]\le1\qquad\text{for every vertex } i
   \text{ and every degree } d,
\]
i.e.\ $V$ has simple spectrum in the sense fixed after Proposition
\ref{prop:grading}.  Over $\mathbb{C}$, where $V$ is the $T$-fixed
representation on $\mathfrak{M}(\la-w\la,\la)$ and the path-length grading is
the $T$-weight grading, this says that $T$ acts with simple spectrum, for
every $w$.  No hypothesis on the characteristic is needed, and the
characteristic-free half of the statement uses neither Proposition
\ref{prop:ST} nor Proposition \ref{prop:grading}.
\end{theorem}

The excluded case $k=1$ of Theorem \ref{thm:typeD} is contained in Theorem
\ref{thm:spin}; the criterion $r+2m\le2$ is in any case automatic there, since
$|S|=1$.

\subsection{What is proved, and over which fields}\label{ss:scope}
Theorem \ref{thm:typeD} holds over every field.  It is worth being precise
about where $\operatorname{char}\ne2$ does and does not appear, because the
role of that hypothesis is easy to misdescribe: it is \emph{not} a standing
hypothesis on the ambient field of the whole argument.  The proof of Theorem \ref{thm:typeD} funnels
through a single existential statement, $(B)$ of Definition \ref{def:B}: at
each of the $9(k-1)$ representatives $\mu$ of Proposition \ref{prop:reps},
\emph{some} submodule of $I_\ka$ has the greedy top truncation as its graded
profile.  Everything else --- the identification $V\cong\soc_\sigma(I_\ka)$,
the reduction to the trivalent node, the profile of $I_\ka$, the reduction of
$(*^\circ)$ to the representatives, the rigidity that turns ``some submodule''
into ``the submodule'', and the induction along $\sigma$ --- is proved directly
over an arbitrary field.  So $\operatorname{char}\ne2$ is a hypothesis on the
field over which one \emph{constructs the witness submodule} for $(B)$, not on
the field over which the theorem is true.  The construction is carried out once
and for all in the generic fibre, and it spreads by specialisation:
$\Pia$ is presented over $\mathbb{Z}$, the witness is saturated over
$\mathbb{Z}_{(p)}$, and its reduction mod $p$ is again a submodule with the
same graded dimensions.  That is Theorem \ref{thm:charp} of
\S\ref{sec:charp}, and it is what removes the hypothesis.

Explicitly:
\begin{itemize}\itemsep1pt
\item \emph{Proved directly in every characteristic}: the Weyl-group
combinatorics of \S\ref{sec:main}; Lemma \ref{lem:cb} and Theorem
\ref{thm:main}; Corollary \ref{cor:pol} (rigidity for submodules of $I_\ka$,
with the self-contained proof given there); Proposition \ref{prop:Pifacts},
whose part \textup{(a)} --- the $\Ext^1$ formula, which carries everything
else --- is \emph{proved} over an arbitrary field in \S\ref{ss:cbproof}
(Lemma \ref{lem:cbseq}, Proposition \ref{prop:ext1}) rather than quoted;
Proposition \ref{prop:dimI}; the inward injectivity Lemma \ref{lem:Id} and the
reduction to the trivalent node (Theorem \ref{thm:lemA}); Lemma \ref{lem:P}
(the profile of $I_\ka$), Lemma \ref{lem:R} (the recursion), Lemma
\ref{lem:C} ($v_t=r+2m$) and Lemma \ref{lem:nodefree} (the node column of a
top truncation, which is what the induction of \S\ref{sec:induction} reads at
the node, and which is deliberately independent of the
$\operatorname{char}\ne2$ Lemma \ref{lem:node}); \emph{the whole of
\S\ref{sec:reps}} --- Definition \ref{def:star}, Lemma \ref{lem:Jinv},
Corollary \ref{cor:reduce}, Proposition \ref{prop:reps}, Definition
\ref{def:B} and \eqref{eq:Brig}, i.e.\ the reduction of $(*^\circ)$ to the
representatives, together with \S\ref{ss:sharp}, namely Proposition
\ref{prop:sharp}, which is deduced from $(*^\circ)$ and Lemma
\ref{lem:Jinv}(b) with no hypothesis on the field --- which lies
\emph{outside} the standing hypothesis
$\operatorname{char}\ne2$ of \S\S\ref{sec:strands}--\ref{sec:T012} and
\S\S\ref{sec:families}--\ref{sec:closure} and is quoted at $p=2$; the
partial-slot closed form of Proposition
\ref{prop:PS}; Theorem \ref{thm:spin}; and the induction of
\S\ref{sec:induction}, given $(B)$.
\item \emph{Proved only for $\operatorname{char}\ne2$, and then transported}:
the strand decomposition and everything built on it --- Proposition
\ref{prop:W2} and (W1)--(W3), Propositions \ref{prop:W4out}, \ref{prop:W4l},
\ref{prop:W4n}, Corollary \ref{cor:W4}, Proposition \ref{prop:T012} and the
facts (T0)--(T2), the definition of the canonical truncation $T(v)$ at a
partially taken multiplicity-$2$ slot, and Theorem \ref{thm:B}, i.e.\ $(B)$
itself.  These are used at $p=2$ only through their consequence $(B)$, and
$(B)$ is transported to $p=2$ in \S\ref{sec:charp}.  None of them is claimed at
$p=2$; indeed Proposition \ref{prop:W2} is \emph{false} at $p=2$ in every
formulation we know (Remark \ref{rem:filtfail}).
\item \emph{Off the logical path}: the whole of \S\ref{sec:staircase}, which is
a word-free complement to Proposition \ref{prop:T} and is used by nothing else
in the paper, as that section says at its head.  Within it the characteristic
splits: Theorem \ref{thm:staircase} and Lemma \ref{lem:intervals} are proved
over an arbitrary field (Theorem \ref{thm:staircase} from Corollary
\ref{cor:pol} and Proposition \ref{prop:T}, the latter available at $p=2$ by
Theorem \ref{thm:charp}), whereas Lemma \ref{lem:U0}, Lemma \ref{lem:U} and
Proposition \ref{prop:tau} retain the hypothesis $\operatorname{char}\ne2$ and
are \emph{not transported at all}, the determinant argument in the proof of
Lemma \ref{lem:U} genuinely degenerating at $p=2$ (Remark
\ref{rem:sec9char}).  Theorem \ref{thm:typeD} does not care either way.
\end{itemize}

Section \ref{sec:verification} separates what is proved from what is merely
checked.  Its first part, \S\ref{ss:evidence}, records item by item everything
in this paper that is \emph{asserted} on the strength of a computation rather
than a proof.  That list is \emph{empty}: every assertion used in the proof of
Theorem \ref{thm:typeD} is proved here or cited; the sharpness of the side
condition in $(*^\circ)$, the last statement to have been offered on
computational evidence, is Proposition \ref{prop:sharp}.  The computations
formerly offered as evidence are retained in \S\ref{ss:checks}, which reports
what they found: counts of weights, of obligations and of slots, over the
ranges stated there.  Those counts are assertions of fact resting on
computation, and we do not claim otherwise; but each is a \emph{check} on a
statement proved elsewhere in this paper, not a step in any proof, and no proof
here cites \S\ref{sec:verification}.

\subsection{What remains of the problem of \texorpdfstring{\cite[\S5.4.2]{DKK}}{[DKK, 5.4.2]}}
\label{ss:remains}
The question of \cite[\S5.4.2]{DKK} is posed for a fixed point of an arbitrary
quiver variety, and the weight $\la$ of \cite[Def.~5.1]{DKK} is an arbitrary
weight; neither the Dynkin type nor the dominant weight $\la$ is restricted
there.  It is therefore worth saying plainly which part of that question the
present paper settles and which part it does not.

Settled here, or trivially:
\begin{itemize}\itemsep1pt
\item \emph{Type $A_n$, every fundamental $\la=\om_k$.}  Every such $\la$ is
minuscule, so every $w$ of minimal length in $wW_\la$ is $\la$-minuscule
\cite[\S2]{Stem} and the spectrum is always simple.  There is nothing to
characterise.
\item \emph{Type $D_n$, every fundamental $\la$.}  The non-spin weights are
Theorem \ref{thm:typeD}, the three minuscule weights $\om_1,\om_{n-1},\om_n$
Theorem \ref{thm:spin}.
\item \emph{The minuscule weights of $E_6$ and $E_7$}, namely $\om_1,\om_6$ and
$\om_7$ respectively, for the same reason as type $A$.
\end{itemize}

Not settled here:
\begin{itemize}\itemsep1pt
\item \emph{Type $E$ at a non-minuscule fundamental weight}: $\om_2,\om_3,\om_4,
\om_5$ in $E_6$, $\om_1,\dots,\om_6$ in $E_7$, and every fundamental weight of
$E_8$, which has no minuscule weight at all.  We know of no criterion in these
cases, and offer none.  What the present paper does supply there is the
machinery rather than the answer: Theorem \ref{thm:main}, which identifies $V$
with the iterated socle $\soc_\sigma(I_\ka)$, is proved for an arbitrary
simply-laced $\Gamma$ and is characteristic-free, so it applies verbatim in
type $E$.  It is the passage from that identification to a closed condition on
$\mu$ --- the analogue of $r+2m\le2$, whose proof here uses the two-leaf
symmetry of $D_n$ throughout \S\S\ref{sec:strands}--\ref{sec:closure} --- that
is missing.
\item \emph{$\la$ not fundamental}, in every type.  Every statement of this
paper assumes $\la=\om_\ka$, Theorem \ref{thm:main} included, and we make no
claim outside that case.  We note only that it need not be independent of the
fundamental case: when $\la=\la_1+\dots+\la_r$ with each $\la_a$ minuscule, the
extremal fixed points are indexed by the tuples
$(\mu_1,\dots,\mu_r)\in\prod_aW\la_a$ summing to $\mu$ and the relevant modules
factor accordingly \cite[Prop.~3.8]{DKK}, so each summand separately has simple
spectrum; but simplicity of the spectrum of the sum is a condition \emph{across}
the summands, and that condition is not addressed here.
\end{itemize}

\subsection{The criterion against the examples of \cite{DKK}}
The first thing to check is that the criterion reproduces the three worked
examples of \cite[\S5.4]{DKK}.  It does.  In each row below $\mu=w\la$ is
computed in Bourbaki coordinates, and the last column is the list of colours
occurring in each graded piece of $V$, read from the top degree down, as in the
Hasse diagrams of \cite[pp.~82--83]{DKK}.

\begin{center}
\renewcommand{\arraystretch}{1.25}
\begin{tabular}{llllcccl}
\hline
\cite{DKK} & $\Gamma$ & $\la$ & $w$ & $\mu$ & $(m,r)$ & $r+2m$ & rows of $V$\\
\hline
Ex.~5.10 & $D_4$ & $\om_2$ & $s_1s_2s_3s_4s_2$ & $(-1,1,0,0)$ & $(1,0)$ & $2$ &
$[2],[1,3,4],[2],[1]$\\
Ex.~5.11 & $D_5$ & $\om_2$ & $s_2s_1s_3s_4s_5s_3s_2$ & $(0,-1,1,0,0)$ & $(1,0)$ & $2$ &
$[2],[1,3],[2,4,5],[3],[2]$\\
Ex.~5.12 & $D_5$ & $\om_2$ & $s_2s_1s_2s_3s_4s_5s_3s_2$ & $(-1,0,1,0,0)$ & $(1,0)$ & $2$ &
$[2],[1,3],[2,4,5],[3],[2],[1]$\\
\hline
\end{tabular}
\end{center}

All three have $r+2m=2\le2$, so Theorem \ref{thm:typeD} predicts simple
spectrum; and indeed no colour is repeated inside any row, which by Proposition
\ref{prop:grading} is exactly simplicity of the spectrum.  The dimension
vectors agree as well: for Ex.~5.10 one gets $v=(2,2,1,1)$, and $v_t=v_2=2=r+2m$
in accordance with Lemma \ref{lem:C}.  Note that Ex.~5.12 is \emph{not} fully
commutative, so full commutativity is not necessary for simple spectrum; the
next example shows it is not sufficient either.

\begin{example}[full commutativity is neither necessary nor sufficient]
\label{ex:fcbad}
In $D_4$ take $w=s_2s_4s_3s_1s_2$, so that $\mu=(0,-1,1,0)$ and $v=(1,3,1,1)$.
Here $m=1$ and $r=1$, so $r+2m=3$ and Theorem \ref{thm:typeD} predicts failure.
Indeed the graded dimensions of $V\subseteq I_2$ are $\{2:1\}$ in degree $4$,
$\{1:1,3:1,4:1\}$ in degree $3$ and $\{2:2\}$ in degree $2$: the rows are
$[2]$, $[1,3,4]$, $[2,2]$, and the colour $2$ is repeated.  This $w$ \emph{is}
fully commutative.  So full commutativity does not imply simple spectrum;
and \cite[Ex.~5.12]{DKK}, whose $w=s_2s_1s_2s_3s_4s_5s_3s_2$ in $D_5$ is
recorded there as \emph{not} fully commutative yet does have simple spectrum,
shows the converse fails too.  Hence the condition fails in both directions.
\end{example}

\subsection{The shape of the proof}
The argument has five movements.

\smallskip\noindent
\emph{(1) From the quiver variety to a module.}  Let $\Pia$ be the
preprojective algebra of $\Gamma$ and $I_\ka$ the injective envelope of the
simple $S_\ka$.  Reading a reduced word for $w$ from the right and recording at
each step $c=\langle\al_i^\vee,\mu_{\mathrm{cur}}\rangle$ produces a sequence
$\sigma$ of colours with multiplicities, and $V\cong\soc_\sigma(I_\ka)$, the
iterated socle of Geiss--Leclerc--Schr\"oer taken along $\sigma$
(Theorem \ref{thm:main}).  Simple spectrum becomes
$\dim V_i[d]\le1$ for all $i$ and $d$, where $[d]$ is the path-length grading.

\smallskip\noindent
\emph{(2) One vertex suffices.}  Inward injectivity of the arrows of $I_\ka$
(Lemma \ref{lem:Id}) transports a repeat at any vertex to a repeat at the
trivalent node $t=n-2$.  So only $V_t$ needs to be inspected
(Theorem \ref{thm:lemA}).

\smallskip\noindent
\emph{(3) The profile at $t$, and the top truncation.}  The colour-$t$
multiplicity sequence of $I_\ka$ is $(1,2,\dots,2,1)$ with $k-1$ twos
(Lemma \ref{lem:P}), and $v_t=r+2m$ (Lemma \ref{lem:C}).  The content of the
theorem is that the colour-$t$ multiplicities of $V$ are the greedy \emph{top}
truncation of that sequence to total $v_t$ (Proposition \ref{prop:T}); then
every multiplicity is $\le1$ precisely when $v_t\le2$.

\smallskip\noindent
\emph{(4) Proposition \ref{prop:T} by induction along $\sigma$.}  The
hypothesis ``$X$ is a top truncation at $t$ and an up-set at each leaf'' is
carried along the socle construction.  Leaf steps are elementary
(\S\ref{ss:leafsteps}); colour-$t$ steps need Claim (G) of
\S\ref{ss:G}, which in turn needs the statement
\[
   (*^\circ)\qquad
   \text{every extremal } M\subseteq I_\ka \text{ with }
   \langle\al_t^\vee,\la-\dimv M\rangle\ge1 \text{ is an up-set at } t-1 .
\]

\smallskip\noindent
\emph{(5) $(*^\circ)$.}  This mentions no word and no order of construction,
and it is invariant under the parabolic $W_{J'}$, $J'=\{1,\dots,t-2\}$.  The
$W_{J'}$-minimal weights satisfying its side condition are $9(k-1)$ explicit
weights, \emph{independently of $n$} (Proposition \ref{prop:reps}).  At each of
them one exhibits a submodule of $I_\ka$ of the right dimension vector; by
rigidity (Corollary \ref{cor:pol}) that submodule \emph{is} $V$, and it is
visibly an up-set at every vertex.  The submodule is built from the two
``strands'' into which the leaf-exchanging automorphism $\tau$ splits $I_\ka$
(\S\ref{sec:strands}), and its closure reduces to two increment inequalities
between integer sequences (\S\ref{sec:closure}).

\subsection*{Acknowledgements}
I thank Vasily Krylov for the question that prompted this work and for
correspondence about \cite{DKK}.

\section{Setup}\label{sec:setup}

\subsection{The preprojective algebra}
Let $\Gamma$ be a simply-laced Dynkin diagram with vertex set $I$, and let
$\overline{Q}$ be the doubled quiver of an orientation $Q$ of $\Gamma$: for
each edge $\{i,j\}$ there are arrows $a_{i\to j}$ and $a_{j\to i}$.  The
preprojective algebra is
\[
   \Pia=\Pia_\Gamma=\mathbb{K}\overline{Q}\Big/
   \Big(\sum_{a\in Q}(aa^*-a^*a)\Big),
\]
over a ground field $\mathbb{K}$.  We fix once and for all the sign convention
$\varsigma(i,j)=\operatorname{sign}(j-i)$ on arrows, so that the relation at a
vertex $j$ reads $\sum_{i\sim j}\varsigma(j,i)\,a_{i\to j}a_{j\to i}=0$.

\emph{The orientation is fixed accordingly}: we take for $Q$ the orientation in
which every arrow points to the \emph{smaller} label, i.e.\ $a_{i\to j}\in Q$
exactly when $j<i$.  With that choice $\varsigma(a):=\varsigma(t(a),s(a))=+1$
precisely for the arrows $a$ of $Q$ and $-1$ for their reverses, which is what
makes the sign convention agree with the presentation
$\sum_{a\in Q}(aa^{*}-a^{*}a)$: cutting that element with $e_j$ on both sides
gives $+a\,a^{*}$ for each arrow $a\in Q$ with head $j$ and $-b\,b^{*}$ for each
reverse $b=a^{*}$ with head $j$ \textup{(}$a\in Q$ with tail $j$\textup{)},
i.e.\ the displayed relation, and \eqref{eq:relj} below.  \emph{For any other orientation the two would differ by
signs}, so the reader checking \eqref{eq:relj} should use this one.  (Nothing
depends on the choice: different orientations give isomorphic $\Pia$, by
rescaling the arrows.)

\emph{All modules in this paper are left $\Pia$-modules}, and all submodules
are $\mathbb{K}$-subspaces closed under every arrow.  (The involution of
$\Pia$ fixing the vertices and sending $a\mapsto\varsigma\,a^{*}$ identifies
$\Pia$ with $\Pia^{\mathrm{op}}$, so nothing is lost by this choice.)  We use
the following two standard facts about $\Pia$ in Dynkin type, and we state
both in full because they are used constantly; their provenance is
discussed in Remark \ref{rem:Pifacts} immediately below.

\begin{proposition}[standard]\label{prop:Pifacts}
Let $\Gamma$ be a simply-laced Dynkin diagram and $\Pia=\Pia_\Gamma$, over an
arbitrary ground field $\mathbb{K}$.
\begin{enumerate}
\item[\textup{(a)}] \textup{(the $\Ext^1$ formula; Crawley-Boevey
\cite[Lem.~1]{CB})}
For all finite-dimensional $\Pia$-modules $M,N$,
\[
   \dim\Ext^1_\Pia(M,N)=\dim\Hom(M,N)+\dim\Hom(N,M)-(\dimv M,\dimv N),
\]
where $(x,y)=x^{\mathsf T}Cy$.  \emph{This part is proved below}
\textup{(}Lemma \ref{lem:cbseq} and Proposition \ref{prop:ext1}\textup{)}
over an arbitrary field, and \cite{CB} is cited for the attribution only.
\item[\textup{(b)}] $\Pia$ is finite dimensional and self-injective.  Each
indecomposable projective $\Pia e_j$ is therefore also injective, and it has
simple socle $S_{\nu(j)}$, where $\nu$ is the Nakayama permutation of $\Pia$,
induced by the diagram automorphism $-w_0$.  Consequently the injective
envelope of $S_\ka$ is
\begin{equation}\label{eq:Ika}
   I_\ka\;\cong\;\Pia e_{\nu^{-1}(\ka)} .
\end{equation}
\end{enumerate}
\end{proposition}

\begin{remark}[provenance of Proposition \ref{prop:Pifacts}]\label{rem:Pifacts}
We are precise about attributions here, and about exactly how much of each
part the paper consumes.
\begin{itemize}\itemsep1pt
\item \textup{(a)} is Crawley-Boevey's $\Ext^1$ formula for modules over a
deformed preprojective algebra \cite[Lem.~1]{CB}, which reads: \emph{``If $M$
and $N$ are finite dimensional $\Pi(Q)$-modules, then
$\dim\Ext^1(M,N)=\dim\Hom(M,N)+\dim\Hom(N,M)-(\dim M,\dim N)$.''}  We do not,
however, rely on it as an authority.  The reason is a mismatch of hypotheses that matters for
this paper: \cite{CB} is a paper about Kleinian singularities and works
throughout over an algebraically closed field of characteristic $0$, whereas
part \textup{(a)} is the one input that carries \emph{everything} that is
characteristic-free here --- Theorem \ref{thm:main} through Lemma
\ref{lem:cb}, and the inequality of Corollary \ref{cor:pol} through which
\S\ref{sec:charp} runs.  The formula is nevertheless true over any field, for
any quiver, by the standard four-term sequence; we write that proof out in
Lemma \ref{lem:cbseq} and Proposition \ref{prop:ext1} below, so that
\textup{(a)} is a proved statement of this paper and \cite{CB} an attribution.
\item \textup{(b)} is sometimes attributed to Crawley-Boevey; it is not his.
Self-injectivity of $\Pia_\Gamma$ in Dynkin type, and the
identification of the Nakayama permutation with the diagram automorphism
induced by $-w_0$, are due to Ringel and Schofield (unpublished manuscript
\emph{Wild algebras with periodic Auslander--Reiten translate}); the first
published proof, with the Nakayama automorphism written out explicitly in each
$ADE$ case, is Brenner--Butler--King \cite{BBK}.  Neither
\cite[\S2]{BK} nor \cite{GLS} contains this statement.
\end{itemize}

Part \textup{(b)} is the only statement quoted rather than proved in this
paper, so we record precisely which of its assertions are consumed and which
are not.  Since $I_\ka$ is \emph{defined} as the injective envelope of
$S_\ka$ inside the finite-dimensional $\Pia$-modules, two of the assertions
that \textup{(b)} would supply cost nothing: $I_\ka$ is injective by
definition, and $\soc I_\ka=S_\ka$ because an injective envelope is an
essential extension, so every nonzero submodule meets $S_\ka$.  Neither
needs \textup{(b)}, and below we appeal to the definition rather than to
\cite{BBK} for them.  What \textup{(b)} is genuinely used for is the
following three items, and nothing else:
\begin{enumerate}\itemsep1pt
\item[\textup{(b1)}] $\Pia_\Gamma$ is finite dimensional, so that the
injective envelope exists inside the finite-dimensional modules and is itself
finite dimensional;
\item[\textup{(b2)}] $I_\ka\cong\Pia e_{\nu^{-1}(\ka)}$ as in \eqref{eq:Ika}.
This is what supplies the path-length grading used throughout, the fact that
$I_\ka$ is generated in degree $0$ \textup{(}fact (G1) in the proof of
Proposition \ref{prop:foldsol}\textup{)}, and the presentation $F/\mathcal{R}$
of $I_\ka$ from which Lemma \ref{lem:R} is proved;
\item[\textup{(b3)}] $\nu$ fixes every chain vertex of $D_n$, i.e.\
$I_k\cong\Pia e_k$ for $1\le k\le n-2$.  This enters once, as the initial
datum $\hat a_0=\delta_{\cdot k}$ of Proposition \ref{prop:foldsol}\textup{(a)}
\textup{(}together with fact (G3) there, which is the same statement read at
the top degree\textup{)}.
\end{enumerate}
What is \emph{not} consumed is the identification of $\nu$ with the
automorphism induced by $-w_0$ beyond \textup{(b3)}.  In particular the
explicit type-$D$ form of $\nu$ on the two leaves --- the identity for $n$
even, the leaf exchange for $n$ odd --- is never used: at a leaf weight the
paper carries $\nu^{-1}(\ka)$ as an unspecified leaf throughout, and
Proposition \ref{prop:foldsol}\textup{(b)} is stated in the folded
coordinates, in which the two leaves are indistinguishable by construction.
So the residue of the citation is: \emph{$\Pia_{D_n}$ is finite dimensional
and self-injective, and its Nakayama permutation fixes the chain vertices.}

Nor is any dimension count imported.  The type-$D$ profile and dimension of
$I_\ka$ are computed outright: Proposition \ref{prop:foldsol} determines the
whole graded profile of $I_\ka$ from Lemma \ref{lem:R} (the recursion),
Lemma \ref{lem:soctop} (the socle sits in the top degree), part
\textup{(a)}, and integer combinatorics, with no hypothesis on the ground
field; the resulting Corollary \ref{cor:dimI} is Proposition
\ref{prop:dimI}.  The numbers, the simplicity of $\soc\Pia e_j$ and the
vertex at which it sits are independently confirmed over $\mathbb{Q}$ and
over $\mathbb{F}_p$ for $p\in\{2,3,5,7\}$, for $D_4$ through $D_{10}$ and
every vertex $j$, by the computations of \S\ref{sec:verification}.
\end{remark}

\subsection{Proof of the \texorpdfstring{$\Ext^1$}{Ext1} formula over an
arbitrary field}\label{ss:cbproof}

We now prove Proposition \ref{prop:Pifacts}(a).  The argument is the standard
four-term sequence for $\Pia$; it is elementary, it is uniform in the ground
field, and --- as will be clear --- it uses neither the Dynkin hypothesis nor
any hypothesis on $\mathbb{K}$.  We record it because everything in this paper
that is asserted over an arbitrary field passes through part \textup{(a)}.

Fix the notation, used in this subsection only.  For an arrow
$a=a_{i\to j}$ of $\overline{Q}$ put $s(a)=i$ (source), $t(a)=j$ (target) and
$a^{*}=a_{j\to i}$, and write $M_a:M_i\to M_j$ for the action of $a$ on a
$\Pia$-module $M$.  (The symbol $t$ for the trivalent node of $D_n$ is
introduced only in Convention \ref{conv:D}, well after this subsection ends,
and $s(\cdot),t(\cdot)$ occur nowhere else.)  Write also, extending the sign
convention of \S\ref{sec:setup} to a one-argument form,
\[
   \varsigma(a):=\varsigma\bigl(t(a),s(a)\bigr)=\operatorname{sign}\bigl(s(a)-t(a)\bigr),
   \qquad\text{so that}\qquad \varsigma(a^{*})=-\varsigma(a),
\]
and the preprojective relation at a vertex $j$ reads
\begin{equation}\label{eq:relj}
   \sum_{a\in\overline{Q}:\ t(a)=j}\varsigma(a)\,M_aM_{a^{*}}=0 .
\end{equation}

\begin{lemma}[the four-term sequence]\label{lem:cbseq}
Let $\Gamma$ be any graph without loops, $\overline{Q}$ the doubled quiver of
an orientation, $\Pia$ the preprojective algebra over an arbitrary field
$\mathbb{K}$, and $M,N$ finite-dimensional $\Pia$-modules.  Put
\[
   C^{0}=\bigoplus_{i\in I}\Hom_{\mathbb{K}}(M_i,N_i),\quad
   C^{1}=\bigoplus_{a\in\overline{Q}}\Hom_{\mathbb{K}}(M_{s(a)},N_{t(a)}),\quad
   C^{2}=\bigoplus_{i\in I}\Hom_{\mathbb{K}}(M_i,N_i),
\]
and define $\al:C^0\to C^1$ and $\beta:C^1\to C^2$ by
\[
   \al(\varphi)_a=N_a\varphi_{s(a)}-\varphi_{t(a)}M_a,
   \qquad
   \beta(\psi)_j=\sum_{t(a)=j}\varsigma(a)\bigl(N_a\psi_{a^{*}}+\psi_aM_{a^{*}}\bigr).
\]
Then $\beta\al=0$, and
\begin{enumerate}
\item[\textup{(i)}] $\ker\al=\Hom_\Pia(M,N)$;
\item[\textup{(ii)}] $\ker\beta/\im\al\cong\Ext^1_\Pia(M,N)$;
\item[\textup{(iii)}] $\operatorname{coker}\beta\cong\Hom_\Pia(N,M)^{*}$.
\end{enumerate}
\end{lemma}

\begin{proof}
\emph{$\beta\al=0$.}  Substituting
$\al(\varphi)_{a^{*}}=N_{a^{*}}\varphi_j-\varphi_{s(a)}M_{a^{*}}$ and
$\al(\varphi)_a=N_a\varphi_{s(a)}-\varphi_jM_a$ (for $t(a)=j$) into $\beta$,
the two cross terms $\pm N_a\varphi_{s(a)}M_{a^{*}}$ cancel and what is left is
\[
   \beta\bigl(\al(\varphi)\bigr)_j
   =\Bigl(\sum_{t(a)=j}\varsigma(a)N_aN_{a^{*}}\Bigr)\varphi_j
   -\varphi_j\Bigl(\sum_{t(a)=j}\varsigma(a)M_aM_{a^{*}}\Bigr)=0
\]
by \eqref{eq:relj} for $N$ and for $M$.

\emph{(i).}  $\al(\varphi)=0$ says exactly $N_a\varphi_{s(a)}=\varphi_{t(a)}M_a$
for every arrow, i.e.\ that $\varphi$ is a morphism of $\Pia$-modules.

\emph{(ii).}  Let $\psi\in C^1$ and define $E(\psi)$ by
$E(\psi)_i=M_i\oplus N_i$ and
$E(\psi)_a(m,n)=\bigl(M_am,\ \psi_am+N_an\bigr)$.  Then
\[
   \sum_{t(a)=j}\varsigma(a)E(\psi)_aE(\psi)_{a^{*}}(m,n)
   =\bigl(0,\ \beta(\psi)_j\,m\bigr),
\]
again by \eqref{eq:relj} for $M$ and $N$; so $E(\psi)$ is a $\Pia$-module if
and only if $\beta(\psi)=0$, and then
$0\to N\to E(\psi)\to M\to0$ is an extension of $\Pia$-modules.  Conversely
every extension of $M$ by $N$ arises this way: choose a $\mathbb{K}$-linear
splitting at each vertex, which identifies the underlying graded vector space
with $M\oplus N$ and writes each arrow in the displayed lower-triangular form.
Finally, an isomorphism $E(\psi)\to E(\psi')$ inducing the identity on $N$ and
on $M$ is precisely a map $f$ with $f_i(m,n)=(m,n+\varphi_{i}m)$, and imposing
$f_{t(a)}E(\psi)_a=E(\psi')_af_{s(a)}$ gives
\[
   \psi_a+\varphi_{t(a)}M_a=\psi'_a+N_a\varphi_{s(a)},
   \qquad\text{i.e.}\qquad
   \psi'_a-\psi_a=\varphi_{t(a)}M_a-N_a\varphi_{s(a)}=-\al(\varphi)_a .
\]
The fibres of $\psi\mapsto[E(\psi)]$ are therefore the cosets of
$\im(-\al)=\im\al$ \textup{(}the sign is immaterial, $\al$ being linear, but we
record it correctly\textup{)}: that map is a surjection
$\ker\beta\to\Ext^1_\Pia(M,N)$ with
kernel $\im\al$, and it is additive because the Baer sum of $E(\psi)$ and
$E(\psi')$ is $E(\psi+\psi')$.

\emph{(iii).}  Use the trace pairing to identify $(C^2)^{*}$ with
$\bigoplus_i\Hom(N_i,M_i)$ and $(C^1)^{*}$ with
$\bigoplus_{a}\Hom(N_{t(a)},M_{s(a)})$.  For $\theta\in\bigoplus_i\Hom(N_i,M_i)$,
\[
   \langle\beta(\psi),\theta\rangle
   =\sum_{a\in\overline{Q}}\varsigma(a)\Bigl[
     \operatorname{tr}\bigl(\psi_{a^{*}}\theta_{t(a)}N_a\bigr)
     +\operatorname{tr}\bigl(\psi_aM_{a^{*}}\theta_{t(a)}\bigr)\Bigr].
\]
Re-indexing the first sum by $a\mapsto a^{*}$ and using $\varsigma(a^{*})=-\varsigma(a)$
and $t(a^{*})=s(a)$, it becomes
\[
   -\sum_{a\in\overline{Q}}\varsigma(a)\,
   \operatorname{tr}\bigl(\psi_a\theta_{s(a)}N_{a^{*}}\bigr),
\]
whence
\[
   \langle\beta(\psi),\theta\rangle
   =\sum_{a\in\overline{Q}}\varsigma(a)\,\operatorname{tr}\Bigl(
     \psi_a\bigl(M_{a^{*}}\theta_{t(a)}-\theta_{s(a)}N_{a^{*}}\bigr)\Bigr).
\]
This vanishes for every $\psi$ if and only if
$M_{a^{*}}\theta_{t(a)}=\theta_{s(a)}N_{a^{*}}$ for every $a$, i.e.\ (replacing
$a$ by $a^{*}$) $M_b\theta_{s(b)}=\theta_{t(b)}N_b$ for every arrow $b$, i.e.\
$\theta\in\Hom_\Pia(N,M)$.  So $(\operatorname{coker}\beta)^{*}=\ker\beta^{*}
=\Hom_\Pia(N,M)$.
\end{proof}

\begin{proposition}\label{prop:ext1}
Proposition \ref{prop:Pifacts}\textup{(a)} holds over an arbitrary field
$\mathbb{K}$: for all finite-dimensional $\Pia$-modules $M,N$, with
$m=\dimv M$ and $n=\dimv N$,
\[
   \dim\Ext^1_\Pia(M,N)=\dim\Hom(M,N)+\dim\Hom(N,M)-(m,n),
   \qquad (m,n)=m^{\mathsf T}Cn .
\]
\end{proposition}

\begin{proof}
Count dimensions in Lemma \ref{lem:cbseq}.  Writing $h=\dim\Hom_\Pia(M,N)$,
$h'=\dim\Hom_\Pia(N,M)$ and $e=\dim\Ext^1_\Pia(M,N)$, parts (i)--(iii) give
$\dim C^0=h+\dim\im\al$, $\dim\ker\beta=\dim\im\al+e$,
$\dim C^1=\dim\ker\beta+\dim\im\beta$ and $\dim C^2=\dim\im\beta+h'$, whence
\[
   \dim C^0-\dim C^1+\dim C^2=h+h'-e .
\]
On the other hand $\dim C^0=\dim C^2=\sum_im_in_i$ and, since the arrows of
$\overline{Q}$ are the ordered pairs $(i,j)$ with $i\sim j$,
$\dim C^1=\sum_{\{i,j\}\in\Gamma}(m_in_j+m_jn_i)$.  As
$C=2\,\mathrm{Id}-A$ with $A$ the adjacency matrix of $\Gamma$,
\[
   \dim C^0-\dim C^1+\dim C^2=2\sum_im_in_i-\sum_{i\sim j}m_in_j
   =m^{\mathsf T}Cn=(m,n).
\]
Comparing the two displays gives $e=h+h'-(m,n)$.
\end{proof}

\begin{remark}
Neither Lemma \ref{lem:cbseq} nor Proposition \ref{prop:ext1} uses that
$\Gamma$ is Dynkin, or that $\mathbb{K}$ is algebraically closed, or anything
about $\operatorname{char}\mathbb{K}$.  Part \textup{(b)} of Proposition
\ref{prop:Pifacts} does use the Dynkin hypothesis and remains an input,
cited to \cite{BBK}; it is used only to identify $I_\ka$ with a graded
$\Pia e_j$, never in a characteristic-sensitive way, and only to the extent
itemised as \textup{(b1)}--\textup{(b3)} in Remark \ref{rem:Pifacts}.
\end{remark}

We label \eqref{eq:Ika} as Proposition \ref{prop:dimI} for reference:

\begin{proposition}\label{prop:dimI}
$\dimv I_\ka=\om_\ka-w_0\om_\ka$ in root coordinates.  In type $D_n$ this
gives
\[
   \dim I_\ka=k(2n-1-k) \quad\text{for } \ka=k\le n-2,
   \qquad
   \dim I_\ka=\tfrac{n(n-1)}{2} \quad\text{for } \ka\in\{n-1,n\}.
\]
\end{proposition}

\begin{proof}
\emph{This proposition is proved outright, over an arbitrary field, in
Corollary \ref{cor:dimI} below}, as a consequence of Proposition
\ref{prop:foldsol}; nothing between here and there uses it in a proof --- the
sole reference, a parenthetical in the proof of Theorem \ref{thm:main}, is
flagged there as not used below --- so there is no circularity.  What follows is only the root-coordinate bookkeeping that
identifies the two displayed integers with the coordinate sum of
$\om_\ka-w_0\om_\ka$, i.e.\ that shows the two formulations of the
proposition to be the same statement.  Now $-w_0$ is
the identity on chain vertices of $D_n$ and either the identity or the leaf
exchange on $\{n-1,n\}$; in both cases $\om_\ka-w_0\om_\ka=\om_\ka+\om_{\ka^*}$
has the same coordinate sum as $2\om_\ka$.  Writing
$\om_k=e_1+\cdots+e_k=\sum_jc_j\al_j$ and comparing coefficients of $e_p$
gives $c_j=\min(j,k)$ for $j\le n-2$ and $c_{n-1}=c_n=k/2$, whence
\[
   \langle\om_k-w_0\om_k\rangle:=2\Bigl(\sum_{j\le n-2}\min(j,k)+k\Bigr)
   =k(k+1)+2k(n-1-k)=k(2n-1-k),
\]
where $\langle\,\cdot\,\rangle$ denotes the coordinate sum.  For $\ka=n$ one has $\om_n=\frac12(e_1+\cdots+e_n)$ and the same comparison
gives $c_j=j/2$ for $j\le n-2$, $c_{n-1}=(n-2)/4$, $c_n=n/4$; the coordinate
sum is $\frac{(n-2)(n-1)}{4}+\frac{n-2}{4}+\frac n4=\frac{n(n-1)}{4}$, and
doubling gives $n(n-1)/2$.  The case $\ka=n-1$ is the same by $\tau$.
\end{proof}

We emphasise that we work exclusively with \emph{submodules of the injective}
$I_\ka$ --- that is what Proposition \ref{prop:ST} below describes --- and
never with quotients of the projective $\Pia e_k$; the identification
\eqref{eq:Ika} is used only to know that $I_\ka$ has the path-length grading
of $\Pia e_{\nu^{-1}(\ka)}$.  (That $\soc I_\ka=S_\ka$ is simple needs no
identification at all: it is the definition of an injective envelope.  See
Remark \ref{rem:Pifacts}.)

For $\ka=k\le n-2$, a chain vertex, $\nu$ fixes
$\ka$ and $I_\ka\cong\Pia e_k$; that is the case of Theorem \ref{thm:typeD} and
of \S\S\ref{sec:strands}--\ref{sec:closure}.  The leaf weights $\ka\in\{n-1,n\}$
occur too, throughout \S\S\ref{sec:setup}--\ref{sec:main}: Lemma
\ref{lem:soctop}, Lemma \ref{lem:R} and Lemma \ref{lem:fold} are stated for an
arbitrary vertex $\ka$; Proposition
\ref{prop:dimI} and Corollary \ref{cor:dimI} state the leaf case explicitly,
and the latter is what Step 1 of Theorem \ref{thm:charp} uses at $p=2$; Lemma
\ref{lem:rigid}\textup{(b)} and Proposition \ref{prop:foldsol}\textup{(b)}
exist only for the leaf case; and Corollary \ref{cor:mfree} and Theorem
\ref{thm:spin} are about it.  There $\nu^{-1}(\ka)$ may be the other leaf, so
below we write $I_\ka\cong\Pia e_{\nu^{-1}(\ka)}$ and keep track of which vertex
carries degree $0$.  We never need to decide which leaf it is --- for the
record, $\nu$ is the identity when $n$ is even and the leaf exchange when $n$
is odd, but that refinement is not used anywhere, and Proposition
\ref{prop:foldsol}\textup{(b)} is stated in folded coordinates precisely
because it does not distinguish the two leaves.

Every $\Pia e_j$ is graded by path length, \emph{and this is the grading used
throughout}.  We write $M_i[d]$ for the degree-$d$ part of a graded module $M$
at vertex $i$.

\begin{lemma}[the socle is the top degree]\label{lem:soctop}
Let $d^{*}$ be the largest degree with $(I_\ka)[d^{*}]\ne0$.  Then
$\soc(I_\ka)=(I_\ka)[d^{*}]$; it is one dimensional and sits at the vertex
$\ka$.  In particular $\soc(I_\ka)=S_\ka$ is concentrated in the top degree,
and the path-length grading is the socle-layer grading read backwards.
\end{lemma}

\begin{proof}
Every arrow raises the path-length degree by $1$, so every element of
$(I_\ka)[d^{*}]$ is killed by every arrow and hence lies in $\soc(I_\ka)$;
thus $0\ne(I_\ka)[d^{*}]\subseteq\soc(I_\ka)$.  The socle is the simple
module $S_\ka$, of dimension $1$ and supported at the vertex $\ka$: this is
not an appeal to Proposition \ref{prop:Pifacts}(b) but the definition of
$I_\ka$, an injective envelope being an essential extension of $S_\ka$, so
that every nonzero submodule of $I_\ka$ meets $S_\ka$ and hence
$\soc(I_\ka)=S_\ka$.  (Only the path-length grading, which is part
\textup{(b2)} of Remark \ref{rem:Pifacts}, is quoted here.)  A
one-dimensional space containing a
nonzero subspace equals it, so $\soc(I_\ka)=(I_\ka)[d^{*}]=\mathbb{K}$, at
vertex $\ka$.
\end{proof}

\begin{lemma}[the recursion]\label{lem:R}
Let $\Gamma$ be simply-laced Dynkin, $\ka$ any vertex, and
$a_d(i)=\dim(I_\ka)_i[d]$.  Then for every $i$ and every $d$ strictly below
the top degree $d^{*}$,
\[
   a_{d+1}(i)=\sum_{j\sim i}a_d(j)-a_{d-1}(i),\qquad
   a_0(i)=\delta_{i,\nu^{-1}(\ka)},\quad a_{-1}=0 .
\]
More precisely, the map
$Y:(I_\ka)_i[d-1]\to\bigoplus_{j\sim i}(I_\ka)_j[d]$, $u\mapsto(a_{i\to j}u)_j$,
is injective for every $i$ and every $d-1$ strictly below $d^{*}$.
\end{lemma}

\begin{proof}
Write $c=\nu^{-1}(\ka)$, so that $I_\ka=F/\mathcal{R}$ with
$F=\mathbb{K}\overline{Q}e_c$ the free module and $\mathcal{R}$ the submodule
spanned by the elements $p\,r_l\,q$, where $l$ runs over the vertices,
$r_l=\sum_{j\sim l}\varsigma(l,j)\,a_{j\to l}a_{l\to j}$ is the preprojective
relation at $l$, and $p$, $q$ are paths $l\to i$ and $c\to l$ respectively.
All of $F$, $\mathcal{R}$ and $I_\ka$ are graded by path length.

Every path $c\to i$ of length $d+1$ factors uniquely as $a_{j\to i}q'$ with
$j\sim i$ and $q'$ a path $c\to j$ of length $d$, so
\[
   \Phi\colon\bigoplus_{j\sim i}F_j[d]\;\xrightarrow{\ \sim\ }\;F_i[d+1],
   \qquad (u_j)_{j\sim i}\longmapsto\sum_{j\sim i}a_{j\to i}u_j ,
\]
is an isomorphism.  We compute $\Phi^{-1}\bigl(\mathcal{R}_i[d+1]\bigr)$.  A
spanning element $p\,r_l\,q$ of $\mathcal{R}_i[d+1]$ has
$|p|+|q|=d-1$, where $|\cdot|$ denotes path length, and there are two cases.
If $|p|\ge1$, write
$p=a_{j\to i}p'$ with $j\sim i$; then $p\,r_l\,q=\Phi$ of the tuple whose
$j$-entry is $p'r_lq\in\mathcal{R}_j[d]$ and whose other entries vanish.  If
$|p|=0$ then $p=e_i$, hence $l=i$ and $q$ is a path $c\to i$ of length
$d-1$, and
\[
   r_i\,q=\sum_{j\sim i}\varsigma(i,j)\,a_{j\to i}\bigl(a_{i\to j}q\bigr)
   =\Phi\Bigl(\bigl(\varsigma(i,j)\,a_{i\to j}q\bigr)_{j\sim i}\Bigr).
\]
So $\Phi^{-1}(\mathcal{R}_i[d+1])=\bigoplus_{j\sim i}\mathcal{R}_j[d]
+\{(\varsigma(i,j)a_{i\to j}q)_j:q\in F_i[d-1]\}$: the first case contributes
only elements that are \emph{already zero one degree down}, in
$(I_\ka)_j[d]$, and the second contributes the image of $Y$ up to the signs
$\varsigma(i,j)$, which rescale the coordinates by $\pm1$ and so change neither the dimension
of the image nor the injectivity of $Y$.  Passing to the
quotient, the induced surjection
\[
   \overline{\Phi}\colon\bigoplus_{j\sim i}(I_\ka)_j[d]
   \longrightarrow(I_\ka)_i[d+1]
\]
has kernel exactly $\im Y$.  This is the point at which ``the preprojective
relations impose one condition for each basis vector of $(I_\ka)_i[d-1]$, and
no others'' is proved rather than assumed.  Hence
$a_{d+1}(i)=\sum_{j\sim i}a_d(j)-\dim\im Y$, and the stated recursion is
equivalent to the injectivity of $Y$.

That injectivity holds below the top degree: a nonzero element of $\ker Y$ is
killed by every arrow out of $i$, hence spans a copy of $S_i$ inside
$\soc(I_\ka)$, which by Lemma \ref{lem:soctop} is a one-dimensional space in
the top degree $d^{*}$.  So $\ker Y=0$ whenever $d-1<d^{*}$.
\end{proof}

This lemma and Lemma \ref{lem:soctop} are used without further comment
throughout.

\subsection{Iterated socles}
For $j\in I$ and a $\Pia$-module $M$ let $\soc_{(j)}(M)$ be the largest
submodule of $M$ that is a direct sum of copies of $S_j$.  Given a submodule
$X\subseteq M$, the \emph{socle step of colour $j$} replaces $X$ by the
preimage in $M$ of $\soc_{(j)}(M/X)$; concretely,
\[
   X\ \longmapsto\ \{\,x\in M:\ a_{j\to i}\,x\in X\ \text{for every } i\sim j\,\}
   \ \cap\ (X+M_j),
\]
which enlarges $X$ only at vertex $j$.  For a sequence $\sigma$ of colours,
$\soc_\sigma(M)$ denotes the result of applying the corresponding steps to
$X=0$ in order.  This is the iterated socle
$\soc_{(j_1,\dots,j_\ell)}$ of Geiss--Leclerc--Schr\"oer \cite[\S2.4]{GLS}.

\subsection{The fixed point, and why it is a submodule}
We use two inputs, which we state in full rather than by number, and which we
keep strictly apart: a \emph{geometric} identification, which is available over
$\mathbb{C}$, and an \emph{algebraic} rigidity property, which needs no
hypothesis on the ground field at all.  The separation matters: the
characteristic-$p$ argument of \S\ref{sec:charp} uses only the second.

\begin{proposition}[Savage--Tingley {\cite[Thm.~4.4, Prop.~4.9]{ST}}]
\label{prop:ST}
Work over $\mathbb{C}$, and let $\mathfrak{L}(v,w)$ be the Lagrangian
subvariety of $\mathfrak{M}(v,w)$.  Then $\mathfrak{L}(v,w)$ is in natural
bijection with the quiver Grassmannian of submodules of dimension vector $v$
inside $\bigoplus_iI_i^{\,w_i}$; the bijection is a bijective algebraic map,
hence a homeomorphism.  Moreover the following are equivalent: $v$ is
$w$-extremal; $\mathfrak{L}(v,w)$ is a single point; that quiver Grassmannian
is a single point; $\bigoplus_iI_i^{\,w_i}$ has a unique submodule of graded
dimension $v$.
\end{proposition}

\begin{remark}\label{rem:STconv}
Two points of bookkeeping.  First, \cite{ST} states Theorem 4.4 as a bijective
algebraic map and a homeomorphism, not as an isomorphism of varieties, and we
have quoted it in that form; nothing below uses more.  Second, the quiver
Grassmannian $\mathrm{Gr}_P(v,q_w)$ of \cite{ST} is \emph{defined} as a space
of submodules of the injective module $q_w$, which is the convention of this
paper (see the paragraph after Proposition \ref{prop:dimI}); no translation
between submodules of an injective and quotients of a projective is needed
anywhere.
Finally, \cite{ST} works throughout over $\mathbb{C}$, and Proposition
\ref{prop:ST} is stated here with that hypothesis.
\end{remark}

\begin{proposition}[rigidity; Baumann--Kamnitzer--Tingley {\cite[\S3]{BKT}}]
\label{prop:rig}
Let $T$ be an object of an essentially small abelian category in which every
object has finite length, and let $\mathrm{Pol}(T)$ be its
Harder--Narasimhan polytope.  If $x$ is a vertex of $\mathrm{Pol}(T)$, then
$T$ has a unique subobject $X$ with $[X]=x$.
\end{proposition}

\begin{remark}[attribution, and the absence of a field hypothesis]
\label{rem:rigidity}
This is stated in \cite[\S3]{BKT} as an unnumbered consequence of Corollary
3.3 there: \emph{``Another noteworthy consequence is the following rigidity
property: if $x$ is a vertex of $\mathrm{Pol}(T)$, then $T$ has a unique
subobject $X$ such that $[X]=x$.''}  Two points deserve emphasis, because both
are easy to get wrong.

\emph{(i) The result is in \cite{BKT}, not in \cite{BK}.}  What
\cite[Rem.~4.5]{BK} proves is the inclusion $\dimv U\in\mathrm{Pol}(T)$ for
every submodule $U\subseteq T$; both the convex-hull description of
$\mathrm{Pol}(T)$ and the description of the submodules realising its extremal
points are there explicitly \emph{deferred} to a paper in preparation of
Baumann, Kamnitzer and Sadanand.  So \cite{BKT} is the load-bearing reference
and \cite{BK} the incidental one.

\emph{(ii) There is no hypothesis on the ground field.}  Section 3 of
\cite{BKT} opens \emph{``We consider an essentially small abelian category
$\mathcal{A}$ such that all objects have finite length''}, and nothing else is
assumed there; the convention \emph{``We fix a base field $K$, which we assume
for convenience to be algebraically closed of characteristic $0$''} appears
only in \cite[\S4.1]{BKT}, where preprojective algebras enter, and does not
govern \S3.  Over any field $\mathbb{K}$ the algebra $\Pia_\Gamma$ of a Dynkin
$\Gamma$ is finite dimensional, so its category of finite-dimensional modules
satisfies the hypothesis of Proposition \ref{prop:rig}.  This is what makes
\S\ref{sec:charp} possible.

For the one case this paper actually needs --- submodules of $I_\ka$ with
$\Gamma$ Dynkin --- we do not even have to rely on \cite{BKT}: Corollary
\ref{cor:pol} below is given a short self-contained proof, valid over every
field.
\end{remark}

\begin{corollary}\label{cor:pol}
Let $\Gamma$ be simply-laced Dynkin, $\mathbb{K}$ \emph{any} field, and
$U\subseteq I_\ka$ a submodule with $v=\dimv U$.  Then
\[
   \delta(v):=2v_\ka-v^{\mathsf T}Cv\ \ge\ 0 ,
\]
and $\delta(v)=(\la,\la)-(\mu,\mu)$ where $\mu=\la-\sum_iv_i\al_i$; in
particular $\delta(v)=0$ whenever $v$ is $w$-extremal.  If $\delta(v)=0$ then
$I_\ka$ has \emph{at most one} submodule of dimension vector $v$.
Consequently, if $v=\la-w\la$ then $I_\ka$ has at most one submodule of
dimension vector $v$; over $\mathbb{C}$ there is exactly one, and it is the
$T$-fixed representation $V$.
\end{corollary}

\begin{proof}
Everything here is characteristic-free.

\emph{The identity.}  With $\la=\om_\ka$ and $\mu=\la-\sum_iv_i\al_i$ one has
$(\la,\la)-(\mu,\mu)=2(\la,\sum_iv_i\al_i)-(\sum_iv_i\al_i,\sum_iv_i\al_i)
=2v_\ka-v^{\mathsf T}Cv$, since
$(\om_\ka,\al_i)=\delta_{i\ka}$ in the normalisation $(\al_i,\al_i)=2$.  If
$\mu=w\la$ then $(\mu,\mu)=(\la,\la)$, so $\delta(v)=0$.

\emph{The inequality.}  Since $I_\ka$ is injective with
$\soc I_\ka=S_\ka$ --- both by the definition of an injective envelope, not
by Proposition \ref{prop:Pifacts}(b) --- the functor
$\Hom_\Pia(-,I_\ka)$ is exact and $\dim\Hom(S_i,I_\ka)=\delta_{i\ka}$, so
counting along a composition series of $U$ gives $\dim\Hom(U,I_\ka)=v_\ka$.
Composition with the inclusion $U\hookrightarrow I_\ka$ embeds
$\End_\Pia(U)$ into $\Hom(U,I_\ka)$, whence $\dim\End(U)\le v_\ka$.  By
Proposition \ref{prop:Pifacts}(a) with $M=N=U$,
\[
   v^{\mathsf T}Cv=(v,v)=2\dim\End(U)-\dim\Ext^1_\Pia(U,U)\ \le\ 2v_\ka ,
\]
which is $\delta(v)\ge0$.

\emph{Uniqueness.}  Let $U\ne U'$ be submodules of $I_\ka$ with
$\dimv U=\dimv U'=v$ and $\delta(v)=0$.  Both $U\cap U'$ and $U+U'$ are
submodules, and $\dimv(U\cap U')+\dimv(U+U')=2v$; write
$\dimv(U\cap U')=v-e$ and $\dimv(U+U')=v+e$.  Since $\delta$ is the quadratic
polynomial $x\mapsto2x_\ka-(x,x)$,
\[
   \delta(v-e)+\delta(v+e)=2\delta(v)-2(e,e)=-2(e,e).
\]
The Cartan matrix of a Dynkin diagram is positive definite, so $(e,e)>0$ unless
$e=0$; and $e\ne0$ would make one of $\delta(v\pm e)$ negative, contradicting
the inequality just proved.  So $e=0$, i.e.\ $U\cap U'$ has dimension vector
$v$, and therefore $U\cap U'=U=U'$.

\emph{Existence over $\mathbb{C}$}, and the identification with $V$, is
Proposition \ref{prop:ST}: for $v$ $w$-extremal that proposition gives a
unique submodule of $I_\ka$ of dimension vector $v$ and a single point of
$\mathfrak{L}(v,\om_\ka)$, which is the $T$-fixed representation $V$.
(Over an arbitrary field existence is not quoted at all: Theorem
\ref{thm:main} constructs such a submodule directly, by an iterated socle.)
\end{proof}

\begin{remark}
Corollary \ref{cor:pol} is exactly the rigidity property of Proposition
\ref{prop:rig} in the case at hand, and we could have quoted \cite[\S3]{BKT}
for it, that section carrying no hypothesis on the field (Remark
\ref{rem:rigidity}).  We prefer the proof above because it is three lines long,
because it makes the absence of a characteristic hypothesis visible rather than
inherited, and because \S\ref{sec:charp} rests on precisely this point.
\end{remark}

We refer to a submodule $U$ with $\dimv U$ extremal as an \emph{extremal
module}, and write $V(\mu)$ for the unique submodule with
$\dimv V(\mu)=\la-\mu$.  Over an arbitrary field its \emph{existence} is not an
input: Theorem \ref{thm:main} below constructs a submodule of $I_\ka$ of
dimension vector $\la-w\la$ by an iterated socle, over any field and with no
appeal to geometry.  So $V(\mu)$ is well defined over every $\mathbb{K}$, and
over $\mathbb{C}$ it is the $T$-fixed point of $\mathfrak{M}(\la-\mu,\la)$.

The last sentence of Corollary \ref{cor:pol} is what we use constantly: to
identify $V(\mu)$ it suffices to \emph{exhibit any submodule} of $I_\ka$ of the
correct dimension vector.

\subsection{The two gradings agree}
\begin{lemma}[the bidegree is determined]\label{lem:bigrade}
$\Pia$ is bigraded by the numbers of $a$- and $a^*$-arrows separately, the
relations being bihomogeneous of bidegree $(1,1)$.  For $\Gamma$ a tree, the
piece $(\Pia e_j)_i[d]$ is pure of a single bidegree, namely
$p=\frac12(d+f(i)-f(j))$, $q=\frac12(d-f(i)+f(j))$, where $f$ is any function
on vertices with $f(\mathrm{head})-f(\mathrm{tail})=1$ on every arrow of $Q$.
\end{lemma}

\begin{proof}
Such an $f$ exists because $\Gamma$ is a tree.  An $a$-arrow raises $f$ by $1$
and its reverse lowers it by $1$, so along any path from $j$ to $i$ one has
$p-q=f(i)-f(j)$; with $p+q=d$ this determines $p,q$.
\end{proof}

\begin{proposition}\label{prop:grading}
The $T$-weight decomposition of $V_i$ coincides with its path-length
decomposition, up to an affine normalisation depending only on $i$.  In
particular
\[
   T \text{ acts on every } V_i \text{ with simple spectrum}
   \iff
   \dim V_i[d]\le 1 \ \text{ for all } i\in I \text{ and all } d .
\]
\end{proposition}

\begin{proof}
Since $\la=\om_\ka$ is fundamental, the framing space is one dimensional and
the framing torus acts on all of $V$ through one character, contributing only a
global shift.  The rest of $T$ acts through the bigrading of $\Pia$, assigning
to $(\Pia e_{\ka'})_i[d]$ a weight $\alpha p+\beta q$ for fixed integers
$\alpha,\beta$.  By Lemma \ref{lem:bigrade} both $p$ and $q$ are affine in $d$
of slope $\pm\frac12$, so $\alpha p+\beta q$ is affine in $d$ of slope
$\frac12(\alpha+\beta)$, injective in $d$ unless $\alpha+\beta=0$.  The latter
is excluded because the symplectic form has nonzero weight --- that being what
makes the $T$-fixed locus isolated.  So the weight determines and is determined
by $d$.  Finally $V$ is a \emph{graded} submodule: by Corollary \ref{cor:pol}
it is the unique submodule of its dimension vector, and uniqueness forces
stability under the whole torus.
\end{proof}

From now on ``simple spectrum'' means $\dim V_i[d]\le1$ for all $i,d$.

\section{The socle-step identity and the main theorem}\label{sec:main}

This section is uniform in type and characteristic-free.

\begin{lemma}\label{lem:cb}
Let $X\subseteq I_\ka$ be a submodule, $v'=\dimv X$ and
$\mu'=\la-\sum_iv'_i\al_i$.  Then for every $i\in I$
\[
   \dim\soc_{(i)}(I_\ka/X)=\langle\al_i^\vee,\mu'\rangle+\dim\topp_i(X).
\]
\end{lemma}

\begin{proof}
Write $I=I_\ka$, $Q=I/X$.  Applying $\Hom_\Pia(S_i,-)$ to
$0\to X\to I\to Q\to 0$ gives
\[
0\to\Hom(S_i,X)\to\Hom(S_i,I)\to\Hom(S_i,Q)\to\Ext^1(S_i,X)\to\Ext^1(S_i,I),
\]
and the last term vanishes as $I$ is injective.  Since $\soc(I_\ka)=S_\ka$ we
have $\dim\Hom(S_i,I)=\delta_{i\ka}$, so
\begin{equation}\label{eq:les}
  \dim\soc_{(i)}(Q)=\delta_{i\ka}-\dim\soc_{(i)}(X)+\dim\Ext^1(S_i,X).
\end{equation}
By the $\Ext^1$ formula, Proposition \ref{prop:Pifacts}(a) --- proved over an
arbitrary field in Proposition \ref{prop:ext1} ---
\[
  \dim\Ext^1_\Pia(M,N)=\dim\Hom(M,N)+\dim\Hom(N,M)-(\dimv M,\dimv N),
\]
where $(x,y)=x^{\mathsf T}Cy$.  With $M=S_i$, $N=X$, and using
$\Hom(S_i,X)=\soc_{(i)}(X)$, $\Hom(X,S_i)=\topp_i(X)$, this reads
$\dim\Ext^1(S_i,X)=\dim\soc_{(i)}(X)+\dim\topp_i(X)-(\al_i,v')$.  Substituting
into \eqref{eq:les}, the socle terms cancel:
$\dim\soc_{(i)}(Q)=\delta_{i\ka}-(\al_i,v')+\dim\topp_i(X)$.  Finally
$(\al_i,v')=\langle\al_i^\vee,\sum_jv'_j\al_j\rangle$ and
$\langle\al_i^\vee,\om_\ka\rangle=\delta_{i\ka}$, so
$\delta_{i\ka}-(\al_i,v')=\langle\al_i^\vee,\mu'\rangle$.
\end{proof}

\begin{corollary}\label{cor:mono}
A socle step of colour $i$ applied to $X\subseteq I_\ka$ enlarges $X$ by
exactly $\langle\al_i^\vee,\mu'\rangle+\dim\topp_i(X)$ dimensions, all at
vertex $i$; in particular by at least $\langle\al_i^\vee,\mu'\rangle$, with
equality iff $S_i$ does not occur in the head of $X$.
\end{corollary}

\begin{remark}\label{rem:repeats}
Repeating a colour is harmless.  If a step of colour $i$ with
$c=\langle\al_i^\vee,\mu'\rangle$ has just added exactly $c$ dimensions,
producing $X'$, those $c$ elements lie in the head of $X'$, so
$\dim\topp_i(X')=c$, while the weight has become $s_i\mu'$ with
$\langle\al_i^\vee,s_i\mu'\rangle=-c$.  Lemma \ref{lem:cb} gives
$\dim\soc_{(i)}(I_\ka/X')=0$.  Thus $\soc_\sigma$ is unchanged if a block
$i^{\,c}$ is replaced by $i^{\,c'}$ for any $c'\ge c$, and ``apply
$\soc_{(i)}$ until it stabilises'' is an equivalent description.
\end{remark}

\begin{definition}[the $c$-trace]\label{def:ctrace}
Let $w=s_{i_1}\cdots s_{i_\ell}$ be reduced.  Set $\mu^{(\ell)}=\la$ and
$\mu^{(p-1)}=s_{i_p}\mu^{(p)}$, and put
$c_p=\langle\al_{i_p}^\vee,\mu^{(p)}\rangle$.  The $c$-weighted sequence is
$\sigma=(i_\ell^{\,c_\ell},\dots,i_1^{\,c_1})$, read from the right end of the
word inwards.
\end{definition}

\begin{lemma}\label{lem:extend}
Let $w$ be of minimal length in $wW_\la$ and $\mu=w\la$.
\begin{enumerate}
\item[(i)] Every $c$-trace value of every reduced word for $w$ satisfies
$c_p\ge1$.
\item[(ii)] If $\mu\ne w_0\la$ there is $i$ with
$\langle\al_i^\vee,\mu\rangle\ge1$, and for any such $i$ one has
$\ell(s_iw)=\ell(w)+1$, with $s_iw$ again minimal in $s_iwW_\la$.
\end{enumerate}
Consequently every reduced word for $w$ extends on the left to a reduced word
for the longest minimal-length representative $w_0^\la$, with all $c$-trace
values $\ge1$.
\end{lemma}

\begin{proof}
(i) Write $u=s_{i_{p+1}}\cdots s_{i_\ell}$, so $\mu^{(p)}=u\la$.  The word being
reduced gives $u^{-1}\al_{i_p}>0$, and $\la$ dominant gives
$c_p=\langle(u^{-1}\al_{i_p})^\vee,\la\rangle\ge0$.  If $c_p=0$ then
$s_{i_p}u\la=u\la$, so $w\la=s_{i_1}\cdots s_{i_{p-1}}u\la$ with an element of
length $\le\ell-1$, contradicting minimality.

(ii) If $\langle\al_i^\vee,\mu\rangle\le0$ for all $i$ then $\mu$ is
antidominant, i.e.\ $\mu=w_0\la$.  Otherwise pick $i$ with
$\langle\al_i^\vee,w\la\rangle>0$; then $\langle\la,w^{-1}\al_i^\vee\rangle>0$,
and $\la$ dominant forces $w^{-1}\al_i>0$, whence $\ell(s_iw)=\ell(w)+1$.  With
$J=\{j:\langle\al_j^\vee,\la\rangle=0\}$ one has $W_\la=W_J$, and $w$ minimal
in $wW_J$ iff $w\al_j>0$ for $j\in J$.  For such $j$, $s_i(w\al_j)<0$ would
force $w\al_j=\al_i$, giving
$\langle\al_i^\vee,w\la\rangle=\langle\al_j^\vee,\la\rangle=0$, a
contradiction.
\end{proof}

\begin{theorem}[Main Theorem]\label{thm:main}
Let $\Gamma$ be simply-laced Dynkin, $\la=\om_\ka$, $w\in W$ of minimal length
in $wW_\la$, and $\sigma$ the $c$-weighted sequence of any reduced word for
$w$.  Then
\[
   \dimv\soc_\sigma(I_\ka)\;=\;v\;=\;\la-w\la,
   \qquad\text{and consequently}\qquad
   V\;\cong\;\soc_\sigma(I_\ka)
\]
as graded representations of the doubled quiver.  In particular the
isomorphism class of $\soc_\sigma(I_\ka)$ is independent of the reduced word.
\end{theorem}

\begin{proof}
Extend the word, by Lemma \ref{lem:extend}, to a reduced word for $w_0^\la$;
let $\sigma_{\mathrm{full}}$ be its $c$-weighted sequence and $\sigma$ the
prefix belonging to $w$.  \emph{Index the blocks in the order in which they are
applied}: $X_0=0$, and $X_p$ is the result of applying the first $p$ blocks, so
that $X_0\subseteq X_1\subseteq\cdots\subseteq X_{\ell_{\mathrm{full}}}
=\soc_{\sigma_{\mathrm{full}}}(I_\ka)$, where $\ell_{\mathrm{full}}$ is the
number of blocks of $\sigma_{\mathrm{full}}$, i.e.\ the length of the word.
Write $i(p)$ for the colour of the $p$-th block applied and $c(p)$ for its
multiplicity, and put $\mu[0]=\la$, $\mu[p]=s_{i(p)}\mu[p-1]$; reading
Definition \ref{def:ctrace} from the right end of the word inwards, this is
exactly the $c$-trace, so
\[
   c(p)=\langle\al_{i(p)}^\vee,\mu[p-1]\rangle\ \ge\ 1
\]
by Lemma \ref{lem:extend}(i), and $\mu[p]=\mu[p-1]-c(p)\al_{i(p)}\in W\la$.

We prove by induction on $p$ that
\begin{equation}\label{eq:lower}
   \dimv X_p=\la-\mu[p]\qquad(0\le p\le \ell_{\mathrm{full}}),
\end{equation}
which for $p=0$ is $X_0=0$.  So let $p\ge1$ and assume
$\dimv X_{p-1}=\la-\mu[p-1]=:v$.  Write $i=i(p)$, $c=c(p)$, and let
$\delta_p:=\dim X_p-\dim X_{p-1}$; a block of colour $i$ enlarges $X$ only at
the vertex $i$, so $\dimv X_p=v+\delta_p\,\epsilon_i$ with $\epsilon_i$ the
$i$-th coordinate vector.

\emph{Lower bound: $\delta_p\ge c$.}  The inductive hypothesis is what makes
Corollary \ref{cor:mono} quotable with the \emph{intended} weight: applied to
$X_{p-1}$, whose weight $\mu'=\la-\sum_jv_j\al_j$ \emph{is} $\mu[p-1]$, the
first letter of the block already adds
$\langle\al_i^\vee,\mu[p-1]\rangle+\dim\topp_i(X_{p-1})\ge c$ dimensions, and
the remaining letters of the block only enlarge $X$ further.

\emph{Upper bound: $\delta_p\le c$.}  Here we use the inequality of Corollary
\ref{cor:pol}, applied to the submodule $X_p\subseteq I_\ka$.  Since
$\mu[p-1]\in W\la$, the same corollary gives $\delta(v)=0$.  With
$\delta(x)=2x_\ka-(x,x)$ and $(\epsilon_i,\epsilon_i)=2$,
$(\epsilon_i,v)=(Cv)_i$,
\[
   \delta\bigl(v+\delta_p\epsilon_i\bigr)
   =\delta(v)+2\delta_p\bigl(\delta_{i\ka}-(Cv)_i\bigr)-2\delta_p^{\,2}
   =2\delta_p\bigl(c-\delta_p\bigr),
\]
because $\delta_{i\ka}-(Cv)_i=\langle\al_i^\vee,\la-\sum_jv_j\al_j\rangle
=\langle\al_i^\vee,\mu[p-1]\rangle=c$.  Corollary \ref{cor:pol} says this is
$\ge0$, i.e.\ $2\delta_p(\delta_p-c)\le0$; and $\delta_p\ge c\ge1>0$, so
$\delta_p\le c$.

Hence $\delta_p=c$ and
$\dimv X_p=v+c\,\epsilon_i=\la-\mu[p-1]+c(p)\al_{i(p)}=\la-\mu[p]$, closing the
induction.  (Note that the two bounds are established \emph{at the same step},
from the inductive hypothesis alone; no later block is invoked.)

Truncating at $\sigma$, i.e.\ stopping at the value of $p$ at which the blocks
of $\sigma$ are exhausted, gives $\dimv\soc_\sigma(I_\ka)=\la-w\la$, which is
the $v$ of the statement.  By
Corollary \ref{cor:pol} the submodule of $I_\ka$ with this dimension vector is
unique and equals $V$.  (Taking $p=\ell_{\mathrm{full}}$ gives
$\dim\soc_{\sigma_{\mathrm{full}}}(I_\ka)=\dim(\la-w_0\la)=\dim I_\ka$ by
Proposition \ref{prop:dimI}, so the full sequence exhausts $I_\ka$; that
consequence is not used below.)
\end{proof}

\begin{corollary}\label{cor:crit}
$T$ acts with simple spectrum if and only if
$\dim\soc_\sigma(I_\ka)_i[d]\le1$ for all $i\in I$, $d$.  The criterion is
decidable in time polynomial in $|I|$ and $\ell(w)$ from any reduced word.
\end{corollary}

\begin{corollary}\label{cor:prefix}
Every $X_p$ of the filtration of Theorem \ref{thm:main} is itself extremal:
$\dimv X_p=\la-\mu[p]$ with $\mu[p]\in W\la$, by \eqref{eq:lower}.  In the
word, $\mu[p]=w'\la$ where $w'$ is the product of the last $p$ letters of the
reduced word, read as an element of $W$.
\end{corollary}

\subsection{The half-line wave, and the spin weights}\label{ss:halfline}

The whole of the multiplicity bookkeeping in this paper rests on one
elementary family of integer sequences, which we isolate now because it is
used twice: here for the spin weights, and in Lemma \ref{lem:P} for the
profile of $I_\ka$.  Nothing in this subsection uses the characteristic.

\begin{lemma}[the half-line wave]\label{lem:halfline}
Let $\rho\ge1$ be an integer.  For $j\ge0$ and $d\in\mathbb{Z}$ put
\[
   A^{\rho}_d(j)=
   \begin{cases}
     1 & |j-\rho|\le d\le j+\rho-2 \text{ and } d\equiv j+\rho\ (2),\\
     0 & \text{otherwise.}
   \end{cases}
\]
Then:
\begin{enumerate}
\item[\textup{(i)}] $A^{\rho}_{d+1}(j)+A^{\rho}_{d-1}(j)
=A^{\rho}_d(j-1)+A^{\rho}_d(j+1)$ for every $j\ge1$ and every $d$;
\item[\textup{(ii)}] $A^{\rho}_d(0)=0$, $A^{\rho}_0(j)=\delta_{j\rho}$ and
$A^{\rho}_d(j)=0$ for $d<0$;
\item[\textup{(iii)}] $A^{\rho}_d(j)=0$ whenever $d>j+\rho-2$; in particular
$A^{\rho}_d(j)=0$ for all $j\le\rho+1$ once $d\ge 2\rho$;
\item[\textup{(iv)}] $\#\{d:A^{\rho}_d(j)=1\}=\min(j,\rho)$.
\end{enumerate}
\end{lemma}

\begin{proof}
Let $u_d(i)=\bigl[\,|i-\rho|\le d\ \text{and}\ i\equiv\rho+d\ (2)\,\bigr]$ for
$i\in\mathbb{Z}$.  An immediate induction on $d$ shows that $u$ is the
solution of the free chain recursion
$u_{d+1}(i)+u_{d-1}(i)=u_d(i-1)+u_d(i+1)$ on $\mathbb{Z}$ with $u_0=\delta_\rho$
and $u_{-1}=0$.  Set $A^{\rho}_d(j)=u_d(j)-u_d(-j)$; being a difference of two
solutions of that recursion (the second reflected), $A^{\rho}$ satisfies it
for every $j$, which is (i).  Under the parity constraint $d\equiv j+\rho$ one
has $u_d(-j)=[\,j+\rho\le d\,]$, so
$u_d(j)-u_d(-j)=[\,|j-\rho|\le d\le j+\rho-2\,]$, which is the displayed
formula.  (ii) is clear ($A^{\rho}_d(0)=u_d(0)-u_d(0)=0$, and $d\ge0$ is
forced by $d\ge|j-\rho|$ unless $j=\rho$, where $d=0$ gives $1$).  (iii) is
read off the formula, together with $j+\rho-2\le2\rho-1$ for $j\le\rho+1$.
For (iv), the degrees of the correct parity in $[\,|j-\rho|,\ j+\rho-2\,]$
number $\frac{(j+\rho-2)-|j-\rho|}{2}+1=\frac{j+\rho-|j-\rho|}{2}
=\min(j,\rho)$.
\end{proof}

\begin{lemma}[folding the two leaves]\label{lem:fold}
Let $\Gamma=D_n$ with $n\ge4$, $t=n-2$, $L=n-1$, and let $\ka$ be any vertex.
Put
\[
   \hat a_d(j)=\dim(I_\ka)_j[d]\ \ (1\le j\le t),
   \qquad
   \hat a_d(L)=\dim(I_\ka)_{n-1}[d]+\dim(I_\ka)_{n}[d],
\]
and let $d^{*}$ be the top degree of $I_\ka$.  Then, with the boundary
conventions $\hat a_d(0):=0$ and $\hat a_d(L+1):=\hat a_d(L-1)$,
\[
   \hat a_{d+1}(j)+\hat a_{d-1}(j)=\hat a_d(j-1)+\hat a_d(j+1)
   \qquad(1\le j\le L,\ \ 0\le d\le d^{*}-1),
\]
with $\hat a_{-1}\equiv0$ and $\hat a_0(j)=\delta_{j,\iota(\ka)}$, where
$\iota(\ka)=\ka$ for a chain vertex and $\iota(\ka)=L$ for a leaf.
\end{lemma}

\begin{proof}
Lemma \ref{lem:R} gives $a_{d+1}(i)=\sum_{j\sim i}a_d(j)-a_{d-1}(i)$ at every
vertex $i$ and every $d<d^{*}$.  At a chain vertex $j<t$ the neighbours are
$j\pm1$, which is the displayed equation.  At $j=t$ the neighbours are $t-1$,
$n-1$ and $n$, and $\hat a_d(t-1)+\hat a_d(L)=\hat a_d(t-1)+\hat a_d(t+1)$
since $L=t+1$.  At a leaf the only neighbour is $t$; adding the two leaf
equations gives $\hat a_{d+1}(L)+\hat a_{d-1}(L)=2\hat a_d(t)
=\hat a_d(L-1)+\hat a_d(L+1)$ under the mirror convention.  The degree-$0$
part of $I_\ka\cong\Pia e_{\nu^{-1}(\ka)}$ is spanned by the idempotent at
$\nu^{-1}(\ka)$; after folding, that vertex is $\iota(\ka)$ whether or not
$\nu$ fixes $\ka$, whence the initial datum $\hat a_0(j)=\delta_{j,\iota(\ka)}$.
\end{proof}

\begin{lemma}[rigidity of the folded profile]\label{lem:rigid}
Let $n\ge4$, $t=n-2$, $L=n-1$, and let $A^{\rho}$ be as in Lemma
\ref{lem:halfline}.
\begin{enumerate}
\item[\textup{(a)}] Let $1\le k\le t$ and put
$\hat A_d(j)=A^{k}_d(j)+A^{k}_{2t-d}(j)$.  If $d\ge0$ and
$\hat A_d(j)=\delta_{jk}$ for every $1\le j\le L$, then $d\in\{0,2t\}$.
\item[\textup{(b)}] If $d\ge0$ and $A^{L}_d(j)=\delta_{jL}$ for every
$1\le j\le L$, then $d\in\{0,2L-2\}$.
\end{enumerate}
\end{lemma}

\begin{proof}
Throughout we use the defining formula: $A^{\rho}_d(j)=1$ exactly when
$|j-\rho|\le d\le j+\rho-2$ and $d\equiv j+\rho\ (2)$.

\emph{(a).}  For $j\le L$ one has $j+k-2\le L+t-2=2t-1$, so
$A^{k}_e(j)=0$ whenever $e>2t-1$; hence $\hat A_d=0$ for $d<0$ (the first
summand vanishes as $d<0$, the second as $2t-d>2t$).  Since
$\hat A_{2t-d}=\hat A_d$, it suffices to treat $0\le d\le t$ and show $d=0$.
So suppose $1\le d\le t$ and $\hat A_d=\delta_{\cdot k}$ on $\{1,\dots,L\}$;
we derive a contradiction in five steps.

\emph{(1) $d\ge t+2-k$, and $k\ge2$.}  If $k+d\le L$ then
$A^{k}_d(k+d)=1$: indeed $|(k+d)-k|=d\le d$, $d\le(k+d)+k-2$ because
$k\ge1$, and $(k+d)+k\equiv d\ (2)$.  As $k+d\ne k$ this contradicts
$\hat A_d=\delta_{\cdot k}$.  Hence $k+d>L=t+1$, i.e.\ $d\ge t+2-k$; with
$d\le t$ this forces $k\ge2$.

\emph{(2) $d\equiv t+k\ (2)$.}  Take $j=L=t+1$, which is $\ne k$ since
$k\le t$.  We have $|L-k|=t+1-k\le d$ by (1), and $d\le t\le L+k-2=t+k-1$;
also $|L-k|\le t\le 2t-d$ and $2t-d\le 2t-(t+2-k)=t+k-2\le L+k-2$.  So both
$A^{k}_d(L)$ and $A^{k}_{2t-d}(L)$ equal $1$ as soon as the common parity
condition $d\equiv 2t-d\equiv L+k\ (2)$ holds.  If it held we would get
$\hat A_d(L)=2\ne0$; so $d\equiv L+k+1=t+k\ (2)$.

\emph{(3) $k=t$.}  Take $j=t$.  Then $|t-k|=t-k\le t+2-k\le d$,
$d\le t\le t+k-2$ (using $k\ge2$), and $d\equiv t+k\ (2)$ by (2); hence
$A^{k}_d(t)=1$ and so $\hat A_d(t)\ge1$, forcing $t=k$.

\emph{(4) $d=2$.}  By (2) and (3), $d$ is even and $2\le d\le t$.  Put
$e=2t-d$ and $j_0=e-t+2=t-d+2$, so $2\le j_0\le t\le L$.  Then
$A^{t}_e(j_0)=1$: $|j_0-t|=d-2\le t\le e$; $e\le j_0+t-2=e$; and
$e\equiv j_0+t=2t-d\equiv e\ (2)$.  So $\hat A_d(j_0)\ge1$, forcing
$j_0=k=t$, i.e.\ $d=2$.

\emph{(5) Contradiction.}  So $k=t$ and $d=2$.  If $t\ge3$ then
$A^{t}_2(t-2)=1$ --- $|{-2}|\le2$, $2\le(t-2)+t-2=2t-4$, and
$2\equiv(t-2)+t\ (2)$ --- with $1\le t-2\ne t$, contradicting
$\hat A_2=\delta_{\cdot t}$.  If $t=2$ then $d=2=2t-d$, so
$\hat A_2(2)=2A^{2}_2(2)=2\ne1$, again a contradiction.

\emph{(b).}  From $A^{L}_d(L)=1$ we get $0\le d\le 2L-2$ and $d\equiv2L$,
i.e.\ $d$ even.  Suppose $2\le d\le2L-4$.  Since $n\ge4$ we have $L\ge3$, so
$1\le L-2<L$, and $A^{L}_d(L-2)=1$: $|(L-2)-L|=2\le d$,
$d\le(L-2)+L-2=2L-4$, and $d\equiv(L-2)+L=2L-2\ (2)$ as both are even.  This
contradicts $A^{L}_d=\delta_{\cdot L}$.  Hence $d=0$ or $d=2L-2$.
\end{proof}

\begin{proposition}[the folded solution]\label{prop:foldsol}
Let $\Gamma=D_n$, $n\ge4$, $t=n-2$, $L=n-1$.
\begin{enumerate}
\item[\textup{(a)}] If $\ka=k$ is a chain vertex with $1\le k\le t$, then the
top degree of $I_\ka$ is $2t$ and
\[
   \hat a_d(j)=A^{k}_d(j)+A^{k}_{2t-d}(j)\qquad(1\le j\le L,\ d\ge0).
\]
\item[\textup{(b)}] If $\ka$ is a leaf, then the top degree of $I_\ka$ is
$2L-2$ and $\hat a_d(j)=A^{L}_d(j)$ for $1\le j\le L$ and $d\ge0$.
\end{enumerate}
\end{proposition}

\begin{proof}
Write $\hat A$ for the proposed right-hand side and $\bar d$ for the asserted
top degree.  We first record three facts that hold in both cases.

\emph{(G1) $I_\ka$ is generated in degree $0$.}  By \eqref{eq:Ika},
$I_\ka\cong\Pia e_{\nu^{-1}(\ka)}$ is spanned by the paths starting at
$\nu^{-1}(\ka)$, so $(I_\ka)[d+1]$ is spanned by the elements $a\,x$ with
$x\in(I_\ka)[d]$ and $a$ an arrow.  Consequently $\hat a_d=0$ forces
$\hat a_{d'}=0$ for every $d'\ge d$, and $d^{*}$ is the unique degree with
$\hat a_{d^{*}}\ne0=\hat a_{d^{*}+1}$.

\emph{(G2) $d^{*}\ge1$.}  Pick $i\sim\ka$.  By Proposition \ref{prop:ext1},
$\dim\Ext^1_\Pia(S_i,S_\ka)=0+0-(\epsilon_i,\epsilon_\ka)=-C_{i\ka}=1\ne0$,
so $S_\ka$ is not injective and $I_\ka\ne S_\ka$.  (This is the only place in
this proof where the $\Ext^1$ formula is used, and Proposition
\ref{prop:ext1} supplies it over an arbitrary field.)

\emph{(G3) $\hat a_{d^{*}}(j)=\delta_{j,\iota(\ka)}$ for $1\le j\le L$.}  By
Lemma \ref{lem:soctop}, $(I_\ka)[d^{*}]=\soc(I_\ka)=S_\ka$ is one dimensional
and sits at the vertex $\ka$.  For a chain vertex $\ka=k\le t$ this is the
assertion; for a leaf $\ka$ the folded value $\hat a_{d^{*}}(L)$ is the sum of
the two leaf dimensions, hence $1$, and all other folded values vanish.

\emph{(a).}  Here $\bar d=2t$ and $\iota(\ka)=k$.  Non-negativity is clear.
Both $d\mapsto A^{k}_d(j)$ and $d\mapsto A^{k}_{2t-d}(j)$ satisfy the
recursion of Lemma \ref{lem:halfline}(i) for every $j\ge1$ (the recursion is
invariant under $d\mapsto2t-d$), and both vanish at $j=0$.  The mirror
condition $\hat A_d(L+1)=\hat A_d(L-1)$, i.e.\ $\hat A_d(t+2)=\hat A_d(t)$,
follows for \emph{every} $d$ from
\begin{equation}\label{eq:mirror}
   A^{k}_d(t+2)=A^{k}_{2t-d}(t)\qquad\text{for all }d,
\end{equation}
applied both as it stands and with $d$ replaced by $2t-d$.  Indeed, $k\le t$
gives $|t+2-k|=t+2-k$ and $|t-k|=t-k$, so
$A^{k}_d(t+2)=[\,t+2-k\le d\le t+k\,]$ while
$A^{k}_{2t-d}(t)=[\,t-k\le2t-d\le t+k-2\,]=[\,t-k+2\le d\le t+k\,]$: the same
interval, and both conditions carry the parity $d\equiv t+k$.  Initial data:
$A^{k}_{2t}(j)=0$ and $A^{k}_{2t+1}(j)=0$ for $j\le L$ by Lemma
\ref{lem:halfline}(iii), since $j+k-2\le L+t-2=2t-1$; hence
$\hat A_0(j)=A^{k}_0(j)=\delta_{jk}$ and $\hat A_{-1}=0$.  Vanishing above
$2t$: for $d>2t$ both summands vanish, the first by
\ref{lem:halfline}(iii) and the second because $2t-d<0$.  Finally
$\hat A_{2t}(j)=A^{k}_{2t}(j)+A^{k}_{0}(j)=\delta_{jk}$, by the same two
facts.

So $\hat A$ is a non-negative solution of the recursion and the initial data
of Lemma \ref{lem:fold}, \emph{valid at every degree}.  Since the recursion
of Lemma \ref{lem:fold} holds for $I_\ka$ at every degree $d\le d^{*}-1$ and
determines $\hat a_{d+1}$ from $\hat a_d,\hat a_{d-1}$, induction on $d$
gives
\[
   \hat a_d=\hat A_d\qquad\text{for every } d\le d^{*}.
\]
If $d^{*}>2t$ this would give $\hat a_{d^{*}}=\hat A_{d^{*}}=0$, absurd; so
$d^{*}\le2t$.  By (G3), $\hat A_{d^{*}}=\hat a_{d^{*}}=\delta_{\cdot k}$ on
$\{1,\dots,L\}$, so Lemma \ref{lem:rigid}(a) gives $d^{*}\in\{0,2t\}$, and
$d^{*}\ne0$ by (G2).  Hence $d^{*}=2t$, and $\hat a_d=\hat A_d$ for
$d\le2t$; for $d>2t$ both sides vanish, by (G1) and by the previous
paragraph.  This is the assertion.

\emph{(b).}  Here $\bar d=2L-2$, $\iota(\ka)=L$ and $\hat A_d(j)=A^{L}_d(j)$,
which is non-negative, vanishes at $j=0$ and satisfies the recursion for all
$j\ge1$.  Its initial data are $\hat A_0(j)=\delta_{jL}=\delta_{j,\iota(\ka)}$
and $\hat A_{-1}=0$, as required.  The mirror condition holds for
$0\le d\le2L-3$: $A^{L}_d(L+1)=[\,1\le d\le2L-1\,]$ and
$A^{L}_d(L-1)=[\,1\le d\le 2L-3\,]$ agree there (both with $d$ odd).  Hence
induction on $d$, using Lemma \ref{lem:fold} at degrees $d\le d^{*}-1$ and
the mirror condition at those degrees, gives
\[
   \hat a_d=\hat A_d\qquad\text{for every } d\le\min(d^{*},2L-2).
\]
Suppose first $d^{*}<2L-2$.  Then $\hat A_{d^{*}}=\hat a_{d^{*}}
=\delta_{\cdot L}$ by (G3), so $d^{*}\in\{0,2L-2\}$ by Lemma
\ref{lem:rigid}(b), hence $d^{*}=0$, contradicting (G2).  So
$d^{*}\ge2L-2$ and $\hat a_d=\hat A_d$ for all $d\le2L-2$.  Suppose next
$d^{*}>2L-2$.  Reading off the defining formula,
$\hat A_{2L-2}(j)=\delta_{jL}$ (the condition $2L-2\le j+L-2$ forces $j\ge L$)
and $\hat A_{2L-3}(j)=\delta_{j,L-1}$ (the condition $2L-3\le j+L-2$ forces
$j\ge L-1$, and $j=L$ has the wrong parity).  Since $2L-2\le d^{*}-1$, Lemma
\ref{lem:fold} applies at degree $2L-2$ and gives, for every $1\le j\le L$,
\[
   \hat a_{2L-1}(j)=\hat A_{2L-2}(j-1)+\hat A_{2L-2}(j+1)-\hat A_{2L-3}(j),
\]
with the conventions $\hat A_{d}(0)=0$ and $\hat A_d(L+1):=\hat A_d(L-1)$.
At $j=L$ this is $2\hat A_{2L-2}(L-1)-\hat A_{2L-3}(L)=0-0=0$; at $j=L-1$ it
is $\hat A_{2L-2}(L-2)+\hat A_{2L-2}(L)-\hat A_{2L-3}(L-1)=0+1-1=0$; and at
$j<L-1$ all three terms vanish.  So $\hat a_{2L-1}=0$, whence $d^{*}\le2L-2$
by (G1) --- a contradiction.  Therefore $d^{*}=2L-2$ and $\hat a=\hat A$.
\end{proof}

\begin{corollary}[the dimension of $I_\ka$ in type $D$]\label{cor:dimI}
Let $\Gamma=D_n$ with $n\ge4$, and let $\mathbb{K}$ be an arbitrary field.
Then
\[
   \dim_{\mathbb{K}}I_\ka=k(2n-1-k)\ \ (\ka=k\le n-2),
   \qquad
   \dim_{\mathbb{K}}I_\ka=\tfrac{n(n-1)}{2}\ \ (\ka\in\{n-1,n\}).
\]
Equivalently, $\dimv I_\ka=\om_\ka-w_0\om_\ka$ in root coordinates.  This is
Proposition \ref{prop:dimI}; it is proved here from Lemma
\ref{lem:R}, Lemma \ref{lem:soctop}, Proposition \ref{prop:ext1} and integer
combinatorics, with no hypothesis on $\mathbb{K}$.
\end{corollary}

\begin{proof}
Sum Proposition \ref{prop:foldsol} over all $d\ge0$ and all $1\le j\le L$;
the folded total is the total dimension, since folding only adds the two
leaf dimensions.  In case (a), $\hat A_d(j)=0$ for $d<0$ and $d\mapsto2t-d$
carries the support of $A^{k}_{\bullet}(j)$ into that of
$A^{k}_{2t-\bullet}(j)$ bijectively, so by Lemma \ref{lem:halfline}(iv)
\[
   \dim I_\ka=\sum_{d\ge0}\sum_{j=1}^{L}\hat A_d(j)
   =2\sum_{j=1}^{L}\min(j,k)
   =2\Bigl(\sum_{j=1}^{k}j+\sum_{j=k+1}^{n-1}k\Bigr)
   =k(k+1)+2k(n-1-k)=k(2n-1-k).
\]
In case (b), Lemma \ref{lem:halfline}(iv) gives
$\sum_{j=1}^{L}\min(j,L)=\sum_{j=1}^{L}j=\tfrac{L(L+1)}{2}
=\tfrac{n(n-1)}{2}$.  The final sentence is the root-coordinate computation
carried out in the proof of Proposition \ref{prop:dimI}.
\end{proof}

\begin{corollary}[multiplicity-freeness at the three minuscule
vertices]\label{cor:mfree}
Let $\Gamma=D_n$, $n\ge4$, and $\ka\in\{1,n-1,n\}$.  Then
$\dim(I_\ka)_i[d]\le1$ for every vertex $i$ and every degree $d$.
\end{corollary}

\begin{proof}
For $\ka=1$ apply Proposition \ref{prop:foldsol}(a) with $k=1$ (legitimate,
since $1\le t$ for $n\ge4$): $A^{1}_d(j)=[\,d=j-1\,]$, so
$\hat a_d(j)=[\,d=j-1\,]+[\,d=2t-j+1\,]$.  For $j\le t$ the two degrees
$j-1\le t-1$ and $2t-j+1\ge t+1$ are distinct, so $\hat a_d(j)\le1$.  At
$j=L=t+1$ they coincide and $\hat a_d(L)=2[\,d=t\,]$, so each of the two
leaves has $\dim(I_1)_{l}[d]=b_d$ with $2b_d\le2$, i.e.\ $b_d\le1$.

For $\ka\in\{n-1,n\}$ apply Proposition \ref{prop:foldsol}(b):
$\hat a_d(j)=A^{L}_d(j)\le1$ for every $j\le L$.  For $j\le t$ this is the
claim directly; for $j=L$ it says that the two leaf dimensions \emph{sum} to
at most $1$, so each is at most $1$.
\end{proof}

\begin{proof}[Proof of Theorem \ref{thm:spin}]
Let $\la=\om_\ka$ with $\ka\in\{1,n-1,n\}$, and let $w$ be of minimal length
in $wW_\la$.  These three weights of $D_n$ are minuscule, so
$\langle\al^\vee,\nu\rangle\in\{0,\pm1\}$ for every root $\al$ and every
weight $\nu$ of the irreducible representation of highest weight $\la$; hence
every $c$-trace value $c_p$ equals $1$, which
is exactly Stembridge's condition that $w$ be $\la$-minuscule
\cite[\S2]{Stem}, and is \cite[Def.~5.1]{DKK} verbatim --- note that
\cite[\S5.1]{DKK} is headed ``dominant minuscule'' while its Definition 5.1
defines ``$\la$-minuscule'', by the condition
$\langle\al_{i_p}^\vee,s_{i_{p+1}}\cdots s_{i_\ell}(\la)\rangle=1$ for
$1\le p\le\ell$.  This is DKK's condition verbatim with their word index
renamed to $p$; they write $k$ for it, which in this paper is the
fundamental-weight label.  It is exactly our $c_p=1$: Definition
\ref{def:ctrace} sets $\mu^{(\ell)}=\la$ and $\mu^{(p-1)}=s_{i_p}\mu^{(p)}$,
so $\mu^{(p)}=s_{i_{p+1}}\cdots s_{i_\ell}\la$ --- an empty product, hence
$\mu^{(\ell)}=\la$, at the endpoint $p=\ell$ --- and therefore
$c_p=\langle\al_{i_p}^\vee,\mu^{(p)}\rangle
=\langle\al_{i_p}^\vee,s_{i_{p+1}}\cdots s_{i_\ell}\la\rangle$, character for
character.  By Stembridge
\cite[Prop.~2.1]{Stem} such a $w$ is in particular fully commutative; we do not
use this.

For the spectrum there is nothing more to do, and in particular no geometry is
needed.  By Theorem \ref{thm:main}, which holds over an arbitrary field,
$V=\soc_\sigma(I_\ka)$ is a \emph{submodule} of $I_\ka$; and it is a
\emph{graded} submodule, because $\soc_\sigma$ is built from socle steps and a
socle step is degree-homogeneous --- $\soc_{(j)}(I_\ka/X)$ is cut out of the
graded module $I_\ka/X$ by the degree-homogeneous conditions
$a_{j\to i}x\in X$, so it is a graded subspace, and $X=0$ is graded to begin
with.  Therefore
\[
   \dim V_i[d]\ \le\ \dim(I_\ka)_i[d]\ \le\ 1
   \qquad\text{for every } i \text{ and } d
\]
by Corollary \ref{cor:mfree}, which is simplicity of the spectrum in the sense
fixed after Proposition \ref{prop:grading}.

No hypothesis on the characteristic was used, and nothing here is quoted from
a source that carries one: Theorem \ref{thm:main} rests on Lemma \ref{lem:cb}
and Corollary \ref{cor:pol}, both proved over an arbitrary field (the former
through Proposition \ref{prop:ext1}), and Corollary \ref{cor:mfree} rests on
Proposition \ref{prop:foldsol}, a statement about the dimensions of $I_\ka$
proved from Lemma \ref{lem:R}, Lemma \ref{lem:soctop}, Proposition
\ref{prop:ext1} and Lemma \ref{lem:rigid} alone.  Over
$\mathbb{C}$, Proposition \ref{prop:ST} together with Corollary
\ref{cor:pol} identifies this $V$ with the $T$-fixed representation and
Proposition \ref{prop:grading} identifies the path-length grading with the
$T$-weight grading, giving the geometric reading of the statement; neither is
used in the argument above.
\end{proof}

\begin{remark}
The route through the heap module of \cite[\S3]{DEKM} and the
claim that the rank function of $H(w)$ is the path-length degree --- is
avoided entirely.  Corollary \ref{cor:mfree} is stronger than what the heap
argument would give: it bounds \emph{every} submodule of $I_\ka$, not only the
extremal ones.
\end{remark}


\section{The injective \texorpdfstring{$I_\ka$}{I} in type \texorpdfstring{$D$}{D}}
\label{sec:profile}

\begin{convention}\label{conv:D}
From now on $\Gamma=D_n$ with $n\ge4$, Bourbaki labelling
$\al_i=e_i-e_{i+1}$ ($i\le n-1$), $\al_n=e_{n-1}+e_n$.  The trivalent node is
$t=n-2$, the two \emph{leaves} are $n-1$ and $n$, and the \emph{chain} is
$\{1,\dots,t\}$.  We fix $\la=\om_k$ with $2\le k\le n-2$ and $\ka=k$, so that
$I_\ka=\Pia e_k$.  The nontrivial diagram automorphism exchanging $n-1$ and $n$
is denoted $\tau$; it fixes every chain vertex, hence fixes $\ka$, hence acts
on $I_\ka$.  For $\mu\in W\la$ we write $v=\la-\mu$ in root coordinates, so
$\mu=\la-\sum_iv_i\al_i$.

We also fix the direction words once and for all: along the chain, the arrow
$a_{j\to j+1}$ points \emph{inward} (towards the node $t$) and $a_{j\to j-1}$
points \emph{outward}; the arrows $a_{t\to l}$ to the leaves and $a_{l\to t}$
from them are called the \emph{node--leaf} arrows.

Finally we fix the symbols that would otherwise collide.  Throughout,
$\varsigma$ is the sign convention on arrows, $\eps\in\{+,-\}$ is a strand
label (\S\ref{sec:strands}), $\eta_i$ denotes a Bourbaki coordinate of a
weight, and $\ell$ is the length of a reduced word or of a colour sequence.
The letter $a$ is reserved for the quiver: the arrows $a_{i\to j}$ and the
multiplicities $a_d(j),\ a^{\eps}_d(j),\ \hat a_d(j)$ of $I_\ka$; the integer
parameter of the nine families of \S\ref{sec:families}, and the corresponding
parameter in $v'_t=2h$ or $2h+1$, is called $h$ instead.  The letter $M$ is
reserved for modules, and $\ms_j:=\min(j,k)$ --- the number of slots of one
strand at a chain vertex $j$ --- is set in upright type to keep it apart.
Two \emph{different} degrees occur and are named differently: $D_j(X)$ is
always the threshold of Definition \ref{def:upset} (the largest degree at
which $X_j$ is a \emph{proper} subspace of $(I_\ka)_j$), while $E_j$ is always
the lowest \emph{taken} degree of a greedy top truncation at $j$, with
$E^{\eps}_j$ its refinement to one strand (Lemma \ref{lem:Deps}); the two are
compared in Remark \ref{rem:nodefree}.  The letter $L$ is used only for the
folded leaf vertex $L=n-1$ of Lemma \ref{lem:fold}; the Lagrangian
subvariety $\mathfrak{L}(v,w)$ of \S\ref{sec:setup} is in fraktur.

Two reuses are left standing because they are separated both by section and by
shape, and no formula contains both.  In \S\S\ref{sec:families}--\ref{sec:closure}
the constant $B:=n-1-k$ never takes an argument, whereas in
\S\ref{sec:staircase} the double-path operators $A(x),B(x),C(x)$ always do;
and the half-line wave $A_d(j)$ of \S\ref{ss:halfline} always carries a degree
subscript.  Likewise $m$ and $r$ of Theorem \ref{thm:typeD}, and the scalars
$m_{\pm}$ of \S\ref{sec:staircase}, are lower-case italic and never carry a
vertex subscript; the operators $P(x)=B(x)+C(x)$ and $Q(x)=B(x)-C(x)$ of
\S\ref{sec:staircase} always carry their argument, as does the extension
module $E(\psi)$ of \S\ref{ss:cbproof}, which is the one other use of the
letter $E$.

Four letters carry more than one meaning, and we flag them rather than pretend
otherwise.  $N$ is a module in \S\ref{sec:setup} and \S\ref{ss:negative}, the
integer $k(2n-1-k)$ in \S\ref{sec:charp}, the saturated lattice $N(v)$ of
\S\ref{sec:charp}, and a multiplicity $N_q$ in \S\ref{sec:reps}.  $P$ is a
projective $P_j$ in \S\ref{sec:setup} and an operator $P(x)$ in
\S\ref{sec:staircase}; the poset of \S\ref{sec:closure} is $\mathcal{P}$.
$\ell$ is the length of a reduced word or colour sequence, and a prime in
\S\ref{sec:charp}; path length is written $|\cdot|$.  Most often reused is
$p$: a prime in \S\ref{sec:charp}, the block index of Theorem \ref{thm:main}
and of the $c$-trace, a path in the proof of Lemma \ref{lem:R}, a Bourbaki
coordinate index in \S\ref{sec:profile} and \S\ref{sec:families}, and the
number of positive signs in Lemma \ref{lem:C}.  Likewise $c$ is $k-q-h$ in
\S\ref{sec:reps}, $\langle\al_t^\vee,\mu'\rangle$ in \S\ref{sec:induction},
$\nu^{-1}(\ka)$ in \S\ref{sec:profile}, and a $c$-trace entry $c_p$.  In every
case the meaning is fixed by the section together with whether the symbol
carries a subscript, an argument or neither.
\end{convention}

\begin{lemma}[the profile of $I_\ka$]\label{lem:P}
The colour-$t$ multiplicity sequence of $I_\ka$ is
\[
   (1,\underbrace{2,\dots,2}_{k-1},1)
   \quad\text{at the degrees}\quad t-k,\ t-k+2,\ \dots,\ t+k,
\]
and $\dim(I_\ka)_t[d]=0$ otherwise.  A chain vertex $i\le t$ carries
multiplicity $2$ in exactly $\max(0,i+k-n+1)$ degrees and multiplicity $\le 2$
throughout.  The two leaves are multiplicity-free, with support
\[
   t-k+1,\ t-k+3,\ \dots,\ t+k-1 \qquad (k \text{ degrees}),
\]
and $a_d(n-1)=a_d(n)$ for every $d$.  The top degree of $I_\ka$ is $2t$.
\end{lemma}

\begin{proof}
Since $\ka=k\le t$ is a chain vertex, Proposition \ref{prop:foldsol}(a)
applies and gives, in the folded notation of Lemma \ref{lem:fold},
\begin{equation}\label{eq:foldedP}
   \hat a_d(j)=A^{k}_d(j)+A^{k}_{2t-d}(j)\qquad(1\le j\le L=t+1),
\end{equation}
with top degree $2t$.  Every assertion of the lemma is read off
\eqref{eq:foldedP}; we record the four readings.  Throughout, degrees at $j$
carry the parity of $j+k$, and by Lemma \ref{lem:halfline}
\[
   A^{k}_d(j)=[\,|j-k|\le d\le j+k-2\,],
   \qquad
   A^{k}_{2t-d}(j)=[\,2t-(j+k-2)\le d\le 2t-|j-k|\,].
\]

\emph{The node.}  At $j=t$ the two intervals are $[\,t-k,\,t+k-2\,]$ and
$[\,t-k+2,\,t+k\,]$.  Their union is $[\,t-k,\,t+k\,]$ and their intersection
is $[\,t-k+2,\,t+k-2\,]$, which contains
$\frac{(t+k-2)-(t-k+2)}{2}+1=k-1$ degrees of the right parity.  So the
colour-$t$ multiplicity sequence is $(1,2^{\,k-1},1)$ at the degrees
$t-k,t-k+2,\dots,t+k$, and $0$ elsewhere.

\emph{Chain vertices.}  For $i\le t$ the two intervals are
$[\,|i-k|,\ i+k-2\,]$ and $[\,2t+2-i-k,\ 2t-|i-k|\,]$, each carrying
$\min(i,k)$ degrees by Lemma \ref{lem:halfline}(iv), so $\hat a_d(i)\le2$
throughout.  As $i,k\le t$ we have $i+k-2\le 2t-|i-k|$, so the intersection is
$[\,\max(|i-k|,\,2t+2-i-k),\ i+k-2\,]$; it is nonempty exactly when
$2t+2-i-k\le i+k-2$ is not the binding constraint, and in all cases the number
of degrees of the right parity in it is
$\max\bigl(0,\ i+k-t-1\bigr)=\max(0,\,i+k-n+1)$.

\emph{The leaves.}  At $j=L=t+1$, using $k\le t$, both intervals are
$[\,t+1-k,\ t+k-1\,]$, so $\hat a_d(L)=2\,[\,t-k+1\le d\le t+k-1\,]$.  Since
$\hat a_d(L)=\dim(I_\ka)_{n-1}[d]+\dim(I_\ka)_n[d]$ and $\tau$ exchanges the
two leaves while fixing $I_\ka$, the two summands are equal; hence
$a_d(n-1)=a_d(n)=b_d=[\,t-k+1\le d\le t+k-1\,]$, so the leaves are
multiplicity-free with the displayed support of $k$ degrees.

\emph{The top degree} is $2t$ by Proposition \ref{prop:foldsol}(a).
\end{proof}

\begin{remark}\label{rem:imagetotal}
The predicted total is genuinely part of the argument, and it is
\emph{termwise non-negative} because \eqref{eq:foldedP} is a sum of two
indicator functions rather than an alternating sum over the images of $k$
under the infinite dihedral group.  Its value, computed in the proof of
Proposition \ref{prop:foldsol}(a), is
$2\sum_{j=1}^{n-1}\min(j,k)=k(2n-1-k)$, which is exactly
$\dim(\la-w_0\la)$ by Proposition \ref{prop:dimI}.
\end{remark}

\begin{lemma}\label{lem:C}
With $m$ and $r$ as in Theorem \ref{thm:typeD} one has $v_t=r+2m$.
\end{lemma}

\begin{proof}
Write $\la-\mu=\sum_jc_je_j$.  Comparing coefficients of $e_j$ in
$\la-\mu=\sum_iv_i\al_i$ gives $c_1=v_1$ and $c_j=v_j-v_{j-1}$ for
$2\le j\le n-2$, so $v_t=\sum_{j\le n-2}c_j$.  Since $k\le n-2$ the support of
$\la$ lies in the first $n-2$ coordinates, whence
$\sum_{j\le n-2}c_j=k-\sum_{j\le n-2}\mu_j$.  Writing $p$ and $m$ for the
numbers of positive and negative signs of $\mu$ among the first $n-2$
coordinates, $\sum_{j\le n-2}\mu_j=p-m$ and $p+m=k-r$, so $v_t=k-(p-m)=r+2m$.
\end{proof}

\subsection{Reduction to the trivalent node}

\begin{lemma}[inward injectivity]\label{lem:Id}
Let $i<t$ be a chain vertex and let $d$ satisfy $\dim(I_\ka)_i[d]\ge2$.  Then
the inward arrow $a_{i\to i+1}:(I_\ka)_i[d]\to(I_\ka)_{i+1}[d+1]$ is injective.
\end{lemma}

\begin{proof}
Suppose $x\ne0$ lies in the kernel.  We first record the degree bound that
drives the whole iteration.  Since $\dim(I_\ka)_i[d]\ge2$, the slot $(d,i)$
lies in the multiplicity-$2$ window of $i$, whose top degree is $i+k-2$ by
Lemma \ref{lem:P}; so
\begin{equation}\label{eq:Iddeg}
   d\le i+k-2,
   \qquad\text{whence}\qquad
   d+(i-k)\ \le\ 2i-2\ \le\ 2t-4\ <\ 2t
\end{equation}
(the last step because $i\le t-1$).  Note also $d\le i+k-2\le 2t-3<2t$.

Now argue as follows.  For $i<t$ the vertex $i$ has exactly the two neighbours
$i-1$ and $i+1$.  By Lemma \ref{lem:R} an element killed by every arrow out of
$i$ is either zero or the socle generator, which by Lemma \ref{lem:soctop}
lives in the top degree $2t$ at the vertex $\ka$; by \eqref{eq:Iddeg}
$d<2t$, so $y:=a_{i\to i-1}(x)\ne0$ (for $i=1$ there is no such
$y$ and we get the contradiction at once).  The preprojective relation at $i$,
$\varsigma\,a_{i+1\to i}a_{i\to i+1}+\varsigma'\,a_{i-1\to i}a_{i\to i-1}=0$, applied to
$x$, has vanishing first term, so $a_{i-1\to i}(y)=0$.  Thus $y\ne0$ lies in
$(I_\ka)_{i-1}[d+1]$ and is killed by \emph{its} inward arrow: the same
situation one vertex further out and one degree higher.

At the general step we sit at $(d+p,\,i-p)$ with a nonzero element killed by
its inward arrow, and we must know that the outward arrow does \emph{not} kill
it, so that the iteration continues.  If both arrows out of $i-p$ killed it,
the element would span a copy of $S_{i-p}$ inside $\soc(I_\ka)$, forcing
$i-p=\ka=k$ \emph{and} $d+p=2t$ by Lemma \ref{lem:soctop}; but $i-p=k$ means
$p=i-k$, and then \eqref{eq:Iddeg} gives $d+p\le2t-4<2t$.  So the outward
arrow survives at every step, and the iteration reaches vertex $1$: it
produces $0\ne z\in(I_\ka)_1[d+i-1]$ killed by $a_{1\to2}$.  Vertex $1$ has
$2$ as its only neighbour, so $z\in\soc(I_\ka)=S_\ka$ and $\ka=1$ ---
contradicting $\ka=k\ge2$ of Convention \ref{conv:D}.
\end{proof}

\begin{theorem}\label{thm:lemA}
If $\dim V_t[d]\le1$ for all $d$, then $\dim V_i[d]\le1$ for all $i\in I$ and
all $d$.
\end{theorem}

\begin{proof}
The two leaves are multiplicity-free in $I_\ka$ by Lemma \ref{lem:P}, so they
never carry a repeat.  Suppose $\dim V_i[d]\ge2$ for a chain vertex $i<t$.
Then $\dim(I_\ka)_i[d]\ge2$, so by Lemma \ref{lem:Id} the inward arrow is
injective on $(I_\ka)_i[d]$, hence on the subspace $V_i[d]$; as $V$ is a
submodule the image lies in $V_{i+1}[d+1]$, so $\dim V_{i+1}[d+1]\ge2$.
Iterating $t-i$ times gives $\dim V_t[d+t-i]\ge2$.
\end{proof}

\begin{remark}\label{rem:nowindow}
The more obvious argument --- containment of multiplicity-$2$ windows --- does
not work.  The multiplicity-$2$ window of vertex $i\le t$ is
$[2L-k-i,\ i+k)$ with $L=n-1$, and its endpoints do satisfy $2L-k-i\ge L-k+1$
and $i+k\le(L-1)+k$, so it is contained in that of $t$ \emph{as an interval}.
But the degrees carried by $i$ obey the parity constraint $d\equiv i+k\pmod2$,
so the degree \emph{sets} need not be nested: in $D_6$ with $\ka=3$ and $t=4$,
vertex $3$ carries multiplicity $2$ exactly at degree $4$ while $t$ carries it
exactly at degrees $3$ and $5$ --- disjoint sets.  And even granting
containment, nothing transports $\dim V_i[d]\ge2$ to $\dim V_t[d']\ge2$.  The
degree shift built into Lemma \ref{lem:Id} is what makes the $D_6$ numbers line
up: $i=3$ at $d=4$ maps to $t$ at $d=5$.
\end{remark}

\subsection{The statement that carries the theorem}

\begin{definition}\label{def:upset}
Let $X\subseteq I_\ka$ be a graded submodule and $j$ a vertex.  We say $X$ is an
\emph{up-set at $j$} if there is a degree $D$ with $X_j[d]=0$ for $d<D$,
$X_j[D]\subseteq(I_\ka)_j[D]$ arbitrary, and $X_j[d]=(I_\ka)_j[d]$ for $d>D$.
At a vertex where $\dim(I_\ka)_j[d]\le1$ this says exactly that $X_j$ occupies a
final segment of degrees; at $t$ we also call it a \emph{top truncation}.

The degree $D$ above is not unique --- it may be replaced by the next
colour-$j$ degree below whenever $X_j[D]$ happens to be full, or by the next
one above whenever $X_j[D]=0$.  \emph{We therefore fix once and for all the
convention}
\[
   D_j(X):=\text{the largest degree at which } X_j \text{ is a proper
   subspace of } (I_\ka)_j,
\]
with $D_j(X):=-\infty$ when $X_j$ is full in every degree.  This is the only
meaning the symbol $D_j(X)$ carries below, and it is the one all later
arguments use: in particular $X_j[D_j(X)]$ is always proper, and
$X_j[d]=(I_\ka)_j[d]$ for every $d>D_j(X)$.
\end{definition}

The following lemma reads off, from the profile $(1,2^{k-1},1)$ of Lemma
\ref{lem:P} alone, everything the induction of \S\ref{sec:induction} needs to
know about the node column of a top truncation.  It is stated for an arbitrary
graded submodule \emph{that is a top truncation at $t$} --- no reduced word,
and no parameter of the nine families of \S\ref{sec:families}: the integer $h$
below is simply $\lfloor v'_t/2\rfloor$, not the $h$ of $\mu(h;x,y,z)$ ---
only the quantities $v'_t$ and $D_t(X)$, which exist for any
such $X$ --- and it is proved over an arbitrary field.  The top-truncation
hypothesis cannot be dropped from \textup{(a)}: if $X_t$ is the line in the
bottom one-dimensional slot, at degree $t-k$, then $v'_t=1$ but
$D_t(X)=t+k$, whereas the table for $v'_t=1$ reads $D=t+k-2$.  The lemma is
placed here because it is used at $p=2$, and because the standing hypothesis
$\operatorname{char}\ne2$ of \S\S\ref{sec:strands}--\ref{sec:T012} and
\S\S\ref{sec:families}--\ref{sec:closure} is never in force for it.

\begin{lemma}[the node column of a top truncation]\label{lem:nodefree}
Let $X\subseteq I_\ka$ be a graded submodule which is a top truncation at $t$,
put $v'_t=\dim X_t$ and let $D=D_t(X)$ be the threshold of Definition
\ref{def:upset}.
\begin{enumerate}
\item[\textup{(a)}] If $v'_t<2k$ then $D$ is finite and
\[
   D=\begin{cases}
     t+k, & v'_t=0,\\
     t+k-2h, & v'_t=2h \text{ with } h\ge1,\\
     t+k-2h-2, & v'_t=2h+1 \text{ with } h\ge0,
   \end{cases}
\]
and
\[
   X_t[D]\ne0\iff v'_t \text{ is even and } v'_t\ge2 ,
\]
in which case $X_t[D]$ is a line in the two-dimensional space
$(I_\ka)_t[D]$; while for $v'_t$ odd, and for $v'_t=0$, one has $X_t[D]=0$.
\item[\textup{(b)}] If moreover $\dimv X$ is extremal and
$c:=\langle\al_t^\vee,\la-\dimv X\rangle\ge1$, then $v'_t\le2k-1$, so the
hypothesis $v'_t<2k$ of \textup{(a)} is automatic.
\end{enumerate}
\end{lemma}

\begin{proof}
\emph{(a).}  By Lemma \ref{lem:P} the colour-$t$ column of $I_\ka$ consists of
a one-dimensional slot at $t+k$, then $k-1$ two-dimensional slots at
$t+k-2,t+k-4,\dots,t-k+2$, then a one-dimensional slot at $t-k$; its total
dimension is $2k$, and there are no other colour-$t$ degrees.  Being a top
truncation, $X_t$ is full above $D$, proper at $D$ and zero below $D$
(Definition \ref{def:upset}); so $v'_t$ is the total dimension of the column
strictly above $D$, plus $\dim X_t[D]$.  Since $v'_t<2k$ the column is not
full, so some degree is proper and $D$ is finite.

If $v'_t=0$ then $X_t=0$ and the largest proper degree is the top one,
$D=t+k$; and $X_t[D]=0$.

If $v'_t=2h$ with $h\ge1$, then $v'_t\le2k-2$, so $h-1\le k-2$ and the degrees
$t+k,t+k-2,\dots,t+k-2(h-1)$ --- the top slot and the top $h-1$
two-dimensional slots --- carry $1+2(h-1)=2h-1$ dimensions in all; one further
dimension is taken at the next degree down, $t+k-2h$, which satisfies
$t+k-2h\ge t-k+2$ and is therefore two-dimensional.  So $D=t+k-2h$ and
$X_t[D]$ is a line in a two-dimensional space: nonzero and proper.

If $v'_t=2h+1$, then $v'_t\le2k-1$, so $h\le k-1$ and the degrees
$t+k,t+k-2,\dots,t+k-2h$ carry $1+2h=v'_t$ dimensions in all and are exactly
the filled ones; the next degree down is $t+k-2h-2\ge t-k$, which is a genuine
colour-$t$ degree and is untaken.  So $D=t+k-2h-2$ and $X_t[D]=0$.

This proves the displayed table, and the equivalence
``$X_t[D]\ne0\iff v'_t$ even and $\ge2$'' is read off it.

\emph{(b).}  Since $\dimv X$ is extremal, $\mu':=\la-\sum_i(\dimv X)_i\al_i$
lies in $W\la=W\om_k$, and for $2\le k\le n-2$ that orbit is
$\{\sum_{i\in S}\eta_ie_i:|S|=k,\ \eta_i=\pm1\}$ (see the first paragraph of
the proof of Proposition \ref{prop:reps}); so we may write
$\mu'=\sum_i\eta_ie_i$ with $\eta_i\in\{0,\pm1\}$ and exactly $k$ nonzero
entries.  By Lemma \ref{lem:C} applied to
$\mu'$, $v'_t=r'+2m'$ where $r'=|S'\cap\{n-1,n\}|$ and
$m'=\#\{i\le n-2: i\in S',\ \eta_i=-1\}$, $S'$ being the support of $\mu'$.
Since $|S'|=k$ we have $m'\le k-r'$, hence
\[
   v'_t=r'+2m'\ \le\ r'+2(k-r')\ =\ 2k-r'\ \le\ 2k ,
\]
with equality throughout only if $r'=0$ and $m'=k$.  In that case every one of
the $k$ elements of $S'$ is a chain coordinate $\le n-2$ carrying
$\eta_i=-1$, and $r'=0$ gives $\eta_{n-1}=\eta_n=0$; so $\eta_{n-2}\in\{0,-1\}$
and
\[
   c=\langle\al_t^\vee,\mu'\rangle=\eta_{n-2}-\eta_{n-1}=\eta_{n-2}\le0 ,
\]
contradicting $c\ge1$.  Hence $v'_t\le2k-1$.
\end{proof}

\begin{remark}\label{rem:nodefree}
Lemma \ref{lem:nodefree} is deliberately \emph{not} a corollary of Lemma
\ref{lem:node}.  The latter is stated inside \S\ref{sec:families}, under the
standing hypothesis $\operatorname{char}\ne2$, for a fixed representative
$\mu=\mu(h;x,y,z)$, and its conclusions are written in the family parameters
$h,q,x,y,z$, which do not exist for an arbitrary extremal submodule of
$I_\ka$.  Its $E_t$ is moreover a \emph{different} quantity: it is the lowest
\emph{taken} colour-$t$ degree, whereas the $D$ of Definition \ref{def:upset}
is the largest degree at which the column is \emph{proper}.  The two agree when
$v'_t$ is even \emph{and positive} --- there the lowest taken slot is the
partial one --- and differ by $2$ when $v'_t$ is odd, where the lowest taken
slot is full and the threshold sits one slot below it.  At $v'_t=0$ there is no
taken slot at all, so the $E_t$ of Lemma \ref{lem:node} does not exist and no
comparison is to be made, while the threshold of Definition \ref{def:upset} is
$t+k$; this is the form recorded in \S\ref{sec:verification}.  Lemma \ref{lem:nodefree} is the statement in
the $D$ of Definition \ref{def:upset}, and it is the one used in
\S\ref{sec:induction}; part \textup{(b)} supplies the boundary hypothesis
$v'_t<2k$, which Lemma \ref{lem:node} secures instead by the family
computation $v_t=2k-2h-s$ with $2h+s\ge q\ge1$.
\end{remark}

\begin{proposition}[top truncation]\label{prop:T}
Let $\mu=w\la$ and $V=V(\mu)$.  Then $V$ is a top truncation at $t$: its
colour-$t$ multiplicities are the greedy top truncation of the sequence of
Lemma \ref{lem:P} to total $v_t$.
\end{proposition}

Proposition \ref{prop:T} is proved in \S\ref{sec:induction}.  We record at once
that it implies the theorem.

\begin{proof}[Proof of Theorem \ref{thm:typeD}, assuming Proposition
\ref{prop:T}]
By Lemma \ref{lem:P} the colour-$t$ profile of $I_\ka$ is $(1,2^{k-1},1)$.  By
Theorem \ref{thm:main} and Corollary \ref{cor:pol}, $V$ is the unique submodule
of $I_\ka$ with $\dimv V=v$, and Proposition \ref{prop:T} applies to it.  The
greedy top truncation of $(1,2^{k-1},1)$ to total $v_t$ is
\[
 v_t=0\mapsto(\,),\ \ 1\mapsto(1),\ \ 2\mapsto(1,1),\ \ 3\mapsto(1,2),
 \ \ 4\mapsto(1,2,1),\ \ 5\mapsto(1,2,2),\ \dots
\]
so every colour-$t$ multiplicity is $\le1$ exactly when $v_t\le2$ (for
$v_t\ge3$ the sequence contains a $2$, since $k\ge2$ guarantees a
multiplicity-$2$ slot).  By Theorem \ref{thm:lemA}, simplicity at $t$ implies
simplicity everywhere, and conversely simplicity everywhere implies simplicity
at $t$.  Lemma \ref{lem:C} identifies $v_t=r+2m$.  Theorem \ref{thm:spin}
covers the spin weights.
\end{proof}

\section{The strand decomposition}\label{sec:strands}

\emph{Throughout \S\ref{sec:strands} and \S\ref{sec:T012}, and again throughout
\S\ref{sec:families} and \S\ref{sec:closure}, we assume
$\operatorname{char}\mathbb{K}\ne2$.}  The intervening \S\ref{sec:reps} is
\emph{not} covered by this hypothesis and does not use it: Definition
\ref{def:star}, Lemma \ref{lem:Jinv}, Corollary \ref{cor:reduce}, Proposition
\ref{prop:reps}, Definition \ref{def:B} and \eqref{eq:Brig} are proved over an
arbitrary field, and they are quoted at $p=2$ in \S\ref{sec:charp}; the
hypothesis lapses at the head of \S\ref{sec:reps} and resumes at the head of
\S\ref{sec:families}, as each of those sections says.  The hypothesis is used
to construct the witness submodule of $(B)$; \S\ref{sec:charp} transports
$(B)$, and with it Theorem \ref{thm:typeD}, to every characteristic.

The automorphism $\tau$ fixes every chain vertex and exchanges the two leaves,
and it acts on $I_\ka$ (Convention \ref{conv:D}).  Since $2$ is invertible,
each graded piece at a \emph{chain} vertex splits into $\tau$-eigenspaces.
Write $a^{\pm}_d(j)=\dim\bigl((I_\ka)_j[d]\bigr)^{\pm}$ for $j\le t$, and call
the two eigen-parts the $+$ \emph{strand} and the $-$ \emph{strand}.  Recall
$t=n-2$ and that the top degree of $I_\ka$ is $2t$ (Lemma \ref{lem:P}).

\begin{proposition}\label{prop:W2}
For every chain vertex $j$ with $1\le j\le t$ and every $d$,
\begin{equation}\label{eq:W2}
   a^{+}_d(j)=
   \begin{cases}
     1 & |j-k|\le d\le j+k-2 \text{ and } d\equiv j+k \ (2),\\
     0 & \text{otherwise},
   \end{cases}
   \qquad
   a^{-}_d(j)=a^{+}_{2t-d}(j).
\end{equation}
Consequently:
\begin{itemize}\itemsep1pt
\item[\textup{(W1)}] each eigenspace is at most one dimensional in every chain
slot, so a multiplicity-$2$ slot is exactly one $+$ line and one $-$ line;
\item[\textup{(W2)}] the two strands are exchanged by the degree reversal
$d\mapsto2t-d$, that is $a^{-}_d(j)=a^{+}_{2t-d}(j)$ at every chain vertex ---
this is the second half of \eqref{eq:W2}, and it is the statement referred to
as \textup{(W2)} elsewhere in the paper;
\item[\textup{(W3)}] each strand occupies a contiguous run of degrees at every
chain vertex, the $+$ strand the low ones and the $-$ strand the high ones,
overlapping exactly in the multiplicity-$2$ window.
\end{itemize}
\end{proposition}

We record the supports once and for all: for $1\le j\le t$,
\begin{equation}\label{eq:supports}
   S^{+}(j)=\bigl[\,|j-k|,\ j+k-2\,\bigr],
   \qquad
   S^{-}(j)=\bigl[\,2t-(j+k-2),\ 2t-|j-k|\,\bigr],
\end{equation}
both in steps of $2$, degrees at $j$ having the parity of $j+k$.  We write
$\lo^{\eps}_j,\hi^{\eps}_j$ for the endpoints, so
\[
   \hi^{+}_j=j+k-2,\qquad \hi^{-}_j=2t-|j-k| .
\]
At the node, $S^{+}(t)\cup S^{-}(t)=[t-k,\,t+k]$, of dimension $2$ on
$[t-k+2,\,t+k-2]$ and of dimension $1$ at exactly the two ends $d=t-k$ ($+$
alone) and $d=t+k$ ($-$ alone); this matches Lemma \ref{lem:P}.

\begin{proof}[Proof of Proposition \ref{prop:W2}]
\emph{Step 1: the eigen-refined recursion.}  The complex of Lemma \ref{lem:R}
at a chain vertex $j$,
\[
   (I_\ka)_j[d-1]\longrightarrow\bigoplus_{i\sim j}(I_\ka)_i[d]
   \longrightarrow (I_\ka)_j[d+1],
\]
consists of $\tau$-equivariant maps: the arrows between chain vertices commute
with $\tau$, and for $j=t$ the two leaf summands are exchanged by $\tau$, so
their sum $J[d]:=(I_\ka)_{n-1}[d]\oplus(I_\ka)_n[d]$ is $\tau$-stable.  As
$\operatorname{char}\ne2$ the complex splits into eigen-parts, and exactness,
together with the injectivity of the first map below the top degree (Lemma
\ref{lem:R}), is inherited by each summand.  The pair $J[d]$ is the induced
representation of $\mathbb{Z}/2$ on $(I_\ka)_{n-1}[d]$, so each of its
eigen-parts has dimension $b_d:=\dim(I_\ka)_{n-1}[d]$.  Counting dimensions in
the $\eps$-part therefore gives, for every chain vertex $j$ and every $d$ with
$0\le d<2t$ \textup{(}the case $d=0$ included: there the first map of the
complex has source $(I_\ka)_j[-1]=0$ and the count is the trivial one, and it
is needed to start the induction below\textup{)},
\begin{equation}\label{eq:eigrec}
   a^{\eps}_{d+1}(j)+a^{\eps}_{d-1}(j)
   =\sum_{i\sim j,\ i\le t}a^{\eps}_d(i)+[\,j=t\,]\,b_d ,
\end{equation}
while Lemma \ref{lem:R} at a leaf reads
\begin{equation}\label{eq:leafrec}
   b_{d+1}+b_{d-1}=a^{+}_d(t)+a^{-}_d(t).
\end{equation}
The generator $e_k$ spans $(I_\ka)_k[0]$ and is $\tau$-fixed, so
$a^{+}_0(j)=\delta_{jk}$, $a^{-}_0\equiv0$, $b_0=0$, and
$a^{\pm}_{-1}=b_{-1}=0$.  Each equation expresses degree $d+1$ in terms of
degrees $d$ and $d-1$, so the system \eqref{eq:eigrec}--\eqref{eq:leafrec} with
these initial data has \emph{at most one} solution, determined by induction on
$d$.  It therefore suffices to check that \eqref{eq:W2} is a solution.

\emph{Step 2: the $+$ strand is the Dirichlet half-line wave.}  Let
$A_d(j):=A^{k}_d(j)$ be the half-line wave of Lemma \ref{lem:halfline} with
$\rho=k$; thus
\[
   A_d(j)=\bigl[\,|j-k|\le d\le j+k-2,\ d\equiv j+k\ (2)\,\bigr],
\]
which is exactly the claimed $a^{+}_d(j)$ of \eqref{eq:W2}, and by Lemma
\ref{lem:halfline}(i),(ii) it satisfies
$A_{d+1}(j)+A_{d-1}(j)=A_d(j-1)+A_d(j+1)$ for every $j\ge1$, with $A_d(0)=0$,
$A_0(j)=\delta_{jk}$ and $A_{-1}=0$.  Then
\eqref{eq:eigrec} for $\eps=+$ and
$1\le j\le t-1$ is the free recursion at $j$; and at $j=t$ it is the free
recursion at $t$ with the value $A_d(t+1)$ in the role of the leaf
contribution.  So the $+$ part of \eqref{eq:eigrec} holds, provided we show
$b_d=A_d(t+1)$.

\emph{Step 3: the $-$ strand.}  Put $a^{-}_d(j):=A_{2t-d}(j)$.  The recursion
$X_{d+1}(j)+X_{d-1}(j)=X_d(j-1)+X_d(j+1)$ is invariant under $d\mapsto 2t-d$,
so $a^{-}$ satisfies \eqref{eq:eigrec} for $\eps=-$ with the leaf value
$A_{2t-d}(t+1)$.  That this is the \emph{same} leaf value as for the $+$ strand,
i.e.\ that
\begin{equation}\label{eq:Bsym}
   A_d(t+1)=A_{2t-d}(t+1)\qquad\text{for every } d,
\end{equation}
is the one symmetry the argument needs, and it is visible from the closed
form: since $k\le t$, $|t+1-k|=t+1-k$ and
$A_d(t+1)=[\,t+1-k\le d\le t+k-1\,]$, an interval symmetric about $t$, so
$d$ lies in it if and only if $2t-d$ does; the parity condition
$d\equiv t+1+k$ is likewise invariant under $d\mapsto2t-d$.  \emph{Without
\eqref{eq:Bsym} the $-$ part of \eqref{eq:eigrec} would not follow from the
$+$ part.}  Its initial data are right: $a^{-}_0(j)=A_{2t}(j)$ vanishes for
$j\le t+1$ because $A_{2t}(j)\ne0$ would force $2t\le j+k-2$, i.e.\
$j+k\ge 2t+2$, while $j\le t+1$ and $k\le t$ give $j+k\le 2t+1$; and
$a^{-}_{-1}(j)=A_{2t+1}(j)=0$ likewise.

\emph{Step 4: the leaves.}  It remains to check that $A_d(t+1)$ satisfies
\eqref{eq:leafrec}, i.e.\ $A_{d+1}(t+1)+A_{d-1}(t+1)=a^{+}_d(t)+a^{-}_d(t)$.
The free recursion at $j=t+1$ gives
$A_{d+1}(t+1)+A_{d-1}(t+1)=A_d(t)+A_d(t+2)$, and
$A_d(t)=a^{+}_d(t)$, so we must show $A_d(t+2)=a^{-}_d(t)=A_{2t-d}(t)$.  Since
$k\le t$ we have $|t+2-k|=t+2-k$ and $|t-k|=t-k$, so
\[
   A_d(t+2)=[\,t+2-k\le d\le t+k\,],
   \qquad
   A_{2t-d}(t)=[\,t-k\le 2t-d\le t+k-2\,]=[\,t-k+2\le d\le t+k\,],
\]
which are the same interval; the parities agree too, both conditions reading
$d\equiv t+k\ (2)$.  Hence $A_\bullet(t+1)$ solves \eqref{eq:leafrec} and, by the
uniqueness of Step 1, \eqref{eq:W2} holds.

\emph{Step 5: (W1) and (W3).}  Both are immediate from \eqref{eq:W2}: each
$a^{\eps}_d(j)\in\{0,1\}$, and each strand occupies the interval
\eqref{eq:supports}.  That $S^{+}(j)$ lies below $S^{-}(j)$ and that the
overlap is the multiplicity-$2$ window is the statement
$\hi^{+}_j=j+k-2\le \hi^-_j$ and $\lo^-_j=2t+2-j-k\ge\lo^+_j=|j-k|$, both
immediate for $j\le t$, $k\le t$.  The overlap
$W_j=[\,2t+2-j-k,\ j+k-2\,]$ has
$o_j:=\max(0,\,j+k-t-1)$ slots, in agreement with Lemma \ref{lem:P}.
\end{proof}

\begin{remark}\label{rem:char2}
This is the only place in the four sections \S\ref{sec:strands},
\S\ref{sec:T012}, \S\ref{sec:families}, \S\ref{sec:closure} where
$\operatorname{char}\ne2$ is used --- the rest of them use it only through
Proposition \ref{prop:W2} and its consequences --- and it is used essentially: at $\operatorname{char}=2$ the operator $\tau$ is unipotent rather
than semisimple and the splitting fails, so (W1)--(W4) are not available.  They
are also not repairable: see Remark \ref{rem:filtfail}.  What
\S\ref{sec:charp} does instead is to dispense with them at $p=2$ altogether.
\end{remark}

We now prove that the arrows of $I_\ka$ transport each strand.

\begin{lemma}\label{lem:hi}
For either strand and every chain vertex $1<j<t$,
$\min\bigl(\hi^{\eps}_{j-1},\hi^{\eps}_{j+1}\bigr)=\hi^{\eps}_j-1$.
\end{lemma}

\begin{proof}
For $\eps=+$, $\hi^{+}_j=j+k-2$ is affine in $j$, so the minimum is
$\hi^+_{j-1}=\hi^+_j-1$.  For $\eps=-$, $\hi^-_j=2t-|j-k|$ and
$\max(|j-1-k|,|j+1-k|)=|j-k|+1$.
\end{proof}

\begin{lemma}[a death propagates]\label{lem:death}
Let $1<j<t$ and let $x$ span the $\eps$-line at $(d,j)$.  Suppose
$a_{j\to i}x=0$ for one chain neighbour $i$ of $j$, the slot $(d+1,i)$ carrying
$\eps$.  Write $i'$ for the other chain neighbour.  Then $a_{j\to i'}x$ spans
the $\eps$-line at $(d+1,i')$, the slot $(d+2,j)$ carries $\eps$, and
$a_{i'\to j}\bigl(a_{j\to i'}x\bigr)=0$.
\end{lemma}

\begin{proof}
As $j<t$ its only neighbours are $i$ and $i'$, so the $\eps$-part of Lemma
\ref{lem:R} (i.e.\ the injectivity of $Y$ restricted to the $\eps$-part) forces
$a_{j\to i'}x\ne0$; it lies in the $\eps$-part of $(I_\ka)_{i'}[d+1]$, which is
at most a line by (W1), so that slot carries $\eps$.  Both slots $(d+1,j\pm1)$
therefore carry $\eps$, whence
$d+1\le\min(\hi^{\eps}_{j-1},\hi^{\eps}_{j+1})=\hi^{\eps}_j-1$ by Lemma
\ref{lem:hi}, i.e.\ $d+2\le\hi^{\eps}_j$; and $d+2\ge d\ge\lo^{\eps}_j$ with
the parity unchanged, so $(d+2,j)$ carries $\eps$.  Finally the preprojective
relation at $j$ reads
$\varsigma_1a_{i\to j}a_{j\to i}x+\varsigma_2a_{i'\to j}a_{j\to i'}x=0$, whose first term
vanishes.
\end{proof}

\begin{proposition}[chain arrows towards the node]\label{prop:W4out}
For either strand, $a_{j\to(j+1)}$ is nonzero on the $\eps$-line at $(d,j)$
whenever $(d+1,j+1)$ carries $\eps$ \textup{(}here $1\le j\le t-1$\textup{)}.
\end{proposition}

\begin{proof}
Suppose $a_{j\to(j+1)}x=0$.  If $j=1$ the only chain neighbour is $2$ and the
$\eps$-part of Lemma \ref{lem:R} is contradicted at once.  Otherwise Lemma
\ref{lem:death} with $i=j+1$, $i'=j-1$ reproduces the same situation at
$(d+1,j-1)$: its arrow towards the node dies and the target $(d+2,j)$ carries $\eps$.
Iterating lowers the vertex to $1$.
\end{proof}

\begin{proposition}[the folded leaf arrow]\label{prop:W4l}
Let $J[d]=(I_\ka)_{n-1}[d]\oplus(I_\ka)_n[d]$.  The induced map
$J[d]^{\eps}\to(I_\ka)_t[d+1]^{\eps}$ is nonzero whenever both sides are.
\end{proposition}

\begin{proof}
The leaves are multiplicity-free (Lemma \ref{lem:P}), so $(I_\ka)_{n-1}[d]$ is
a line, spanned by $u$ say, and $J[d]^{\eps}$ is spanned by
$(u,\eps\,\tau u)$, whose image is $(1+\eps\tau)z$ with $z=a_{(n-1)\to t}u$.
As $n-1$ has the single neighbour $t$, Lemma \ref{lem:R} gives $z\ne0$, and the
relation at $n-1$ reads $a_{t\to(n-1)}z=0$.  If $\dim(I_\ka)_t[d+1]=1$ the
claim is exactly $z\ne0$.

So let that slot be $2$-dimensional and suppose $z$ were a $\tau$-eigenvector
of eigenvalue $\eps$.  Applying $\tau$ to $a_{t\to(n-1)}z=0$ gives
$a_{t\to n}z=0$, so both leaf arrows kill $z$; by Lemma \ref{lem:R} at $t$,
$y:=a_{t\to(t-1)}z\ne0$, and $y$ lies in the $\eps$-part at $(d+2,t-1)$, which
is therefore nonzero.  The relation at $t$ applied to $z$ has both leaf terms
zero, so $a_{(t-1)\to t}y=0$.  Finally $(d+3,t)$ carries $\eps$: for $\eps=+$
because $d+2\le\hi^{+}_{t-1}=\hi^{+}_t-1$; for $\eps=-$ because the slot
$(d+1,t)$ then carries $+$ as well, so $d+1\le\hi^{+}_t=t+k-2$ and hence
$d+3\le t+k=\hi^{-}_t$.  Thus the arrow towards the node at $(d+2,t-1)$ dies with its
target carrying $\eps$, contradicting Proposition \ref{prop:W4out}.  So $z$ is
not an eigenvector and both of its eigen-components are nonzero.
\end{proof}

\begin{proposition}[outward node arrows]\label{prop:W4n}
Let $z$ span the $\eps$-eigenline of $(I_\ka)_t[d]$ and let $l\in\{n-1,n\}$
satisfy $(I_\ka)_l[d+1]\ne0$.  Then $a_{t\to l}z\ne0$.
\end{proposition}

\begin{proof}
Suppose $a_{t\to l}z=0$.

\emph{(1) Both leaf arrows die.}  $\tau$ exchanges the leaves and the two
arrows out of $t$, and $\tau z=\eps z$; applying $\tau$ to
$a_{t\to(n-1)}z=0$ gives $\eps\,a_{t\to n}z=0$, so both leaf arrows kill $z$
(and symmetrically).

\emph{(2) The chain arrow cannot die too.}  If $a_{t\to(t-1)}z=0$ as well then
all three arrows out of $t$ kill $z$, so $\mathbb{K}z$ is a copy of $S_t$ inside
$\soc(I_\ka)=S_\ka$, which sits in the top degree $2t$; hence $d=2t$.  But the
leaves have top degree $t+k-1\le 2t-1$, so $(I_\ka)_l[d+1]=0$, against the
hypothesis.

\emph{(3)} So $y:=a_{t\to(t-1)}z\ne0$.  It sits at $(d+1,t-1)$, and $\tau$
commutes with $a_{t\to(t-1)}$, so $y$ is again an $\eps$-eigenvector; by (W1)
the slot $(d+1,t-1)$ carries $\eps$.  The relation at $t$,
\[
   \varsigma_1a_{(t-1)\to t}a_{t\to(t-1)}z+\varsigma_2a_{(n-1)\to t}a_{t\to(n-1)}z
   +\varsigma_3a_{n\to t}a_{t\to n}z=0,
\]
has its last two terms zero by (1), so $a_{(t-1)\to t}y=0$.

\emph{(4)} By Proposition \ref{prop:W4out} applied to the arrow
$(t-1)\to t$ at degree $d+1$, the slot $(d+2,t)$ does \emph{not} carry $\eps$,
i.e.\ $d=\hi^{\eps}_t$.

\emph{(5) Both values of $\eps$ are impossible.}  For $\eps=+$, $\hi^+_t=t+k-2$,
so $d+1=t+k-1>t+k-3=\hi^+_{t-1}$ and the slot $(d+1,t-1)$ does not carry $+$,
contradicting (3).  For $\eps=-$, $\hi^-_t=t+k$, so $d+1=t+k+1$ exceeds the
leaf top degree $t+k-1$ and $(I_\ka)_l[d+1]=0$, contradicting the hypothesis.
\end{proof}

\begin{corollary}[\textup{(W4)}]\label{cor:W4}
Every arrow of $I_\ka$ restricts to an isomorphism between consecutive nonzero
slots of the same strand.  Precisely: for every arrow $a_{j\to i}$ and every
degree $d$ such that the slots $(d,j)$ and $(d+1,i)$ both carry $\eps$
\textup{(}with the folded convention $J[\cdot]^{\eps}$ at the leaves\textup{)},
the map $a_{j\to i}$ is nonzero, hence an isomorphism, on the $\eps$-parts.
\end{corollary}

\begin{proof}
There are four classes of arrow.

\emph{Chain inward \textup{(}towards the node\textup{)}, $j\to j+1$ with $j\le t-1$}: Proposition
\ref{prop:W4out}.

\emph{Leaf inward, $l\to t$}: Proposition \ref{prop:W4l}.

\emph{Node outward, $t\to l$}: Proposition \ref{prop:W4n}.

\emph{Chain outward \textup{(}away from the node\textup{)}, $j\to j-1$ with $2\le j\le t$}: suppose the $\eps$-line at
$(d,j)$ is killed, the slot $(d+1,j-1)$ carrying $\eps$.  If $j<t$, Lemma
\ref{lem:death} (with $i=j-1$, $i'=j+1$) produces the same situation at
$(d+1,j+1)$, again with a dying outward arrow; iterating raises the vertex to
$t$.  So we may assume $j=t$: the $\eps$-line at $(D,t)$ is killed by
$a_{t\to(t-1)}$ and $(D+1,t-1)$ carries $\eps$.  The relation at $t$ then makes
the composite $J[D+1]^\eps\to(I_\ka)_t[D+2]^\eps$ vanish on the image of the
$\eps$-line, while Lemma \ref{lem:R} makes that image nonzero in $J^{\eps}$;
by Proposition \ref{prop:W4l} this is a contradiction \emph{provided} the slot
$(D+2,t)$ carries $\eps$.  We check that it does, except in one case which we
handle separately.

The slot $(D+1,t-1)$ carries $\eps$, so $D+1\le\hi^{\eps}_{t-1}$.  For
$\eps=+$, $\hi^+_{t-1}=\hi^+_t-1$, so $D+2\le\hi^+_t$ and $(D+2,t)$ carries
$+$: no exception.  For $\eps=-$, $\hi^-_{t-1}=2t-|t-1-k|$ and
$\hi^-_t=2t-(t-k)=t+k$; the condition $D+1\le\hi^-_{t-1}$ reads $D\le\hi^-_t$
(when $k\le t-1$; when $k=t$ it reads $D\le \hi^-_t-2$), while $(D+2,t)$ carries
$-$ iff $D\le\hi^-_t-2$.  So the \emph{only} uncovered arrival is
$D=\hi^{-}_t=t+k$, and it occurs only for $k\le t-1$.

At $D=t+k$ we argue directly and do not use Proposition \ref{prop:W4l} at all.
The leaves are multiplicity-free with top degree $t+k-1$ (Lemma \ref{lem:P}),
so $(I_\ka)_{n-1}[D+1]=(I_\ka)_n[D+1]=0$, since $D+1=t+k+1>t+k-1$.  If in
addition $a_{t\to(t-1)}z=0$ for the vector $z$ spanning the $-$line at $(D,t)$,
then all three arrows out of $t$ kill $z$, so $\mathbb{K}z$ is a copy of $S_t$
inside $\soc(I_\ka)=S_\ka$, which lies in the top degree $2t$.  But
$D=t+k<2t$, because this configuration requires $k\le t-1$.  Contradiction.
\end{proof}

\begin{remark}
The last paragraph is the same socle argument that proves Lemma \ref{lem:R},
and it is characteristic-free.  It is needed: the target slot $(D+2,t)$ really
is empty at $D=\hi^-_t$, so the appeal to Proposition \ref{prop:W4l} would be
vacuous there.
\end{remark}

\section{Ranks of arrows between chain slots}\label{sec:T012}

By Corollary \ref{cor:W4} the rank of an arrow $(d,j)\to(d+1,i)$ between chain
slots is the number of strands carried by \emph{both} slots.  The following
three consequences are what the closure arguments of \S\ref{sec:closure} use.
We call the $+$ line of a multiplicity-$2$ chain slot its
\emph{$\tau$-eigenline} (it is the $\tau=+1$ eigenspace).

\begin{proposition}\label{prop:T012}
Let $i,j\le t$ be adjacent chain vertices and $d$ a degree.  Then:
\begin{itemize}\itemsep1pt
\item[\textup{(T0)}] $\operatorname{rank}\bigl(a_{j\to i}:(I_\ka)_j[d]\to
(I_\ka)_i[d+1]\bigr)=\min\bigl(\dim(I_\ka)_j[d],\ \dim(I_\ka)_i[d+1]\bigr)$;
\item[\textup{(T1)}] if $\dim(I_\ka)_j[d]=1$ and $\dim(I_\ka)_i[d+1]=2$ then
the image of $a_{j\to i}$ is the $\tau$-eigenline of $(I_\ka)_i[d+1]$;
\item[\textup{(T2)}] if $\dim(I_\ka)_j[d]=2$ and $\dim(I_\ka)_i[d+1]=1$ then
$a_{j\to i}$ annihilates the $\tau$-eigenline of $(I_\ka)_j[d]$.
\end{itemize}
Moreover the $\tau$-eigenline of a multiplicity-$2$ chain slot is \emph{not}
annihilated by an arrow into a $2$-dimensional chain slot \textup{(}Corollary
\ref{cor:W4}\textup{)}, nor by an arrow into a nonzero leaf slot
\textup{(}Proposition \ref{prop:W4n}\textup{)}.
\end{proposition}

\begin{proof}
By Corollary \ref{cor:W4} the rank is $|\Sigma(d,j)\cap\Sigma(d+1,i)|$, where
$\Sigma$ denotes the set of strands present.  So (T0) says exactly that the two
sets $\Sigma(d,j),\Sigma(d+1,i)$ are \emph{nested}; and (T1), (T2) say that in
the two mixed cases the single strand present at the $1$-dimensional slot is
$+$, respectively $-$.  All three are therefore statements about the intervals
\eqref{eq:supports}, and we verify them.  Throughout $i,j,k\le t$ and
$|i-j|=1$; degrees at $j$ have the parity of $j+k$.

\emph{(T0).}  Suppose the sets are not nested; then one slot carries only $+$
and the other only $-$.

\emph{Case A: $+$ only at $(d,j)$, $-$ only at $(d+1,i)$.}  From $+$ at
$(d,j)$, $d\ge|j-k|\ge|i-k|-1$, so $d+1\ge|i-k|$; hence the failure of $+$ at
$(d+1,i)$ must be $d+1>i+k-2$, i.e.\ $d\ge i+k-2$, while $+$ at $(d,j)$ gives
$d\le j+k-2$.  If $i=j+1$ this reads $d\le i+k-3<i+k-2\le d$, absurd.  If
$i=j-1$ then $d\in\{i+k-2,i+k-1\}$ and the parity $d\equiv j+k=i+k+1$ forces
$d=i+k-1=j+k-2$; now $-$ at $(d+1,i)=(i+k,i)$ needs $i+k\ge2t-i-k+2$, i.e.\
$i+k\ge t+1$, while the absence of $-$ at $(d,j)$ needs $j+k-2<2t-j-k+2$, i.e.\
$j+k\le t+1$ --- impossible for $j=i+1$.

\emph{Case B: $-$ only at $(d,j)$, $+$ only at $(d+1,i)$.}  From $+$ at
$(d+1,i)$, $d\le i+k-3$.  The absence of $+$ at $(d,j)$ is either $d>j+k-2$,
which with $d\le i+k-3$ forces $j<i-1$, impossible; or $d<|j-k|$, in which case
$-$ at $(d,j)$ gives $2t-j-k+2\le d<|j-k|$, reading $2t+2<2j$ if $j\ge k$ and
$2t+2<2k$ if $j<k$, both impossible.

\emph{(T1).}  Here the source carries one strand and the target two.  Suppose
the source strand were $-$.  This is \emph{not} the configuration of Case B:
there the target carried $+$ only, whereas here it carries both strands.  But
the proof of Case B never uses the absence of $-$ at $(d+1,i)$.  Reading it
again, it uses exactly three facts --- that $+$ is present at $(d+1,i)$, that
$+$ is absent at $(d,j)$, and that $-$ is present at $(d,j)$ --- and derives a
contradiction from those alone.  All three hold in the present situation: the
target is $2$-dimensional, so it carries $+$; and the source is
$1$-dimensional carrying $-$, so it carries $-$ and not $+$.  The same
contradiction therefore applies, and the source strand is $+$.  By Corollary
\ref{cor:W4} the image is then the $+$ line of the target.

\emph{(T2).}  Here the source carries both strands and the target one; suppose
the target strand were $+$.  The absence of $-$ at $(d+1,i)$ is either
$d\le 2t-i-k$, which against $d\ge2t-j-k+2$ (from $-$ at $(d,j)$) gives
$i\le j-2$; or $d\ge2t-|i-k|$, which against $d\le j+k-2$ (from $+$ at $(d,j)$)
gives $2t+2\le i+j$ when $i\ge k$ and $2t+2\le j-i+2k\le2t+1$ when $i<k$.  All
are impossible.  So the target carries $-$ only, and the $+$ line of the source
dies.
\end{proof}

\begin{remark}
Proposition \ref{prop:T012} is stated and proved only for arrows \emph{between
chain vertices}; its proof is entirely the interval combinatorics of
$S^{\eps}(i),S^{\eps}(j)$ with $i,j\le t$, and it says nothing about an arrow
out of, or into, a leaf.  The one leaf arrow needed later, namely the
injectivity of $a_{l\to t}$ on a leaf slot, is Lemma \ref{lem:R} verbatim: at a
leaf $l$ the only neighbour is $t$, so the map $Y$ of Lemma \ref{lem:R}
\emph{is} $a_{l\to t}$, and the leaf support tops out at $t+k-1\le 2t-1$,
strictly below the top degree $2t$, so the hypothesis of Lemma \ref{lem:R} holds
at every leaf slot without exception.  Since a leaf slot is at most a line
(Lemma \ref{lem:P}), the image of $a_{l\to t}$ on a nonzero leaf slot is
a line.
\end{remark}

\section{Removing the word: the statement \texorpdfstring{$(*^\circ)$}{(*o)}
and its representatives}\label{sec:reps}

\emph{This section is characteristic-free.}  The hypothesis
$\operatorname{char}\ne2$ declared at the head of \S\ref{sec:strands} is
suspended here and is used nowhere in this section: Lemma \ref{lem:Jinv} rests
on Theorem \ref{thm:main}, Corollary \ref{cor:reduce} on Lemma \ref{lem:Jinv},
Proposition \ref{prop:reps} is sign-vector combinatorics in the Weyl orbit, and
\eqref{eq:Brig} is Corollary \ref{cor:pol} --- all proved over an arbitrary
field.  This matters: \S\ref{sec:charp} quotes Definition \ref{def:star},
Corollary \ref{cor:reduce}, Proposition \ref{prop:reps}, Definition
\ref{def:B} and \eqref{eq:Brig} at $p=2$.  The hypothesis resumes at the head
of \S\ref{sec:families}.

The colour-$t$ step of the induction of \S\ref{sec:induction} needs to know
that the partial module is an up-set at the chain neighbour $t-1$.  Since the
partial module is itself extremal (Corollary \ref{cor:prefix}), the requirement
can be stated with no reference to $w$ or to any order of construction.

\begin{definition}\label{def:star}
$(*^\circ)$ is the statement: \emph{every extremal submodule $M\subseteq I_\ka$
with $\langle\al_t^\vee,\la-\dimv M\rangle\ge1$ is an up-set at $t-1$.}
\end{definition}

The side condition is essential and cannot be dropped: there exist extremal
modules that are not up-sets at $t-1$, and every one of them has
$\langle\al_t^\vee,\la-\dimv M\rangle=0$ (Proposition \ref{prop:sharp}; the
existence is exhibited in \S\ref{sec:verification}).  So any reduction of
$(*^\circ)$ to special weights must preserve the side condition.

\subsection{The parabolic}
For $\mu\in W\la$ write $c_t(\mu)=\langle\al_t^\vee,\mu\rangle$.  In $D_n$ the
vertices adjacent to $t=n-2$ are exactly $t-1$, $n-1$ and $n$.  Put
\[
   J'\;=\;I\setminus\{t-1,\,t,\,n-1,\,n\}\;=\;\{1,2,\dots,t-2\}.
\]

\begin{lemma}\label{lem:Jinv}
Let $\mu\in W\la$.
\begin{enumerate}
\item[\textup{(a)}] For $j\in J'$ we have $c_t(s_j\mu)=c_t(\mu)$.
\item[\textup{(b)}] For \emph{any} $j\in I$ that is a descent of $\mu$
\textup{(}i.e.\ $\langle\al_j^\vee,\mu\rangle\le-1$\textup{)}, the modules
$V(\mu)$ and $V(s_j\mu)$ have the same graded profile at every vertex other
than $j$.  In particular, for $j\in J'$ they agree at $t-1$ and at $t$, and for
$j=t$ they agree at $t-1$.
\end{enumerate}
\end{lemma}

\begin{proof}
(a) For $j\in J'$ the simple roots $\al_j$ and $\al_t$ are orthogonal, so
$s_j\al_t^\vee=\al_t^\vee$ and $c_t(s_j\mu)=\langle s_j\al_t^\vee,\mu\rangle
=c_t(\mu)$.  (b) Let $\nu=s_j\mu$; since $j$ is a
descent, $\langle\al_j^\vee,\nu\rangle\ge1$ and $\nu$ is shorter, and by
Theorem \ref{thm:main} $V(\mu)=\soc_{(j)}V(\nu)$.  A socle step of colour $j$
enlarges a submodule only at vertex $j$.  This argument is independent of $j$.
For $j\in J'$ one has $t-1\notin J'$ and $t\notin J'$, so $j\ne t-1,t$; and
$t\ne t-1$.
\end{proof}

\begin{remark}
This is the only choice of parabolic that works.  For the larger parabolic
$J=I\setminus\{t-1,t\}$, which also leaves the profile at $t-1$ unchanged, the
side condition is \emph{not} invariant: for a descent $j\in\{n-1,n\}$, which is
adjacent to $t$,
$c_t(s_j\mu)=c_t(\mu)+\langle\al_j^\vee,\mu\rangle\le c_t(\mu)-1$, so the
$W_J$-minimal representative of an orbit can fall outside the hypothesis set
--- and does so for the large majority of instances
(\S\ref{sec:verification}).  With $J'$ the side condition is literally
invariant and the reduction is closed by construction.
\end{remark}

\begin{corollary}\label{cor:reduce}
$(*^\circ)$ holds for all $\mu$ if and only if it holds for every
$W_{J'}$-minimal $\mu\in W\la$ with $c_t(\mu)\ge1$.
\end{corollary}

\begin{proof}
Let $\mu$ satisfy $c_t(\mu)\ge1$ and let $\mu_0$ be the $W_{J'}$-minimal
element of $W_{J'}\mu$.  Any chain $\mu\to s_{j_1}\mu\to\cdots\to\mu_0$ of
descents $j_i\in J'$ keeps $c_t\ge1$ and keeps the profile at $t-1$ unchanged,
by Lemma \ref{lem:Jinv}.  Being an up-set at $t-1$ depends only on that
profile.
\end{proof}

\subsection{Sharpness of the side condition}\label{ss:sharp}

$(*^\circ)$ says that $c_t\ge1$ forces the up-set property at $t-1$.  The
converse-shaped statement --- that the side condition is exactly right, i.e.\
that the extremal modules which fail to be up-sets at $t-1$ are confined to the
single value $c_t=0$ --- is also true, and is a short consequence of
$(*^\circ)$ itself together with Lemma \ref{lem:Jinv}(b).

\begin{proposition}[sharpness]\label{prop:sharp}
Granted $(*^\circ)$ \textup{(}proved in \S\ref{sec:closure}\textup{)}: every
extremal submodule $M\subseteq I_\ka$ that is \emph{not} an up-set at $t-1$
satisfies $\langle\al_t^\vee,\la-\dimv M\rangle=0$.
\end{proposition}

\begin{proof}
Put $\mu=\la-\dimv M\in W\la$; by Corollary \ref{cor:pol} the extremal
submodule with that dimension vector is unique, so $M=V(\mu)$, and
$c_t(M)=c_t(\mu)$.  If $c_t(\mu)\ge1$ then $(*^\circ)$ makes $M$ an up-set at
$t-1$, contrary to hypothesis; so $c_t(\mu)\le0$.

Suppose $c_t(\mu)\le-1$.  Then $t$ is a descent of $\mu$, so Lemma
\ref{lem:Jinv}(b) applies with $j=t$: writing $\nu=s_t\mu\in W\la$, the modules
$V(\mu)$ and $V(\nu)$ have the same graded profile at every vertex other than
$t$, in particular at $t-1$.  But
$c_t(\nu)=\langle\al_t^\vee,s_t\mu\rangle=-\langle\al_t^\vee,\mu\rangle
=-c_t(\mu)\ge1$, and $V(\nu)$ is extremal, so $(*^\circ)$ makes $V(\nu)$ an
up-set at $t-1$.  Being an up-set at $t-1$ depends only on the graded profile
at $t-1$, so $V(\mu)=M$ is an up-set at $t-1$ --- again contrary to hypothesis.

Hence $c_t(\mu)=0$.
\end{proof}

\begin{remark}
The value $c_t\le-1$ is not vacuous: since
$c_t(\mu)=\eta_{n-2}-\eta_{n-1}$ in the coordinates of Proposition
\ref{prop:reps}, $c_t$ takes all five values $-2,\dots,2$ on $W\om_k$, and over
$D_4$--$D_8$ exactly $2565$ of the $7464$ extremal modules have $c_t\le-1$
(\S\ref{ss:checks}).  So Proposition \ref{prop:sharp} is not the triviality
``$c_t\ge0$ always'' --- that statement is \emph{false}, and $(*^\circ)$ alone
gives only $c_t\le0$.  The work is done by the reflection $s_t$,
which converts a negative $c_t$ into a positive one at no cost at $t-1$
precisely because the socle step of colour $t$ moves nothing at $t-1$.
\end{remark}

\subsection{The representatives}

\begin{proposition}\label{prop:reps}
Let $2\le k\le n-2$.  The $W_{J'}$-minimal $\mu\in W\om_k$ with $c_t(\mu)\ge1$
are exactly the weights
\begin{equation}\label{eq:reps}
   \mu(h;x,y,z)\;=\;e_1+\cdots+e_h\;-\;(e_{n-2-c}+\cdots+e_{n-3})
   \;+\;x\,e_{n-2}+y\,e_{n-1}+z\,e_n,
\end{equation}
where $(x,y,z)\in\{0,\pm1\}^3$ satisfies $x-y\ge1$, and $h,c\ge0$ with
$h+c=k-q$, $q:=\#\{\text{nonzero among }x,y,z\}$.  The admissible triples are
the nine
\[
\begin{array}{ccc}
(1,0,0) & (1,0,1) & (1,0,-1)\\
(1,-1,0) & (1,-1,1) & (1,-1,-1)\\
(0,-1,0) & (0,-1,1) & (0,-1,-1),
\end{array}
\]
with $q=1$ for two of them, $q=2$ for five and $q=3$ for two, and their number
is
\[
   \sum_q N_q\,(k-q+1)=2k+5(k-1)+2(k-2)=9(k-1),
\]
\emph{independently of $n$}, and with no exception at $k=n-2$.  The automorphism
$\tau$ acts by $z\mapsto-z$, pairing the nine triples into six $\tau$-classes.
\end{proposition}

\begin{proof}
For $k\le n-2$ one has
$W\om_k=\{\sum_{i\in S}\eta_ie_i:|S|=k,\ \eta_i=\pm1\}$, so a weight is a vector
$(\eta_1,\dots,\eta_n)$ with $\eta_i\in\{0,\pm1\}$ and exactly $k$ nonzero
entries.  The pairings are $\langle\al_j^\vee,\mu\rangle=\eta_j-\eta_{j+1}$ for
$j\le n-1$ and $\langle\al_n^\vee,\mu\rangle=\eta_{n-1}+\eta_n$.

$W_{J'}$ is generated by $s_1,\dots,s_{t-2}=s_{n-4}$, hence is the symmetric
group permuting the coordinates $\eta_1,\dots,\eta_{n-3}$ and fixing
$\eta_{n-2},\eta_{n-1},\eta_n$.  Minimality in the orbit means no descent in
$J'$, i.e.\ $\eta_j\ge\eta_{j+1}$ for $1\le j\le n-4$: the sequence
$(\eta_1,\dots,\eta_{n-3})$ is weakly decreasing, hence of the form
$1^h0^b(-1)^c$ with $h+b+c=n-3$.  That is precisely the displayed shape, with
$x=\eta_{n-2}$, $y=\eta_{n-1}$, $z=\eta_n$.  The side condition is
$c_t(\mu)=\eta_{n-2}-\eta_{n-1}=x-y\ge1$, which with $x,y\in\{0,\pm1\}$ leaves
$(x,y)\in\{(1,0),(1,-1),(0,-1)\}$ and $z$ free, i.e.\ the nine listed triples.
Counting nonzero entries gives $h+c+q=k$, so $h+c=k-q$ and $h$ ranges over
$0,\dots,k-q$, which is $k-q+1$ values; and $b=n-3-h-c=n-3-k+q\ge0$ since
$k\le n-2$ and $q\ge1$.  Summing over the $N_1=2$, $N_2=5$, $N_3=2$ triples
gives $9(k-1)$, uniformly in $n$.  When $k=n-2$ the value $b=0$ is attained and
the triple $(h;x,y,z)=(k-1;1,0,0)$ gives
$\mu=e_1+\cdots+e_{k}=\la$; this is a \emph{genuine} member of the list and
nothing is to be subtracted for it, since $\la$ is $W_{J'}$-minimal (its
Bourbaki coordinates $1^k0^{n-k}$ are weakly decreasing) and satisfies the side
condition, $c_t(\la)=\eta_{n-2}-\eta_{n-1}=[\,k=n-2\,]-0=1\ge1$; so nothing is
subtracted at $k=n-2$.
Finally $\tau$ fixes $e_i$ for $i\le n-1$ and
negates $e_n$, so it sends $\mu(h;x,y,z)$ to $\mu(h;x,y,-z)$.
\end{proof}

\begin{remark}\label{rem:mu_a}
The $\tau$-fixed family $(x,y,z)=(1,0,0)$ consists of the $k$ weights
$\mu_h=e_1+\cdots+e_h-(e_{n-2-c}+\cdots+e_{n-3})+e_{n-2}$, $h+c=k-1$.  With one exception these
are exactly the weights whose \emph{only} descent is $t-1$; the exception is
$k=n-2$, $h=k-1$, where $e_{n-2}=e_k$ and $\mu_{k-1}=e_1+\cdots+e_k=\la$, which
is dominant and has no descent at all.  (Proposition \ref{prop:reps} states the
list correctly in either case: the condition there is $W_{J'}$-minimality
together with $c_t(\mu)\ge1$, and $\la$ satisfies both when $k=n-2$.)  They are the
$W_J$-minimal weights for the larger parabolic $J=I\setminus\{t-1,t\}$, and
they are not enough: the reduction through $J$ is not closed, and the other
$8k-9$ representatives are genuinely new.  In particular their dimension
vectors are mostly leaf-unbalanced, so Proposition \ref{prop:tau} below does not
dispose of them.
\end{remark}

The statement we shall prove at the representatives is stronger than
$(*^\circ)$ requires.

\begin{definition}\label{def:B}
$(B)$ is the statement: \emph{for every $\mu$ of Proposition \ref{prop:reps},
the graded multiplicities of $V(\mu)$ are the greedy top truncation of those of
$I_\ka$ at every vertex.}
\end{definition}

A greedy top truncation is an up-set, so $(B)$ gives $(*^\circ)$ at the
representatives, and hence, by Corollary \ref{cor:reduce}, in general.  By
Corollary \ref{cor:pol}, $(B)$ is equivalent to:

\begin{equation}\label{eq:Brig}
   \text{for each representative } \mu, \text{ \emph{some} submodule of }
   I_\ka \text{ has the greedy top truncation as its graded profile.}
\end{equation}

Indeed, such a submodule has dimension vector $v=\la-\mu$, and by rigidity it
is then the unique one, namely $V(\mu)$.  \emph{This removes the socle
construction from the problem completely}: what remains is to exhibit a
subspace of $I_\ka$, closed under all arrows, whose graded profile is greedy
from the top.  That is done in \S\ref{sec:closure}.

\section{The nine families}\label{sec:families}

\emph{From here to the end of \S\ref{sec:closure} we again assume
$\operatorname{char}\mathbb{K}\ne2$}, as declared at the head of
\S\ref{sec:strands}: everything below rests on the strand decomposition.

Fix a representative $\mu=\mu(h;x,y,z)$ as in \eqref{eq:reps}, write
$q=\#\{\text{nonzero among }x,y,z\}$, $b=n-3-k+q$ (the number of zero
coordinates among $\eta_1,\dots,\eta_{n-3}$), $c=k-q-h$, and $v=\la-\mu$ in
root coordinates.  Everything below is controlled by the two integers
$q+x\bmod2$ and $z$.

\subsection{The dimension vector}

\begin{lemma}\label{lem:vclosed}
With the above notation,
\begin{equation}\label{eq:vclosed}
   v_j=\bigl(\min(j,k)-h\bigr)^{+}+\bigl(j-h-b\bigr)^{+}
   \quad (1\le j\le t-1),
   \qquad
   v_t=2k-2h-q-x,
\end{equation}
and
\begin{equation}\label{eq:vleaf}
   v_{n-1}=\tfrac{v_t-y+z}{2},\qquad v_n=\tfrac{v_t-y-z}{2} .
\end{equation}
In particular $v(h;x,y,-z)$ is $v(h;x,y,z)$ with the two leaf entries
exchanged, and every chain entry is independent of $z$.
\end{lemma}

\begin{proof}
Write $X_p=\la_p-\mu_p$.  From \eqref{eq:reps},
\[
   X_p=0\ (p\le h),\quad X_p=[p\le k]\ (h<p\le h+b),\quad
   X_p=[p\le k]+1\ (h+b<p\le n-3),
\]
and $X_{n-2}=[k=t]-x$, $X_{n-1}=-y$, $X_n=-z$.  Comparing coefficients of
$e_p$ in $\sum_jv_j\al_j$ gives $v_j=X_1+\cdots+X_j$ for $j\le t$, whence the
first formula of \eqref{eq:vclosed} for $j\le n-3=t-1$.  At $j=t$,
\[
   v_t=v_{t-1}+[k=t]-x .
\]
If $k\le t-1$ then $v_{t-1}=(k-h)+(t-1-h-b)=(k-h)+c=2k-2h-q$ (using
$t-1-h-b=k-h-q=c$), and $v_t=2k-2h-q-x$.  If $k=t$ then
$v_{t-1}=(t-1-h)+c=2k-2h-q-1$ and $v_t=v_{t-1}+1-x=2k-2h-q-x$.  For the leaves,
the coefficient of $e_{n-1}$ in $\sum_jv_j\al_j$ is $-v_t+v_{n-1}+v_n$ and that
of $e_n$ is $-v_{n-1}+v_n$; equating with $X_{n-1}=-y$, $X_n=-z$ gives
$v_{n-1}+v_n=v_t-y$ and $v_{n-1}-v_n=z$, i.e.\ \eqref{eq:vleaf}.
\end{proof}

\subsection{Occupation numbers}
For $1\le j\le t$ put
\[
   \ms_j=\min(j,k),\qquad o_j=\max(0,\ j+k-t-1)=\max(0,\ j-B),\quad B:=n-1-k,
   \qquad u_j=v_j-\ms_j .
\]
By Proposition \ref{prop:W2}, each strand has $\ms_j$ slots at $j$, of which
$o_j$ lie in the multiplicity-$2$ window $W_j$; above $W_j$ there are
$\ms_j-o_j$ minus-only slots and below it $\ms_j-o_j$ plus-only slots.  Taking
$v_j$ dimensions greedily from the top exhausts the minus-only slots first,
then fills window slots; writing $r_j:=v_j-(\ms_j-o_j)=u_j+o_j$, it fills
$\lfloor r_j/2\rfloor$ window slots completely and, if $r_j$ is odd and
$r_j<2o_j$, leaves exactly one window slot half filled.  Consequently the
numbers of occupied slots in the two strands are
\begin{equation}\label{eq:sform}
   s^{+}_j=\max\Bigl(0,\ u_j,\ \Bigl\lceil\frac{u_j+o_j}{2}\Bigr\rceil\Bigr),
   \qquad
   s^{-}_j=v_j-s^{+}_j ,
\end{equation}
and in each strand the occupied slots are the top ones.  We also set
\begin{equation}\label{eq:g}
   g_j:=s^{+}_j-u_j
   =\max\Bigl(0,\ -u_j,\ \Bigl\lceil\frac{o_j-u_j}{2}\Bigr\rceil\Bigr),
\end{equation}
the \emph{dual occupation}, obtained from $s^+$ by $u\mapsto-u$.  Then
identically $s^{-}_j=\ms_j-g_j$.

\begin{lemma}\label{lem:Deps}
Define, for every chain vertex $j$ and either strand,
\[
   E^{\eps}_j:=\hi^{\eps}_j-2s^{\eps}_j+2 .
\]
When $s^{\eps}_j>0$ this is the lowest occupied degree of strand $\eps$ at
$j$; when $s^{\eps}_j=0$ it equals $\hi^{\eps}_j+2$, which lies above the
whole support $S^{\eps}(j)$.  \emph{In either case, a slot $(d,j)$ of strand
$\eps$ is occupied if and only if $d\ge E^{\eps}_j$}, and
\[
   E^{+}_j=j+k-2u_j-2g_j,\qquad E^{-}_j=2t+2-j-k+2g_j .
\]
\end{lemma}

\begin{proof}
The occupied slots of strand $\eps$ at $j$ are its top $s^{\eps}_j$ slots,
i.e.\ the degrees $d\in S^{\eps}(j)$ with
$d\ge\hi^{\eps}_j-2(s^{\eps}_j-1)=E^{\eps}_j$; for $s^{\eps}_j=0$ this
condition is $d\ge\hi^{\eps}_j+2$, which no slot satisfies, as required.  The
formulas: $E^+_j=(j+k-2)-2s^+_j+2=j+k-2s^+_j=j+k-2u_j-2g_j$ since
$s^+_j=u_j+g_j$.  For the minus strand, $s^-_j=\ms_j-g_j$ and
$E^-_j=2t-|j-k|-2(\ms_j-g_j)+2$; if $j\le k$ this is
$2t-k+j-2j+2g_j+2=2t+2-j-k+2g_j$, and if $j\ge k$ it is
$2t-j+k-2k+2g_j+2=2t+2-j-k+2g_j$.  Note that no case distinction on
$s^{\eps}_j$ was made, so both formulas are valid unconditionally.
\end{proof}

\begin{proposition}[the closure criterion, strandwise]\label{prop:crit}
Let occupations $s^{\pm}_j$ ($0\le s^{\eps}_j\le \ms_j$) be given at every chain
vertex $j$, defining an occupied set of strand slots as in Lemma
\ref{lem:Deps}, and suppose that at every chain edge $j\to j+1$
\begin{equation}\label{eq:twobounds}
   \Delta s^{+}_j=\Delta u_j+\Delta g_j\in\{0,1\}
   \qquad\text{and}\qquad
   \Delta g_j\in\{0,1\},
\end{equation}
where $\Delta f_j:=f_{j+1}-f_j$.  Then, for either strand $\eps$, either
direction of the edge, and every degree $d$:
\begin{center}
\emph{if the $\eps$-slot at $(d,j)$ is occupied and the slot $(d+1,i)$
\textup{(}$i=j\pm1$\textup{)} carries $\eps$, then that slot is occupied too.}
\end{center}
Consequently the graded subspace which at each chain vertex is the direct sum
of the top $s^{+}_j$ slots of the $+$ strand and the top $s^{-}_j$ slots of
the $-$ strand is carried into itself by every chain arrow; and no obligation
runs from a \emph{taken} chain slot to an \emph{untaken} nonzero chain slot,
except possibly from a slot in which only one strand is occupied to a slot
carrying only the other strand.
\end{proposition}

\begin{proof}
By Lemma \ref{lem:Deps} the two bounds \eqref{eq:twobounds} say exactly
\[
   E^{+}_{j+1}-E^{+}_j=1-2\Delta s^{+}_j\in\{-1,1\},
   \qquad
   E^{-}_{j+1}-E^{-}_j=-1+2\Delta g_j\in\{-1,1\},
\]
so $|E^{\eps}_{j+1}-E^{\eps}_j|\le1$ for both strands.  Now let the
$\eps$-slot at $(d,j)$ be occupied, so $d\ge E^{\eps}_j$, and let $(d+1,i)$
carry $\eps$ with $i=j\pm1$.  Then
$d+1\ge E^{\eps}_j+1\ge E^{\eps}_i$, so the target is occupied.  The same
argument with the roles of $j$ and $i$ exchanged covers the other direction.

For the two consequences: by Corollary \ref{cor:W4} a chain arrow maps the
$\eps$-line of $(d,j)$ isomorphically onto that of $(d+1,i)$ when both slots
carry $\eps$, and to $0$ otherwise, so the displayed statement is precisely
closure under chain arrows.  Finally, if $(d,j)$ is taken in strand $\eps$ and
$(d+1,i)$ is a nonzero chain slot that carries $\eps$, the displayed statement
makes it taken; so an untaken nonzero target forces $(d+1,i)$ not to carry
$\eps$, which at a chain slot means it carries only $-\eps$.  If the source is
full it carries $-\eps$ as well and is occupied there, so the target is taken
after all.
\end{proof}

\begin{remark}[why this form, and why it applies to $T^{\,l}$]
\label{rem:critform}
Two points.  First, $E^{\eps}_j$ is now \emph{defined at every chain vertex},
including those where $s^{\eps}_j=0$, for which reading $E^{\eps}_j$ as
``the lowest occupied degree'' is ill-posed on an unoccupied strand, and the
statement of Proposition \ref{prop:crit} is asserted at every chain edge.
Lemma \ref{lem:Deps} shows that the two closed formulas hold unconditionally,
which is what makes the uniform statement legitimate.  Second, only the
\emph{sufficiency} of \eqref{eq:twobounds} is used anywhere in this paper, and
the conclusion as stated is a statement about \emph{which slots are occupied},
not about which line sits in a partial slot.  This matters because Theorem
\ref{thm:B} applies the conclusion to $T^{\,l}(v)$, which is \emph{not} a sum
of two strand top-segments; see Lemma \ref{lem:samepattern}.
\end{remark}

\begin{remark}
It is tempting, and wrong, to phrase the minus condition as
$|\Delta s^{-}_j|\le1$.  Since $\hi^{-}_j=2t-|j-k|$ is a tent rather than a
ramp, the correct condition is one-sided and changes side at $j=k$:
$\Delta s^-\in\{0,1\}$ for $j<k$ and $\Delta s^-\in\{-1,0\}$ for $j\ge k$.  The
symmetric bound $|\Delta s^-|\le1$ is strictly weaker on the wrong side of the
tent, and is not equivalent to up-closedness.  The formulation of Proposition
\ref{prop:crit} in terms of the thresholds $E^{\eps}_j$ of Lemma \ref{lem:Deps}
--- equivalently, of $\Delta s^{+}$ and $\Delta g$ --- is uniform in $\eps$ and
has no case split; the asymmetry lives entirely inside $\hi^{\eps}$.
\end{remark}

\subsection{The two increment bounds}

\begin{lemma}\label{lem:increments}
For every one of the nine families:
\begin{enumerate}
\item[\textup{(i)}] $\Delta u_j\in\{-1,0,1\}$ and $\Delta o_j\in\{0,1\}$ at
every chain edge $j\to j+1$, $1\le j\le t-1$;
\item[\textup{(ii)}] $\Delta u_j=-1$ occurs only for $j+1\le h$, or for
$j+1=t$ with $x=1$;
\item[\textup{(iii)}] $\Delta u_j=+1$ forces $j\ge h+b$;
\item[\textup{(iv)}] $u_{t-1}\le o_{t-1}+1$.
\end{enumerate}
\end{lemma}

\begin{proof}
For $j+1\le t-1$, \eqref{eq:vclosed} gives
$v_{j+1}-v_j=[\,h<j+1\le k\,]+[\,j+1>h+b\,]\in\{0,1,2\}$ and
$\ms_{j+1}-\ms_j=[\,j+1\le k\,]$, so
$\Delta u_j=[\,h<j+1\le k\,]+[\,j+1>h+b\,]-[\,j+1\le k\,]\in\{-1,0,1\}$;
it is $-1$ only when $[\,h<j+1\le k\,]=0=[\,j+1>h+b\,]$ and
$[\,j+1\le k\,]=1$, i.e.\ $j+1\le h$.  It is $+1$ only when
$[\,j+1>h+b\,]=1$, i.e.\ $j\ge h+b$.  At the last edge $j+1=t$,
Lemma \ref{lem:vclosed} gives $v_t-v_{t-1}=[k=t]-x$ and
$\ms_t-\ms_{t-1}=[\,t\le k\,]=[k=t]$, so $\Delta u_{t-1}=-x$.  That proves
(i)--(iii); $o_j=\max(0,j-B)$ is nondecreasing with steps in $\{0,1\}$.

For (iv): $o_{t-1}=\max(0,t-1+k-t-1)=k-2$ since $k\ge2$.  If $k\le t-1$ then
$u_{t-1}=v_{t-1}-k=(2k-2h-q)-k=k-2h-q$; if $k=t$ then
$u_{t-1}=v_{t-1}-(t-1)=(2k-2h-q-1)-(k-1)=k-2h-q$ as well.  So
$u_{t-1}-o_{t-1}=2-2h-q\le1$ because $q\ge1$.
\end{proof}

\begin{proposition}[increasing bound]\label{prop:incr}
$\Delta u_j+\Delta g_j=\Delta s^{+}_j\in\{0,1\}$ at every chain edge, for every
one of the nine families.
\end{proposition}

\begin{proof}
\emph{Upper bound.}  In \eqref{eq:sform} the three arguments are $0$ (constant),
$u_j$ (rising by $\Delta u\le1$) and $\lceil(u_j+o_j)/2\rceil$ (rising by
$\lceil(\Delta u+\Delta o)/2\rceil\le1$, since $\Delta u+\Delta o\le2$).  A
maximum of finitely many quantities each rising by at most $1$ rises by at most
$1$.

\emph{Lower bound.}  A maximum can only fall if some argument falls, so by
Lemma \ref{lem:increments} we need only examine the two exceptional steps.

\emph{(a) $j+1\le h$.}  Then $j<h\le k-q<k$, so $\ms_j=j$ and $v_j=0$, giving
$u_j=-j$.  Also $j\le h-1\le k-q-1$ and $B=n-1-k\ge1$, so
$j-B\le k-q-1-(n-1-k)=2k-q-n\le n-4-q<0$ and $o_j=0$.  Hence
$u_j+o_j=-j\le0$, so $s^{+}_j=0$; the same computation gives $s^{+}_{j+1}=0$.
No drop.

\emph{(b) $j+1=t$ with $x=1$.}  Here $\Delta u=-1$ and, since $k\ge2$,
$\Delta o=o_t-o_{t-1}=(k-1)-(k-2)=1$.  So $u_j+o_j$ is unchanged and the third
argument of \eqref{eq:sform} is unchanged; the first is constant; only $u$
falls.  Therefore $s^{+}$ can drop only if the maximum at $j=t-1$ was attained
\emph{strictly} by $u_{t-1}$, i.e.\
$u_{t-1}>\lceil(u_{t-1}+o_{t-1})/2\rceil$, which forces $u_{t-1}\ge o_{t-1}+2$.
This contradicts Lemma \ref{lem:increments}(iv).  No drop.

\emph{(c) $j+1=t$ with $x=0$.}  Then $\Delta u=0$ and $\Delta o=1$, so no
argument falls.
\end{proof}

\begin{proposition}[decreasing bound]\label{prop:decr}
$\Delta g_j\in\{0,1\}$ at every chain edge, for every one of the nine families.
\end{proposition}

\begin{proof}
Write $g^{1}_j=-u_j$ and $g^{2}_j=\lceil(o_j-u_j)/2\rceil$, so
$g_j=\max(0,\,g^{1}_j,\,g^{2}_j)$, and put
$\delta:=\Delta o-\Delta u$.

\emph{Upper bound.}  $g^{1}$ rises by $-\Delta u\le1$; $g^{2}$ rises
by $\lceil\delta/2\rceil\le1$ since $\delta\le 1-(-1)=2$; the third argument is
the constant $0$.

\emph{Lower bound.}  If $\Delta u\le0$ then $g^{1}$ does not fall and
$\delta\ge0$, so $g^{2}$ does not fall: $g$ does not fall.  So assume
$\Delta u=+1$; by Lemma \ref{lem:increments}(iii), $j\ge h+b$, and then
\eqref{eq:vclosed} gives $u_j=j-2h-b$ (indeed $j\ge h+b\ge h$ and $h\le k-q<k$,
so $(\min(j,k)-h)^+-\min(j,k)=-h$ and $(j-h-b)^+=j-h-b$).

\emph{Case $\Delta o=1$, i.e.\ $\delta=0$.}  Then $g^{2}$ is unchanged and the
third argument is constant, so $g$ can only fall if $g_j=g^{1}_j$ strictly
exceeds $\max(0,g^{2}_j)$.  Putting $U=-u_j$, this requires $U\ge1$ and
$U-1\ge\lceil(o_j+U)/2\rceil$, i.e.\ $U\ge o_j+2$, i.e.\ $u_j+o_j\le-2$.  We
show this is impossible.  $\Delta o=1$ means $j+1>B$, so $j\ge B$ and
$o_j=j-B\ge0$.  Since $B=n-1-k$ and $b=n-3-k+q$ we have $B=b+2-q$, and
therefore
\[
   u_j+o_j=(j-2h-b)+(j-B)=2j-2h-b-B\ \ge\ 2(h+b)-2h-b-B=b-B=q-2\ \ge\ -1,
\]
using $j\ge h+b$ and $q\ge1$.  So $u_j+o_j\ge-1>-2$.

\emph{Case $\Delta o=0$.}  Since $o$ is nondecreasing with
$o_{j+1}=\max(0,j+1-B)$, $\Delta o=0$ forces $o_{j+1}=0$, hence $o_j=0$ and
$j+1\le B$.  Combined with $j\ge h+b=h+B+q-2$ this gives $h+q\le1$, so
$q=1$, $h=0$ and $j=B-1=b$.  Then $u_j=j-2h-b=0=o_j$, so $g_j=0$; and
$u_{j+1}=1$, $o_{j+1}=0$, so
$g_{j+1}=\max(0,-1,\lceil-1/2\rceil)=0$.  No fall.

\emph{The last edge $j=t-1\to t$.}  There $\Delta u=-x\le0$, so the first case
of the lower bound applies.
\end{proof}

\subsection{The node and the leaves}

\begin{lemma}\label{lem:node}
Let $E_t$ denote the lowest \emph{taken} degree at $t$ in the greedy top
truncation to total $v_t>0$, and $E_l=t+k+1-2v_l$ the lowest taken degree at a
leaf $l$ with $v_l>0$.  Put $s:=q+x$.  Then
\begin{align}
   E_t&=t-k+2h+q+x+[\,q+x \text{ odd}\,],\label{eq:Dt}\\
   E_{n-1}-E_t&=y-z+[\,q+x \text{ even}\,],
   \qquad
   E_{n}-E_t=y+z+[\,q+x \text{ even}\,],\label{eq:offsets}
\end{align}
and the slot $(E_t,t)$ is \emph{partial} \textup{(}taken to dimension $1$ out
of $2$\textup{)} if and only if $v_t$ is even, if and only if $q+x$ is even.
Moreover the two offsets in \eqref{eq:offsets} are always $+1$ or $-1$, never
anything else, and a leaf offset $-1$ occurs exactly when $z\ne0$, or when
$y=-1$ and $q+x$ is odd.
\end{lemma}

\begin{proof}
By Lemma \ref{lem:P} the node column has one $1$-dimensional slot at degree
$t+k$, then $k-1$ two-dimensional slots at $t+k-2,\dots,t-k+2$, then one
$1$-dimensional slot at $t-k$.  Taking $v_t$ dimensions from the top, the
lowest taken degree is $t+k-2\lceil(v_t-1)/2\rceil$, and the last slot reached
is partial exactly when $v_t$ is even and $v_t<2k$, i.e.\ when the slot reached
is two-dimensional.  By Lemma \ref{lem:vclosed}, $v_t=2k-2h-s$, so $v_t\equiv s
\pmod 2$, and $v_t<2k$ because $2h+s\ge q\ge1$.  If $s$ is even then
$\lceil(v_t-1)/2\rceil=v_t/2$ and $E_t=t+k-v_t=t-k+2h+s$; if $s$ is odd then
$\lceil(v_t-1)/2\rceil=(v_t-1)/2$ and $E_t=t+k-v_t+1=t-k+2h+s+1$.  That is
\eqref{eq:Dt}.

For the leaves, Lemma \ref{lem:P} gives the leaf support
$t-k+1,\dots,t+k-1$, whence $E_l=t+k+1-2v_l$, and by \eqref{eq:vleaf}
\[
   E_{n-1}=t+k+1-(v_t-y+z)=t-k+2h+s+1+y-z,
\]
so $E_{n-1}-E_t=1+y-z-[\,s\text{ odd}\,]=y-z+[\,s\text{ even}\,]$; the formula
for $E_n$ is the same with $z\mapsto-z$.

Finally, the nine families satisfy $x-y\ge1$ with $x\le1$, so $y\le0$ and
$y\in\{0,-1\}$, and $z\in\{0,\pm1\}$.  Enumerating the nine triples:
\[
\begin{array}{llll}
(1,0,0):\ s=2,& (+1,+1);\qquad
&(1,0,\pm1):\ s=3,& (\mp1,\pm1);\\
(1,-1,0):\ s=3,& (-1,-1);
&(1,-1,\pm1):\ s=4,& (\mp1,\pm1);\\
(0,-1,0):\ s=1,& (-1,-1);
&(0,-1,\pm1):\ s=2,& (\mp1,\pm1).
\end{array}
\]
All eighteen values are $\pm1$, and the pattern of $-1$'s is as stated.
\end{proof}

\begin{proposition}[the partial-slot closed form]\label{prop:PS}
At every chain vertex $j$ with $1\le j\le t$ the greedy top truncation has
\emph{at most one} partial slot.  It exists if and only if
\[
   |u_j|<o_j \quad\text{and}\quad u_j+o_j \text{ is odd},
\]
and then its degree is
\begin{equation}\label{eq:PS}
   d_j\;=\;t-u_j .
\end{equation}
No partial slot occurs at a leaf.  This statement is characteristic-free, being
a statement about the greedy truncation of the multiplicity profile alone.
\end{proposition}

\begin{proof}
The window $W_j=[\,2t+2-j-k,\ j+k-2\,]$ has $o_j$ slots, sitting below the
$\ms_j-o_j$ minus-only slots and above the $\ms_j-o_j$ plus-only slots.  Greedy
from the top exhausts the minus-only slots first; writing
$r=v_j-(\ms_j-o_j)=u_j+o_j$ for what is left, it fills $\lfloor r/2\rfloor$
window slots completely, and leaves one window slot half filled precisely when
$r$ is odd and $0<r<2o_j$, i.e.\ $u_j>-o_j$ and $u_j<o_j$.  That slot is the
$\frac{r+1}{2}$-th from the top of $W_j$, of degree
\[
   \max(W_j)-2\Bigl(\frac{r+1}{2}-1\Bigr)=(j+k-2)-(r-1)=j+k-1-r .
\]
Existence forces $o_j>0$, so $o_j=j+k-t-1$ and $j+k-1-r=j+k-1-u_j-(j+k-t-1)
=t-u_j$.  A leaf slot is at most a line by Lemma \ref{lem:P}, hence taken
wholly or not at all.
\end{proof}

\begin{corollary}\label{cor:parity}
Along the region $j\ge\max(B,h+b)$ one has $u_j=j-2h-b$ and $o_j=j-B$, so
$u_j+o_j=2j-2h-b-B\equiv b+B\equiv q\pmod2$.  Hence the parity of $r_j$ is
constant and equal to that of $q$: if $q$ is even there is no partial slot at
any chain vertex $j\le t-1$, and if $q$ is odd the partial chain vertices form
a contiguous interval ending at $t-1$.  Outside that region no partial slot can
occur: for $j<h+b$ one would need $B<b$, i.e.\ $q>2$, and then
$h+B=h+b-1$ leaves no integer $j$ with $h+B<j<h+b$; and for $j<h$ the column is
empty.
\end{corollary}

\section{Closure of the truncation}\label{sec:closure}

We must exhibit, for each representative $\mu$, a graded submodule of $I_\ka$
whose profile is the greedy top truncation; by \eqref{eq:Brig} this proves
$(B)$ and hence $(*^\circ)$.

A greedy top truncation is not determined by its multiplicities: at a partial
slot a \emph{line} in a two-dimensional space must be chosen.  We therefore
define two candidate subspaces.

\begin{definition}\label{def:Tv}
Let $v$ be the dimension vector of a representative.  The \emph{canonical top
truncation} $T(v)$ is the graded subspace of $I_\ka$ which at every vertex is
full above the greedy threshold, zero below it, and, at a partial slot, is the
$\tau$-eigenline of that slot.  When some leaf $l$ has $E_l=E_t-1$, the
\emph{leaf-pinned truncation} $T^{\,l}(v)$ is defined in the same way except
that at $(E_t,t)$ it carries the line $\im(a_{l\to t})$, and at each further
partial slot it carries the transport of that line along the path of
Proposition \ref{prop:L1} below.  (This second clause presupposes that
$(E_t,t)$ is itself partial; by Lemma \ref{lem:node} that holds, together with
the existence of such a leaf, exactly in the four families named below, and
$T^{\,l}(v)$ is defined only there.)  In both cases $\dimv T=v$ by
construction.
\end{definition}

By Lemma \ref{lem:node} a leaf with $E_l=E_t-1$ exists, \emph{and the node slot
$(E_t,t)$ is partial}, exactly in the four families with $q+x$ even and
$z\ne0$, namely $(0,-1,\pm1)$ and $(1,-1,\pm1)$.  Note that $\im(a_{l\to t})$ is
a line, by Lemma \ref{lem:R} at the leaf (see the Remark after Proposition
\ref{prop:T012}), and that it is \emph{not} a $\tau$-eigenvector: by
Proposition \ref{prop:W4l} the folded leaf map is nonzero on both eigen-parts,
so the image of the one-dimensional $(I_\ka)_l[E_t-1]$ has nonzero component in
each eigenline.  In those four families $T(v)$ is therefore \emph{not} a
submodule and the pinned version is forced.

\begin{lemma}[$T$ and $T^{\,l}$ have the same taken pattern]\label{lem:samepattern}
$T(v)$ and $T^{\,l}(v)$ have the same graded dimension at every slot, namely
the greedy top truncation of the profile of $I_\ka$ to total $v_j$ at each
vertex $j$.  In particular a slot is full, partial or untaken for $T(v)$ if and
only if it is so for $T^{\,l}(v)$; the two subspaces differ only in \emph{which}
line occupies each partial slot.
\end{lemma}

\begin{proof}
Immediate from Definition \ref{def:Tv}: both are ``full above the greedy
threshold, zero below it, one line at a partial slot'', and the threshold and
the set of partial slots depend only on the multiplicity profile of $I_\ka$
and on $v$, not on the choice of line.  A line has dimension $1$ in either
case.
\end{proof}

Consequently every statement about the \emph{taken/untaken} pattern of $T(v)$
--- in particular every consequence of Proposition \ref{prop:crit}, which by
Remark \ref{rem:critform} is such a statement --- holds verbatim for
$T^{\,l}(v)$.  We use this silently below, and it is the only way Proposition
\ref{prop:crit} is ever applied in the four pinned families, where the
hypothesis ``a sum of two strand top-segments'' fails for $T^{\,l}(v)$ itself
but holds for the auxiliary $T(v)$ with the same occupations $s^{\pm}_j$.

\subsection{The obligation census}
Call an \emph{obligation} a pair consisting of a taken slot $(d,j)$ and a slot
$(d+1,i)$ with $i\sim j$; closure of $T$ means that $a_{j\to i}$ maps the taken
part of $(d,j)$ into the taken part of $(d+1,i)$.  We classify obligations by
the state of source and target: \emph{full}, \emph{line} (a partial slot), or
\emph{empty}.

\begin{proposition}\label{prop:partials}
For the nine families the partial slots are as follows.  \textup{(}The node
entry $(E_t,t)$ presupposes $v_t>0$, so that the $E_t$ of Lemma
\ref{lem:node} exists; this holds at every representative except the single
degenerate one $(x,y,z)=(1,0,0)$ with $h=k-1$, which is treated at the end of
the proof.\textup{)}
\[
\renewcommand{\arraystretch}{1.2}
\begin{array}{lccl}
\text{family} & q & q+x & \text{partial slots}\\\hline
(1,0,0) & 1 & \text{even} & \text{a chain interval ending at } t-1,\ \text{and } (E_t,t)\\
(0,-1,0) & 1 & \text{odd} & \text{a chain interval ending at } t-1;\ (E_t,t)\text{ is full}\\
(1,-1,0) & 2 & \text{odd} & \text{none}\\
(1,0,\pm1) & 2 & \text{odd} & \text{none}\\
(0,-1,\pm1) & 2 & \text{even} & \text{exactly one: } (E_t,t)\\
(1,-1,\pm1) & 3 & \text{even} & \text{a chain interval ending at } t-1,\ \text{and } (E_t,t)
\end{array}
\]
\end{proposition}

\begin{proof}
Combine Corollary \ref{cor:parity} (chain partials occur iff $q$ is odd, and
then form an interval ending at $t-1$) with Lemma \ref{lem:node} (the node slot
is partial iff $q+x$ is even).

For $(1,0,0)$ and $(0,-1,0)$, where $q=1$, the interval may be empty, and we
record exactly when.  Here $b=n-2-k$ and $B=b+1$, and for
$j\ge\max(B,h+b)$ Corollary \ref{cor:parity} gives $u_j=j-2h-b$ and
$o_j=j-B=j-b-1$, so $r_j=2(j-h-b)-1$.  Thus $r_j>0$ iff $j>h+b$, and
$r_j<2o_j$ reads $-2h<1$, i.e.\ $h\ge1$.  So the partial vertices are exactly
the $j$ with $h+b+1\le j\le t-1$ when $h\ge1$, and there are none when $h=0$;
and since $h+b+1=h+n-1-k$ and $t-1=n-3$, the range is empty also when
$h\ge k-1$.  As $h\le k-q=k-1$, the conclusion is:
\begin{center}
\emph{the chain interval is empty exactly when $h=0$ or $h=k-1$; in
particular it is always empty when $k=2$.}
\end{center}
(The statement in the table is that the partial slots form a chain interval
ending at $t-1$, possibly empty, and that is what is used below; no later
argument depends on the interval being nonempty in these two families.)

The node entry of the table needs the same care, and only in the family
$(1,0,0)$.  By Lemma \ref{lem:vclosed}, $v_t=2k-2h-q-x$ is decreasing in $h$,
so its least value over the range $0\le h\le k-q$ of Proposition
\ref{prop:reps} is attained at $h=k-q$ and equals $q-x$; running over the nine
triples, $q-x$ is $0$ for $(1,0,0)$, $1$ for $(0,-1,0)$, $(1,0,\pm1)$ and
$(1,-1,0)$, and $2$ for $(0,-1,\pm1)$ and $(1,-1,\pm1)$.  Hence $v_t>0$ at
every representative save one: the family $(1,0,0)$ at the top of its range,
$h=k-1$, where $v_t=0$.  There the hypothesis $v_t>0$ of Lemma \ref{lem:node}
fails --- the colour-$t$ column is untaken in every degree, $E_t$ does not
exist, and $\mu=e_1+\cdots+e_{k-1}+e_{n-2}$ has $v=\la-\mu=\al_k+\cdots+
\al_{t-1}$ supported on the chain alone, with $v_t=v_{n-1}=v_n=0$.  The table
row is to be read accordingly: at $h=k-1$ there is no partial slot at all,
neither on the chain (previous paragraph) nor at the node.  Nothing below is
affected.  The only consumer of this proposition is Proposition
\ref{prop:L5}(a), which concerns $(0,-1,\pm1)$, where $h\le k-2$ and so
$v_t=2k-2h-2\ge2$; and in Theorem \ref{thm:B} the family $(1,0,0)$ is
discharged entirely through Proposition \ref{prop:crit} (classes (ii), (iii)
and (vi)) and Proposition \ref{prop:leafclasses}, neither of which mentions
$E_t$ --- the latter derives $v_l>0$, respectively $v_t>0$, from the source of
the obligation being taken, so that at $v_t=0$ its two cases are vacuous ---
while class (v) requires $z\ne0$.
\end{proof}

\begin{theorem}\label{thm:B}
$(B)$ holds: for every representative $\mu(h;x,y,z)$ of Proposition
\ref{prop:reps} there is a graded submodule of $I_\ka$ of dimension vector
$v=\la-\mu$ whose profile at every vertex is the greedy top truncation, namely
$T(v)$ in the five families with $q+x$ odd or $z=0$, and $T^{\,l}(v)$ in the
four families with $q+x$ even and $z\ne0$.  Consequently $V(\mu)$ is that
module, and $(*^\circ)$ holds.
\end{theorem}

\begin{proof}
Every obligation of $T$ falls into one of six classes, and each class has its
own discharge.

\smallskip\noindent
\textbf{(i) Target full.}  Nothing to check.  This covers all full$\to$full and
line$\to$full obligations, at chain arrows and at leaf arrows alike.

\smallskip\noindent
\textbf{(ii) full$\to$partial on the chain, in the five families with $q+x$
odd or $z=0$.}  In these five families the module built is $T(v)$ and every
partial slot carries the $\tau$-eigenline.  Here the source slot is
$1$-dimensional and taken, the target $2$-dimensional and half taken.  By (T0)
the arrow has rank $1$, and by (T1) its image is the $\tau$-eigenline of the
target --- which is what $T(v)$ carries there.  That the thresholds line up so
that no other configuration arises is the strandwise statement of Proposition
\ref{prop:crit}, i.e.\ the two increment bounds \eqref{eq:twobounds},
established in Propositions \ref{prop:incr} and \ref{prop:decr}.

\emph{In the remaining four families --- those with $q+x$ even and $z\ne0$,
where the module built is $T^{\,l}(v)$ and the partial slots carry the pinned
line $\im(a_{l\to t})$ and its transports, which by \textup{(N6)} is provably
\emph{not} a $\tau$-eigenvector --- the discharge just given is unavailable,
and it is the negation of what $T^{\,l}(v)$ carries.  In those four families
this class is instead \emph{empty}: no full$\to$partial obligation occurs
along a chain arrow at all.  That is Proposition \ref{prop:L5}
below, and it is what makes Theorem \textup{\ref{thm:B}} true there.}

\smallskip\noindent
\textbf{(iii) line$\to$partial on the chain, for the eigenline families
$(1,0,0)$ and $(0,-1,0)$.}  Here every partial slot carries the $+$ line, so
$T(v)$ is at every chain vertex the direct sum of a top segment of the $+$
strand and a top segment of the $-$ strand, and Proposition \ref{prop:crit}
applies verbatim: closure follows from \eqref{eq:twobounds}, which holds by
Propositions \ref{prop:incr} and \ref{prop:decr}.

\smallskip\noindent
\textbf{(iv) line$\to$partial on the chain, for $(1,-1,\pm1)$.}  Here the
partial slots carry the transported line, which is not an eigenvector, so
$T^{\,l}(v)$ is \emph{not} a sum of two strand parts and (iii) does not apply.
This class is discharged by Propositions \ref{prop:L1}, \ref{prop:L3} and
\ref{prop:L4} below: the partial slots form a path with $(E_t,t)$ an endpoint,
these obligations are exactly its edges, and each edge has rank $2$ by (T0),
hence is an isomorphism, so the pinned line propagates uniquely along the path.

\smallskip\noindent
\textbf{(v) the single full$\to$partial leaf obligation.}  This is
$a_{l\to t}:(E_t-1,l)\to(E_t,t)$ and occurs exactly in the four families with
$q+x$ even and $z\ne0$ (Lemma \ref{lem:node}).  Its image is a line by Lemma
\ref{lem:R} at the leaf, and $T^{\,l}(v)$ is defined to carry precisely that
line; so the obligation holds by construction.

For $(0,-1,\pm1)$ that is the only partial slot at all, so no further
propagation is needed.  The line at $(E_t,t)$ has \emph{three} obligations out
of it, one for each neighbour of $t$: to $(E_t+1,t-1)$, to $(E_t+1,l)$ and to
$(E_t+1,l')$, where $l'$ is the other leaf.  All three targets are full or
zero, hence all three fall in class (i).  Indeed, by Proposition
\ref{prop:L5}(a) the lowest taken degree at $t-1$ is $E_t+1$ and $t-1$ carries
no partial slot, so $(E_t+1,t-1)$ is full or zero; and at each leaf $E_{l}\le
E_t+1$ by Lemma \ref{lem:node}, so a nonzero leaf slot at $E_t+1$ is taken,
hence full (a leaf slot is at most a line).  (For the leaf $l$ itself one may
also see the obligation vanish outright: $a_{t\to l}a_{l\to t}=0$ is the
preprojective relation at $l$, whose only neighbour is $t$.)

\smallskip\noindent
\textbf{(vi) line$\to$empty on the chain.}  Let the source be the partial slot
$(d,j)$, carrying by Proposition \ref{prop:PS} the $+$ line (the lowest
occupied slot of the $+$ strand, by Lemma \ref{lem:Deps}), and let the target
$(d+1,i)$, $i=j\pm1$, be a nonzero untaken chain slot.

We first show the target is $1$-dimensional.  Suppose it were
$2$-dimensional; then it carries $+$.  If $s^{+}_i>0$ then, by the proof of
Proposition \ref{prop:crit}, $|E^{+}_i-E^{+}_j|=1$ and $E^{+}_j=d$, so
$E^{+}_i\le d+1$ and the target is occupied in the $+$ strand, hence taken ---
a contradiction.  So $s^{+}_i=0$.  If $i=j+1$ this gives
$\Delta s^{+}_j=-s^{+}_j\le-1$, contradicting Proposition \ref{prop:incr}.  If
$i=j-1$, use Proposition \ref{prop:PS}: the partial slot at $j$ has degree
$d=t-u_j$ and lies in the window, so $o_j=j+k-t-1$ and $s^{+}_j=\lceil
r_j/2\rceil$ with $r_j=u_j+o_j$; the hypothesis that $(d+1,i)$ carries $+$
reads $d+1\le\hi^{+}_{j-1}=j+k-3$, i.e.\ $t-u_j+1\le o_j+t-2$, i.e.\
$r_j\ge3$, whence $s^{+}_j\ge2$ and $\Delta s^{+}_{j-1}=s^{+}_j-s^{+}_{j-1}
=s^{+}_j\ge2$, again contradicting Proposition \ref{prop:incr}.  So the target
is $1$-dimensional.

Now (T2) applies: the source is a $2$-dimensional chain slot, the target a
$1$-dimensional one, and (T2) says the arrow annihilates the $\tau$-eigenline
of the source, which is the line $T(v)$ carries.  So the obligation is
discharged.  (This class occurs only in the two families $(1,0,0)$ and
$(0,-1,0)$, where the module is $T(v)$; for $(1,-1,\pm1)$ there is no
line$\to$empty obligation at all, by Proposition \ref{prop:L4}.)

\smallskip
It remains to show that \emph{no other class occurs}: no full$\to$empty
obligation anywhere, and no line$\to$partial or line$\to$empty obligation at a
leaf.  On the chain, full$\to$empty is the last consequence recorded in
Proposition \ref{prop:crit}: a full source slot has both strands occupied, so
whichever strand the nonzero target carries, that strand is occupied there
too.  This is a statement about occupations only, so by Lemma
\ref{lem:samepattern} it holds for $T^{\,l}(v)$ as well as for $T(v)$; the
occupations satisfy \eqref{eq:twobounds} by Propositions \ref{prop:incr} and
\ref{prop:decr} in all nine families.  For $(1,-1,\pm1)$ the corresponding
conclusion for the transported partial slots is Proposition \ref{prop:L4}, and
the absence of full$\to$partial chain obligations in all four pinned families
is Proposition \ref{prop:L5}.  The three leaf classes are Proposition
\ref{prop:leafclasses}.

Thus $T(v)$, respectively $T^{\,l}(v)$, is closed under every arrow, i.e.\ is a
graded submodule, with $\dimv T=v$.  By Corollary \ref{cor:pol} it is the unique
submodule of that dimension vector, hence equals $V(\mu)$; and its profile is
greedy from the top at every vertex, which is $(B)$.  By
Corollary \ref{cor:reduce}, $(*^\circ)$ follows.
\end{proof}

\subsection{The leaf classes}

\begin{proposition}\label{prop:leafclasses}
For every family instance $(n,k,h;x,y,z)$ the greedy truncation has
\textup{(F1)} no line$\to$partial obligation at a leaf, \textup{(F2)} no
full$\to$empty obligation at a leaf, and \textup{(F3)} no line$\to$empty
obligation at a leaf.
\end{proposition}

\begin{proof}
An obligation is ``at a leaf'' if one end of its arrow is a leaf slot.

\emph{(F1).}  A partial target has dimension $2$, so by Lemma \ref{lem:P} it is
not a leaf slot; hence the target is the node and the \emph{source} is a leaf
slot, which has dimension $\le1$ and is therefore full or empty, never a line.

\emph{(F2), (F3).}  Two directions.

\emph{Source at the leaf}, $a_{l\to t}:(d,l)\to(d+1,t)$ with the source taken.
Then $v_l>0$ and $d\ge E_l$; by Lemma \ref{lem:node} $E_l\ge E_t-1$, so
$d+1\ge E_t$.  The node column takes every nonzero slot of degree $\ge E_t$, so
the target is taken: full or partial, never empty.

\emph{Source at the node}, $a_{t\to l}:(d,t)\to(d+1,l)$ with the source taken,
so $v_t>0$ and $d\ge E_t$.  If $v_l>0$, the leaf column takes every nonzero
slot of degree $\ge E_l$, and $E_l\le E_t+1\le d+1$ by Lemma \ref{lem:node}; so
the target is taken, hence full, never empty.  If $v_l=0$, then by
\eqref{eq:vleaf} either $v_t=y-z$ or $v_t=y+z$; the nine families have
$x-y\ge1$ with $x\le1$, so $y\le0$, and $z\in\{0,\pm1\}$, whence $v_t\le1$.
Since the source is taken, $v_t=1$, so $E_t=t+k$ is the top node degree, the
only one taken, and its leaf target sits at degree $t+k+1$, beyond the leaf top
degree $t+k-1$: there is no obligation at all.
\end{proof}

\subsection{The path of partial slots for
\texorpdfstring{$(1,-1,\pm1)$}{(1,-1,+-1)}}

Throughout this subsection $(x,y,z)=(1,-1,z)$ with $z=\pm1$, so $q=3$, $x=1$,
$b=n-k$, $B=b-1$, $h\le k-3$, and by Lemma \ref{lem:vclosed}
$u_j=\max(-\min(j,h),\ j-2h-b)$ for $j\le t-1$.

\begin{lemma}\label{lem:partialset}
The partial slots of the greedy truncation are exactly
\[
   \mathcal{P}=\bigl\{(d_j,\,j)\ :\ h+b\le j\le t-1\bigr\}\ \cup\
   \bigl\{(E_t,\,t)\bigr\},
   \qquad d_j=t+2h+b-j,
\]
and $E_t=d_{t-1}+1=2h+b+2$.  The chain part is never empty.
\end{lemma}

\begin{proof}
By Proposition \ref{prop:PS} a vertex $j\le t-1$ is partial iff
$r_j=u_j+o_j$ is odd and $0<r_j<2o_j$, where $o_j=\max(0,j-B)=\max(0,j-b+1)$.

If $j<b$ then $o_j=0$ and there is no partial slot.  If $b\le j<h+b$ then
$u_j=-\min(j,h)$ and $r_j=(j-b+1)-\min(j,h)$; for $j\ge h$ this is
$j-h-b+1\le0$, and for $j<h$ it is $1-b\le-1$ since $b=n-k\ge2$.  Either way
$r_j\le0$: no partial slot.  If $j\ge h+b$ then $u_j=j-2h-b$ and $o_j=j-b+1$,
so $r_j=2(j-h-b)+1$ is odd and positive, and $r_j<2o_j$ reads $-2h<1$, always
true.  So every $j\in[h+b,t-1]$ is partial, at degree
$d_j=t-u_j=t+2h+b-j$ by \eqref{eq:PS}.

At the node, $q+x=4$ is even, so $(E_t,t)$ is partial and, by \eqref{eq:Dt},
$E_t=t-k+2h+4$.  Since $b=n-k=t+2-k$ we get $E_t=2h+b+2=d_{t-1}+1$.  Finally
$h\le k-q=k-3$ gives $h+b\le k-3+n-k=n-3=t-1$, so the chain part is nonempty.
\end{proof}

\begin{proposition}[\textup{(L1)}]\label{prop:L1}
$\mathcal{P}$, with the slot adjacency of $I_\ka$, is a path, and $(E_t,t)$ is
one of its two endpoints.
\end{proposition}

\begin{proof}
Two slots $(d,j)$ and $(d',i)$ are joined by an arrow iff $i\sim j$ and
$|d-d'|=1$.  By Lemma \ref{lem:partialset} the vertices carrying a partial slot
are the consecutive $h+b,h+b+1,\dots,t-1,t$, and the degrees
$d_{h+b}=t+h,\ d_{h+b+1}=t+h-1,\ \dots,\ d_{t-1}=2h+b+1$ decrease by exactly
one per step, with $E_t=d_{t-1}+1$.  So consecutive vertices give adjacent
slots, and non-consecutive chain vertices $i,j$ have
$|d_i-d_j|=|i-j|\ge2$, hence no arrow; the leaves carry no partial slot by
Proposition \ref{prop:PS}.  Therefore $\mathcal{P}$ is the path
\[
   (d_{h+b},h+b)\ \text{---}\ (d_{h+b+1},h+b+1)\ \text{---}\cdots\text{---}\
   (d_{t-1},t-1)\ \text{---}\ (E_t,t),
\]
whose endpoints are $(t+h,h+b)$ and $(E_t,t)$.
\end{proof}

\begin{proposition}[\textup{(L3)}]\label{prop:L3}
The line$\to$partial obligations of $T^{\,l}(v)$ are exactly the edges of the
path $\mathcal{P}$, each traversed once.
\end{proposition}

\begin{proof}
An obligation is line$\to$partial iff its source is the partial slot of its
vertex $j$ and its target is the partial slot of a neighbour $i$; by
Proposition \ref{prop:PS} each vertex has at most one partial slot, so this
says $d_i=d_j+1$ with $i\sim j$.  With $d_j=t+2h+b-j$ on the chain and
$E_t=d_{t-1}+1$, the possibilities are $i=j-1$ (requiring $j-1\ge h+b$) and
$(j,i)=(t-1,t)$; whereas $i=j+1$ gives $d_i=d_j-1$, the node-to-chain step
$(j,i)=(t,t-1)$ gives $d_{t-1}=E_t-1$, and the leaves carry no partial slot.
Those are exactly the edges of $\mathcal{P}$, each traversed once, in the
direction away from vertex $t-1$.
\end{proof}

\begin{proposition}[\textup{(L4)}]\label{prop:L4}
No line$\to$empty obligation occurs for $T^{\,l}(v)$ in the family
$(1,-1,\pm1)$.
\end{proposition}

\begin{proof}
Such an obligation would be $\bigl((d_j,j),(d_j+1,i)\bigr)$ with $(d_j+1,i)$
nonzero in $I_\ka$ but untaken, i.e.\ $d_j+1<E_i$, where
$E_i$ is the lowest taken degree at $i$.  At a partial vertex
$E_i=d_i$.  We run over the neighbours of a partial vertex $j$, i.e.\
$h+b\le j\le t$.
\begin{itemize}\itemsep1pt
\item $i=j-1$ with $j-1\ge h+b$: $E_i=d_i=d_j+1$, so the target is
taken (partial): a path edge.
\item $i=j+1\le t-1$: $E_i=d_i=d_j-1\le d_j+1$: taken, and full.
\item $(j,i)=(t-1,t)$: $E_t=d_{t-1}+1$: a path edge.
\item $(j,i)=(t,t-1)$: $E_{t-1}=d_{t-1}=E_t-1\le E_t+1$: taken.
\item $j=t$, $i$ a leaf: by Lemma \ref{lem:node} the offsets are
$E_l-E_t=\mp z$ or $\pm z$, i.e.\ $\pm1$, so $E_l\le E_t+1$ and the target
$(E_t+1,l)$ is taken whenever it exists.
\item $j=h+b$, $i=h+b-1\ge1$: the only nontrivial case.  Here $d_j=t+h$ and the
target is $(t+h+1,\,h+b-1)$.  At $i=h+b-1$ we have
$o_i=\max(0,i-B)=\max(0,h)=h$ and, since $i=h+b-1\ge h$ (as $b\ge2$),
$u_i=-\min(i,h)=-h$; so $r_i=u_i+o_i=0$ and $i$ carries no partial slot, and
$v_i=\ms_i-o_i$ is exactly the number of minus-only slots at $i$.  Hence the
greedy stops at the bottom minus-only slot, of degree
\[
   E_i=\bigl(2t-|i-k|\bigr)-2\bigl(\ms_i-h\bigr)+2 .
\]
Using $k=t+2-b$ and $i=h+b-1=h+t+1-k$, this equals $t+h+1$ whether $i\le k$ or
$i>k$: for $i\le k$ it is $2t-k+i-2i+2h+2=2t-k-i+2h+2=t+h+1$, and for $i>k$ it
is $2t-i+k-2k+2h+2=2t-i-k+2h+2=t+h+1$.  So the obligation lands exactly on
$E_i$: the target is taken, and full.
\end{itemize}
Hence no line$\to$empty obligation occurs.
\end{proof}

\subsection{No full\texorpdfstring{$\to$}{->}partial chain obligation in the
four pinned families}\label{ss:L5}

This subsection isolates and proves the statement that class (ii) of Theorem
\ref{thm:B} cannot supply in the four families with $q+x$ even and $z\ne0$.
It is a genuine gap in the class list otherwise: in those families the partial
slots of $T^{\,l}(v)$ carry $\im(a_{l\to t})$ and its transports, which by
Proposition \ref{prop:W4l} (see \textup{(N6)}) is \emph{not} a
$\tau$-eigenvector, so the discharge ``(T1) makes the image the
$\tau$-eigenline, which is what the module carries'' asserts exactly the
opposite of the truth.  A single full$\to$partial chain obligation in those
families would make $T^{\,l}(v)$ fail to be a submodule and Theorem
\ref{thm:B} --- hence Theorem \ref{thm:typeD} --- false.  We prove that the
class is empty.  Nothing below is a census: the argument is the closed-form
threshold arithmetic of \S\ref{sec:families}, run once for $(0,-1,\pm1)$ and
once for $(1,-1,\pm1)$.

Throughout, for a vertex $j$ we write $E_j$ for the lowest taken
degree of the greedy top truncation at $j$ (so $E_j=d_j$ at a
partial vertex, the partial slot being the last one the greedy reaches).

\begin{proposition}[\textup{(L5)}]\label{prop:L5}
Let $\mu=\mu(h;x,y,z)$ belong to one of the four families with $q+x$ even and
$z\ne0$, i.e.\ $(x,y,z)\in\{(0,-1,\pm1),(1,-1,\pm1)\}$.  Then:
\begin{enumerate}
\item[\textup{(a)}] for $(x,y,z)=(0,-1,\pm1)$ the unique partial slot is
$(E_t,t)$, and $E_{t-1}=E_t+1$;
\item[\textup{(b)}] for $(x,y,z)=(1,-1,\pm1)$ every chain obligation whose
target is a partial slot has a partial slot as its source, i.e.\ is an edge of
the path $\mathcal{P}$ of Proposition \ref{prop:L1};
\item[\textup{(c)}] in all four families, \emph{no full$\to$partial obligation
occurs along a chain arrow}.  Equivalently, every obligation of $T^{\,l}(v)$
whose target is a partial slot is either the leaf obligation
$a_{l\to t}:(E_t-1,l)\to(E_t,t)$ of class \textup{(v)} or an edge of
$\mathcal{P}$.
\end{enumerate}
\end{proposition}

\begin{proof}
\emph{(a).}  Here $q=2$, $x=0$, $y=-1$, $z=\pm1$.  Since $q$ is even,
Corollary \ref{cor:parity} gives no partial slot at any chain vertex
$j\le t-1$; and since $q+x=2$ is even, Lemma \ref{lem:node} makes $(E_t,t)$
partial.  So $(E_t,t)$ is the unique partial slot, as recorded in Proposition
\ref{prop:partials}, and by \eqref{eq:Dt}
\[
   E_t=t-k+2h+q+x+[\,q+x\text{ odd}\,]=t-k+2h+2 .
\]

We compute $E_{t-1}$.  By Lemma \ref{lem:vclosed} (and its proof)
$v_{t-1}=2k-2h-q-[k=t]=2k-2h-2-[k=t]$, while $\ms_{t-1}=\min(t-1,k)$ and
$o_{t-1}=k-2$ (Lemma \ref{lem:increments}(iv)), so the number of minus-only
slots at $t-1$ is
\[
   \ms_{t-1}-o_{t-1}=
   \begin{cases}k-(k-2)=2,&k\le t-1,\\ (t-1)-(k-2)=1,&k=t,\end{cases}
   \qquad\text{i.e.}\qquad 2-[k=t].
\]
Greedy from the top exhausts these first, leaving
$r_{t-1}=v_{t-1}-(\ms_{t-1}-o_{t-1})=2k-2h-4$ dimensions, which is even and
non-negative (as $h\le k-q=k-2$) and at most $2o_{t-1}=2k-4$.  So the greedy
fills exactly $k-h-2$ window slots \emph{completely} and leaves none partial
--- confirming Corollary \ref{cor:parity} --- and occupies
$(2-[k=t])+(k-h-2)$ degrees in all.  Since
$\hi^{-}_{t-1}=2t-|t-1-k|=t+k+1-2[k=t]$, the lowest of them is
\[
   E_{t-1}=\hi^{-}_{t-1}-2\bigl((2-[k=t])+(k-h-2)-1\bigr)
   =\bigl(t+k+1-2[k=t]\bigr)-\bigl(2k-2h-2-2[k=t]\bigr)=t-k+2h+3,
\]
the terms in $[k=t]$ cancelling.  Hence $E_{t-1}=E_t+1$, as claimed.

\emph{(b).}  Here $q=3$, $x=1$, $b=n-k$, $B=b-1$, $h\le k-3$.  By Lemma
\ref{lem:partialset} the partial slots are $(d_j,j)$ for $h+b\le j\le t-1$,
with $d_j=t+2h+b-j$, together with $(E_t,t)$ where $E_t=d_{t-1}+1$; and the
leaves carry none (Proposition \ref{prop:PS}).  An obligation with target a
partial slot $(d_i,i)$ has source $(d_i-1,j)$ for some neighbour $j$ of $i$.
We run over all of them.  This is exactly the bullet list in the proof of
Proposition \ref{prop:L4}, read in the opposite direction: there we asked,
for each partial vertex, what lies one degree \emph{above} at each neighbour;
here we ask what lies one degree \emph{below}.  The same four threshold
computations answer both questions.

\begin{itemize}\itemsep1pt
\item \emph{$i\le t-1$, $j=i+1\le t-1$.}  Then $d_{i+1}=d_i-1$, so the source
slot $(d_i-1,i+1)$ \emph{is} the partial slot at $i+1$: a line.  This is a
path edge.
\item \emph{$i=t-1$, $j=t$.}  The source is $(d_{t-1}-1,t)=(E_t-2,t)$, and
$E_t=E_t>E_t-2$, so the source is untaken: there is no obligation.
\item \emph{$i\le t-1$, $j=i-1\ge h+b$.}  Then $E_{i-1}=d_{i-1}
=d_i+1>d_i-1$, so the source is untaken: no obligation.
\item \emph{$i=h+b$, $j=h+b-1$.}  This vertex carries no partial slot, and the
last bullet of the proof of Proposition \ref{prop:L4} computes
$E_{h+b-1}=t+h+1=d_{h+b}+1=d_i+1>d_i-1$: the source is untaken, so
again there is no obligation.  (The vertex exists, since $b=n-k\ge2$ gives
$h+b-1\ge1$.)
\item \emph{$i=t$, $j=t-1$.}  The source is $(E_t-1,t-1)=(d_{t-1},t-1)$, the
partial slot at $t-1$: a line.  This is the remaining path edge.
\item \emph{$i=t$, $j$ a leaf.}  These are the two leaf obligations into
$(E_t,t)$; by Lemma \ref{lem:node} exactly one of them is nonzero, namely the
one with $E_l=E_t-1$, and it is the class \textup{(v)} obligation.  Neither is
a chain arrow.
\end{itemize}
So every chain obligation into a partial slot has a partial source, which is
(b).

\emph{(c).}  In the two families of (a) the only partial slot is $(E_t,t)$, so
any full$\to$partial chain obligation would have source $(E_t-1,t-1)$; but
$E_{t-1}=E_t+1>E_t-1$, so that slot is untaken and no obligation
starts there.  (The other two neighbours of $t$ are leaves, whose obligations
are not chain arrows.)  In the two families of (b), every chain obligation
into a partial slot has a \emph{line} source, never a full one.  In all four
families, therefore, the class is empty.
\end{proof}

\begin{remark}[the hypotheses actually used]\label{rem:L5hyp}
Proposition \ref{prop:L5} draws on exactly four earlier facts, and each holds
in the stated range.  (1) Corollary \ref{cor:parity}, that chain partial slots
occur iff $q$ is odd --- used for $(0,-1,\pm1)$, where $q=2$.  (2) Lemma
\ref{lem:node}, that $(E_t,t)$ is partial iff $q+x$ is even and that the leaf
offsets are $\pm1$ --- both families have $q+x$ even.  (3) Lemma
\ref{lem:vclosed} and Lemma \ref{lem:increments}(iv), for $v_{t-1}$ and
$o_{t-1}=k-2$; the latter needs $k\ge2$, which is Convention \ref{conv:D}.
(4) Lemma \ref{lem:partialset} and the last bullet of Proposition
\ref{prop:L4}, for $(1,-1,\pm1)$; these need $h\le k-3$ and $b=n-k\ge2$, both
guaranteed by Proposition \ref{prop:reps} with $q=3$ and $k\le n-2$.  In
particular the proof never assumes that the chain interval of partial slots is
nonempty in the two families of (a) --- it is empty there, by Corollary
\ref{cor:parity} --- nor that it is nonempty in (b), where Lemma
\ref{lem:partialset} shows that it is.
\end{remark}

\begin{remark}[what a census would and would not give]\label{rem:L5census}
An exhaustive census over $D_4$--$D_{44}$, all $2\le k\le n-2$, all nine
families and all $h$ --- some $4.8\times10^{7}$ obligations --- confirms
Proposition \ref{prop:L5}: the class ``full$\to$partial along a chain arrow''
is populated only in the family $(1,0,0)$, never in the four pinned families,
and in the pinned families every chain obligation into a partial slot has a
line source (\S\ref{sec:verification}).  That census is a check, not the
proof; the proof is the threshold arithmetic above, which is uniform in $n$,
$k$ and $h$.
\end{remark}


\section{The induction along \texorpdfstring{$\sigma$}{sigma}}
\label{sec:induction}

We now prove Proposition \ref{prop:T}.  Fix $\mu=w\la$ and a reduced word for
$w$, and let $0=X_0\subseteq X_1\subseteq\cdots\subseteq X_\ell=V$ be the
filtration of Theorem \ref{thm:main}, indexed by blocks.  Each $X_p$ is
extremal (Corollary \ref{cor:prefix}).  We carry along the hypothesis
\[
   (H)\qquad
   X \text{ is a top truncation at } t \text{ and an up-set at each of the two
   leaves,}
\]
which holds for $X_0=0$.  Steps at colours $j\notin\{t,n-1,n\}$ change $X$ only
at $j$, so $(H)$ is untouched by them.  We write
$v'=\dimv X$, $\mu'=\la-\sum_iv'_i\al_i=\sum_i\eta_ie_i$ with
$\eta_i\in\{0,\pm1\}$, and $D=D_t(X)$ for the threshold of Definition
\ref{def:upset}.

\subsection{Weight arithmetic at the node}

\begin{lemma}\label{lem:B0}
For any graded submodule $X\subseteq I_\ka$ with $\dimv X$ extremal,
\[
   v'_{n-1}+v'_n=v'_t-\eta_{n-1},
   \qquad
   v'_{n-1}-v'_n=\eta_n,
\]
and, writing $c:=\langle\al_t^\vee,\mu'\rangle$,
\[
   c=[k=t]-2v'_t+v'_{t-1}+v'_{n-1}+v'_n .
\]
If $c\ge1$ then $\eta_{n-1}\le 0$, and
\begin{equation}\label{eq:B1}
   v'_{t-1}\ \ge\ v'_t+1+\eta_{n-1}-[k=t]\ \ge\ v'_t-[k=t].
\end{equation}
\end{lemma}

\begin{proof}
The coefficient of $e_{n-1}$ in $\sum_jv'_j\al_j$ is $-v'_t+v'_{n-1}+v'_n$ and
that of $e_n$ is $-v'_{n-1}+v'_n$; since $\la$ has support in the first $n-2$
coordinates, these equal $-\eta_{n-1}$ and $-\eta_n$.  The formula for $c$ is
$\langle\al_t^\vee,\la\rangle-\sum_ja_{tj}v'_j$.  Next,
$c=\langle\al_t^\vee,\mu'\rangle=\eta_{n-2}-\eta_{n-1}$ with
$\eta_i\in\{0,\pm1\}$, so $c\ge1$ forces $\eta_{n-1}\le0$.  Substituting
$v'_{n-1}+v'_n=v'_t-\eta_{n-1}$ into $c\ge1$ gives \eqref{eq:B1}.
\end{proof}

\begin{lemma}\label{lem:B2}
Suppose $X$ is an up-set at $t-1$ and put
$\mathrm{cum}:=\sum_{d\ge D+1}\dim(I_\ka)_{t-1}[d]$.  If
$v'_{t-1}\ge\mathrm{cum}$ then $X_{t-1}[D+1]$ is full; and if
$v'_{t-1}\ge\mathrm{cum}-1$ then $X_{t-1}[D+1]$ is deficient by at most one
dimension.
\end{lemma}

\begin{proof}
Let $D'$ be the threshold of $X$ at $t-1$.  If $D'\ge D+1$ then $X_{t-1}$ is
supported in degrees $\ge D'$ and is proper at $D'$, so
$v'_{t-1}\le\mathrm{cum}-1$; contrapositively $v'_{t-1}\ge\mathrm{cum}$ forces
$D'<D+1$, i.e.\ $X_{t-1}[D+1]$ full.  If $D'=D+1$ the deficiency at $D+1$ is
$\dim(I_\ka)_{t-1}[D+1]-\dim X_{t-1}[D+1]$ and
$v'_{t-1}\le\mathrm{cum}-(\text{that deficiency})$, so
$v'_{t-1}\ge\mathrm{cum}-1$ bounds it by $1$; if $D'>D+1$ then
$v'_{t-1}\le\mathrm{cum}-1-\dim(I_\ka)_{t-1}[D+1]$, which combined with
$v'_{t-1}\ge\mathrm{cum}-1$ forces $(I_\ka)_{t-1}[D+1]=0$ and the slot is full
vacuously.
\end{proof}

\begin{lemma}[the profile bound at $t-1$]\label{lem:cumbound}
Assume $(H)$, and let $h\ge0$.
\begin{enumerate}
\item[\textup{(i)}] If $v'_t=2h$, so that $D=t+k-2h$, then
$\mathrm{cum}\le\max(0,h-1)+\max(0,h+1-[k=t])$, which equals $2h-[k=t]$ when
$h\ge1$ and $1-[k=t]$ when $h=0$.
\item[\textup{(ii)}] If $v'_t=2h+1$, so that $D=t+k-2h-2$, then
$\mathrm{cum}\le 2h+2-[k=t]$.
\end{enumerate}
\end{lemma}

\begin{proof}
By \eqref{eq:foldedP} the multiplicity profile at $t-1$ is the sum of two
indicator intervals: the \emph{low} one on $[\,|t-1-k|,\ t+k-3\,]$ and the
\emph{high} one on $[\,t-k+3,\ 2t-|t-1-k|\,]$, each with $\min(t-1,k)$ slots,
in steps of $2$.  Degrees at $t-1$ have the parity of $t-1+k$.  (When
$\operatorname{char}\ne2$ these two intervals are the supports of the $+$ and
$-$ strands of Proposition \ref{prop:W2}, but the counting below uses only the
profile, so this lemma is valid over every field.)

(i) $D+1=t+k-2h+1\equiv t-1+k$.  The low interval contributes the slots
$d\le t+k-3$ with $d\ge D+1$, of which there are
$\lfloor((t+k-3)-(t+k-2h+1))/2\rfloor+1=h-1$ when $h\ge1$ and none otherwise,
capped at $\min(t-1,k)$: that is $\max(0,h-1)$.  The high interval tops out at
$2t-|t-1-k|$, which is $t+k+1$ if $k\le t-1$ and $2t-1$ if $k=t$; the number of
its slots $\ge D+1$ is $h+1$, respectively $h$, when nonnegative: that is
$\max(0,h+1-[k=t])$, again capped.  Summing gives the bound, and for $h\ge1$
the two terms are $h-1$ and $h+1-[k=t]$.

(ii) $D+1=t+k-2h-1$.  The same count gives $h$ from the low interval and
$h+2-[k=t]$ from the high interval, whence $\mathrm{cum}\le 2h+2-[k=t]$.
\end{proof}

\subsection{Leaf steps}\label{ss:leafsteps}

\begin{lemma}\label{lem:leafstep}
Assume $(H)$ and let the step have colour $j\in\{n-1,n\}$.  Then $(H)$ is
preserved.
\end{lemma}

\begin{proof}
The leaf column is multiplicity-free with degrees $t+k-1>t+k-3>\cdots>t-k+1$
(Lemma \ref{lem:P}), and by $(H)$ the module occupies the top $v'_j$ of them,
so the largest unfilled leaf degree is $D_j=t+k-1-2v'_j$.  We have
$c_j=\langle\al_j^\vee,\mu'\rangle=v'_t-2v'_j\ge1$ because the word is reduced
(Lemma \ref{lem:extend}(i)).  If $v'_t=2h$ this forces $v'_j\le h-1$ and
$D_j\ge t+k-2h+1>D$; if $v'_t=2h+1$ it forces $v'_j\le h$ and
$D_j\ge t+k-2h-1>D$.  Either way $D_j>D$.

The socle step at $j$ adds exactly the unfilled leaf slots $(e,j)$ for which
$a_{j\to t}\bigl((I_\ka)_j[e]\bigr)\subseteq X_t[e+1]$.  If $e\ge D$ then
$e+1>D$ and $X_t[e+1]$ is full by $(H)$, so the condition holds.  If
$e<D-1$ then $X_t[e+1]=0$ by $(H)$, while $a_{j\to t}$ is injective on the leaf
slot by Lemma \ref{lem:R}, so the condition fails.  The degree $e=D-1$ is not
decided by this dichotomy: there $X_t[e+1]=X_t[D]$ is \emph{proper} by
Definition \ref{def:upset}, so the inclusion
$a_{j\to t}\bigl((I_\ka)_j[D-1]\bigr)\subseteq X_t[D]$ may hold or fail, and
$D-1$ need not be added.  Hence, the unfilled degrees being exactly those
$\le D_j$, the added set \emph{contains} $[D,\,D_j]$ and \emph{is contained
in} $[\max(D-1,\ \text{bottom}),\,D_j]$.  Both of these are top segments of
the unfilled block; and since the only leaf degree in which they differ is
$D-1$, the added set is one of the two, hence itself a top segment of the
unfilled block.
Adding a top segment of the unfilled block to a final segment of degrees leaves
a final segment.  So $X$ is still an up-set at $j$, and it is unchanged at $t$
and at the other leaf.
\end{proof}

\subsection{Colour-\texorpdfstring{$t$}{t} steps: Claim G and Claim F}
\label{ss:G}

\begin{proposition}[Claim G]\label{prop:G}
Assume $(H)$, that $\dimv X$ is extremal, that
$c=\langle\al_t^\vee,\mu'\rangle\ge1$, and $(*^\circ)$ for $X$ \textup{(}whose
hypotheses $X$ meets, being extremal with $c\ge1$; $(*^\circ)$ itself is
available from Corollary \ref{cor:reduce} --- which is characteristic-free ---
together with $(B)$, and $(B)$ comes from Theorem \ref{thm:B} when
$\operatorname{char}\ne2$ and from Step 4 of Theorem \ref{thm:charp} at
$p=2$, where Theorem \ref{thm:B} is not claimed\textup{)}.  Then:
\begin{itemize}\itemsep1pt
\item[\textup{(G1)}] if $X_t[D]\ne0$ then every neighbour of $t$ is full at
degree $D+1$;
\item[\textup{(G2)}] if $X_t[D]=0$ then at most one neighbour of $t$ is
deficient at degree $D+1$, and that deficiency is one dimensional; moreover if
some neighbour is deficient then $\dim(I_\ka)_t[D]=2$.
\end{itemize}
\end{proposition}

\begin{proof}
Under $(H)$ the column at $t$ is a top truncation, and $X$ is extremal with
$c\ge1$, so Lemma \ref{lem:nodefree}(b) gives $v'_t<2k$ and Lemma
\ref{lem:nodefree}(a) then gives
\[
   X_t[D]\ne0
   \iff v'_t \text{ is even and } \ge2 ,
\]
and $D=t+k-2h$ when $v'_t=2h$ with $h\ge1$, while $D=t+k-2h-2$ when
$v'_t=2h+1$, and $D=t+k$ when $v'_t=0$.  (Lemma \ref{lem:node} is \emph{not}
what is being used here: it is stated only for the nine representatives and
only for $\operatorname{char}\ne2$, and its $E_t$ is the lowest taken degree
rather than the threshold of Definition \ref{def:upset}; see Remark
\ref{rem:nodefree}.)

\smallskip\noindent\emph{(G1): $v'_t=2h$, $h\ge1$.}
\emph{The leaves.}  By Lemma \ref{lem:B0}, $\eta_{n-1}\le0$ and
$v'_{n-1}+v'_n=2h-\eta_{n-1}\in\{2h,2h+1\}$ with difference $\eta_n$.  If the
sum is $2h$ its parity forces $\eta_n=0$ and both leaves equal $h$; if the sum
is $2h+1$ then $\eta_n=\pm1$ and the leaves are $\{h,h+1\}$.  Either way
$\min(v'_{n-1},v'_n)\ge h$.  Now $D+1=t+k-2h+1$ is the $h$-th leaf degree from
the top, and an up-set at a leaf is full there exactly when $v'_l\ge h$.  So
both leaves are full at $D+1$.
\emph{The chain.}  By \eqref{eq:B1}, $v'_{t-1}\ge2h-[k=t]$, and by Lemma
\ref{lem:cumbound}(i) with $h\ge1$, $\mathrm{cum}\le 2h-[k=t]$.  By
$(*^\circ)$, $X$ is an up-set at $t-1$, so Lemma \ref{lem:B2} applies and
$X_{t-1}[D+1]$ is full.

\smallskip\noindent\emph{(G2), first branch: $v'_t=2h+1$.}
Here $D+1=t+k-2h-1$ is the $(h+1)$-st leaf degree from the top, so a leaf is
full at $D+1$ iff $v'_l\ge h+1$.  By Lemma \ref{lem:B0},
$v'_{n-1}+v'_n=2h+1-\eta_{n-1}$ with difference $\eta_n$, and $\eta_{n-1}\le0$.
\begin{itemize}\itemsep1pt
\item If $\eta_{n-1}=0$ the sum is $2h+1$, odd, so $\eta_n=\pm1$ and the leaves
are $\{h,h+1\}$: exactly one leaf is deficient, by exactly one dimension.  But
then \eqref{eq:B1} reads $v'_{t-1}\ge 2h+2-[k=t]\ge\mathrm{cum}$ by Lemma
\ref{lem:cumbound}(ii), so the chain neighbour is \emph{full}.
\item If $\eta_{n-1}=-1$ the sum is $2h+2$, even, so $\eta_n=0$ and both leaves
equal $h+1$: both leaves are full.  Only the chain neighbour can be deficient,
and \eqref{eq:B1} gives $v'_{t-1}\ge 2h+1-[k=t]\ge\mathrm{cum}-1$, so by Lemma
\ref{lem:B2} it is deficient by at most one dimension.
\end{itemize}
Either way at most one neighbour is deficient, by at most one dimension.

For the last clause, suppose a neighbour is deficient.  We have
$\dim(I_\ka)_t[D]=1$ only for $D\in\{t-k,t+k\}$, and $D=t+k$ is excluded since
$v'_t\ge1$.  So suppose $D=t-k$, i.e.\ $v'_t=2k-1$.  Then
$v'_t=r'+2m'$ (Lemma \ref{lem:C} applied to $\mu'$) is odd, so $r'=1$ and
$m'=k-1$; thus all $k-1$ chain entries of $\operatorname{supp}\mu'$ carry sign
$-1$, whence $\eta_{n-2}\le0$, and $c=\eta_{n-2}-\eta_{n-1}\ge1$ forces
$\eta_{n-2}=0$ and $\eta_{n-1}=-1$, hence $\eta_n=0$ and $r'=1$ is accounted
for by $n-1$.  So both leaves have $v'_l=k$ and are full (the leaf column has
exactly $k$ slots).  For the chain, $\mathrm{cum}$ is the whole column at $t-1$
above degree $D+1=t-k+1$; the column starts at $|t-1-k|$, which is $t-1-k$ when
$k\le t-1$ and $1$ when $k=t$, so $\mathrm{cum}=2k-1$ when $k\le t-1$ (the
single bottom $+$ slot is excluded) and $\mathrm{cum}=2k-2$ when $k=t$.  Now
\eqref{eq:B1} gives $v'_{t-1}\ge 2k-1-[k=t]$, which is $\ge\mathrm{cum}$ in
both cases, so by Lemma \ref{lem:B2} the chain neighbour is full too.  No
neighbour is deficient after all --- contradiction.  Hence
$\dim(I_\ka)_t[D]=2$.

\smallskip\noindent\emph{(G2), second branch: $v'_t=0$.}
Here $X_t=0$ and $D=t+k$, the multiplicity-one top slot.  We show that no
neighbour is deficient at all.
\begin{itemize}\itemsep1pt
\item \emph{The leaves are empty.}  If $0\ne u\in X_{n-1}[e]$, its only arrow
lands in $X_t[e+1]=0$, so $u$ spans a copy of the leaf simple inside
$\soc(I_\ka)=S_\ka$; but $\ka=k\le n-2$ is a chain vertex.  So $X_{n-1}=0$ and
likewise $X_n=0$: $v'_{n-1}=v'_n=0$.
\item \emph{The chain neighbour is not empty.}  By Lemma \ref{lem:B0},
$c=[k=t]+v'_{t-1}\ge1$, so $v'_{t-1}\ge1-[k=t]$.
\item \emph{$\mathrm{cum}\le1$.}  Here $D+1=t+k+1$.  At $t-1$ the low interval
tops out at $t+k-3$ and the high interval at $2t-|t-1-k|$, which is $t+k+1$ if
$k\le t-1$ and $t+k-1$ if $k=t$.  So $\mathrm{cum}=1-[k=t]$ exactly.
\item \emph{No deficiency.}  By Lemma \ref{lem:P} the leaf top degree is
$t+k-1<t+k+1=D+1$, so $(I_\ka)_l[D+1]=0$ for both leaves and they are full
vacuously.  For the chain: if $k=t$ then $\mathrm{cum}=0$ and
$(I_\ka)_{t-1}[D+1]=0$, again full vacuously; if $k\le t-1$ then
$v'_{t-1}\ge1=\mathrm{cum}$ and Lemma \ref{lem:B2} gives fullness.
\end{itemize}
So (G2) holds vacuously in this branch, and the last clause is vacuous too.
\end{proof}

\begin{corollary}[Claim F]\label{cor:F}
Assume $(H)$, $c\ge1$, \emph{and that $\dimv X$ is extremal}.  If $X_t[D]$ is
a proper subspace of $(I_\ka)_t[D]$ then $\soc_{(t)}(I_\ka/X)$ has a nonzero
component in degree $D$.
\end{corollary}

\begin{remark}\label{rem:Fextremal}
The extremality hypothesis is not decorative and cannot be removed: by
\textup{(N3)} of \S\ref{ss:negative}, Claim F is \emph{false} for arbitrary
graded submodules satisfying $(H)$ and $c\ge1$.  It enters through
Proposition \ref{prop:G}, whose proof invokes $(*^\circ)$ for $X$, and
$(*^\circ)$ is a statement about \emph{extremal} modules only (Definition
\ref{def:star}).  At the one place where Corollary \ref{cor:F} is used --- the
proof of Proposition \ref{prop:T} --- the module $X=X_p$ is a prefix of the
filtration of Theorem \ref{thm:main} and is therefore extremal by Corollary
\ref{cor:prefix}, so the hypothesis holds.
\end{remark}

\begin{proof}
An element $z\in(I_\ka)_t[D]$ lies in $\soc_{(t)}(I_\ka/X)$ iff all three
images $a_{t\to i}z$ ($i\in\{t-1,n-1,n\}$) lie in $X$.  If no neighbour is
deficient at $D+1$ --- which is (G1), and also the case $v'_t=0$ of (G2) ---
then all three conditions are automatic and the whole of
$(I_\ka)_t[D]/X_t[D]\ne0$ lies in the socle.  Otherwise, by (G2), exactly one
neighbour is deficient, by exactly one dimension, and $\dim(I_\ka)_t[D]=2$.  So
membership in the socle is cut out by one linear condition on a
two-dimensional space, and a nonzero solution exists; it is not contained in
$X_t[D]$ for dimension reasons unless $X_t[D]$ is the whole solution space,
which is impossible because $X_t[D]=0$ in the (G2) branch.
\end{proof}

\subsection{Preservation of \texorpdfstring{$(H)$}{(H)} at colour-$t$ steps}

\begin{lemma}\label{lem:ab}
Assume $(H)$.  Then
\begin{enumerate}
\item[\textup{(a)}] $\soc_{(t)}(I_\ka/X)$ vanishes in every degree $<D-2$;
\item[\textup{(b)}] if $\soc_{(t)}(I_\ka/X)$ is nonzero in degree $D-2$ then
$\dim X_t[D]=1$ and $\dim(I_\ka)_t[D]=2$.
\end{enumerate}
\end{lemma}

\begin{proof}
For (a), let $0\ne x\in(I_\ka)_t[d]$ represent an element of
$\soc_{(t)}(I_\ka/X)$ with $d<D-2$.  Then $X_t[d+2]=0$ by $(H)$, so every
double path $a_{j\to t}a_{t\to j}x$ vanishes.  Not all $a_{t\to j}x$ can
vanish: otherwise $x$ spans a copy of $S_t$ inside $\soc(I_\ka)=S_\ka$, forcing
$d=2t$, whereas $d<D\le t+k\le 2t$.  So some $y=a_{t\to j}x\ne0$ satisfies
$a_{j\to t}y=0$.  If $j$ is a leaf, $y$ is killed by its only arrow and hence
lies in $\soc(I_\ka)=S_\ka$, forcing $\ka$ to be a leaf --- excluded, as
$\ka=k\le n-2$.  So $j=t-1$.  Before propagating we record the degree bound.  By $(H)$ and
Definition \ref{def:upset}, $D\le t+k$ (the top colour-$t$ degree, Lemma
\ref{lem:P}), so $d<D-2\le t+k-2$; the element $y=a_{t\to(t-1)}x$ sits at
degree $d+1\le t+k-2$, at the vertex $t-1$.  Since $k\le t$,
\begin{equation}\label{eq:abdeg}
   d+1+p\ \le\ t+k-2+p\ <\ 2t \qquad\text{for every } 0\le p\le t-k-1 ,
\end{equation}
and more crudely $d+1\le t+k-2\le 2t-2<2t$.  (Note that the relevant bound is
this one and not the leaf top degree $t+k-1$: $y$ sits at the chain vertex
$t-1$, and the bound that applies comes from $d<D-2$ and $D\le t+k$.)

Now propagate outward exactly as in Lemma \ref{lem:Id}: the relation at $t-1$
applied to $y$ has its $t$-term zero, so
$a_{(t-2)\to(t-1)}\bigl(a_{(t-1)\to(t-2)}y\bigr)=0$, and
$a_{(t-1)\to(t-2)}y\ne0$ because otherwise $y$ would be killed by both arrows
out of $t-1$ and hence lie in $\soc(I_\ka)$, which by Lemma \ref{lem:soctop}
forces the vertex to be $\ka=k$ and the degree to be $2t$; the degree is
$d+1<2t$.  At the general step we sit at $(d+1+p,\,t-1-p)$, and the same
alternative arises: the element is killed by both its arrows only if
$t-1-p=k$ and $d+1+p=2t$.  The first gives $p=t-1-k$, and then
\eqref{eq:abdeg} gives $d+1+p\le 2t-3<2t$.  So the iteration never stops
early, reaches vertex $1$, whose only neighbour is $2$, and lands in
$\soc(I_\ka)=S_\ka$, forcing $\ka=1$ --- excluded, as $k\ge2$.

For (b), let $x$ be nonzero in $\soc_{(t)}(I_\ka/X)$ at degree $D-2$.  If
$X_t[D]$ were $0$, the double paths of $x$ would all land in $X_t[D]=0$ and the
argument just given applies verbatim with $d=D-2$, a contradiction.  So
$X_t[D]\ne0$; it is proper by the definition of $D$, hence a line in a
two-dimensional space.
\end{proof}

\begin{proof}[Proof of Proposition \ref{prop:T}]
We show $(H)$ is preserved by every socle step.  Steps at colours
$\notin\{t,n-1,n\}$ do not change $X$ at $t$ or the leaves; leaf steps are
Lemma \ref{lem:leafstep}.  Let the step have colour $t$, write $Y$ for the
result, and put $s:=\dim\bigl(\soc_{(t)}(I_\ka/X)\bigr)[D-2]\ge0$.

\emph{Case $0$ (no-op).}  If $\soc_{(t)}(I_\ka/X)=0$ then $Y=X$ and there is
nothing to prove.  This occurs exactly at the repetitions inside a block of
$\sigma$: by Remark \ref{rem:repeats}, a colour-$t$ socle step is idempotent,
so only the first application in a block can contribute.

\emph{Effective steps.}  Assume $Y\ne X$.  Then $c=\langle\al_t^\vee,\mu'\rangle
\ge1$ by Lemma \ref{lem:extend}(i); moreover $\dimv X=\dimv X_p$ is extremal
by Corollary \ref{cor:prefix}, so Corollary \ref{cor:F} applies.  We
record the new profile degree by degree:
\begin{itemize}\itemsep1pt
\item degrees $<D-2$ are untouched by Lemma \ref{lem:ab}(a) and were $0$ by
$(H)$, so stay $0$;
\item degrees $>D$ were full and remain full;
\item at degree $D-2$: $X_t[D-2]=0$ by $(H)$, so $\dim Y_t[D-2]=s$;
\item at degree $D$: $X_t[D]$ is proper by the definition of $D$, so by
Corollary \ref{cor:F} the step contributes something there, and $Y_t[D]\ne0$.
\end{itemize}
Now split on $s$.
\begin{itemize}\itemsep1pt
\item \emph{$s>0$.}  The step reaches $D-2$, so Lemma \ref{lem:ab}(b) gives
$\dim X_t[D]=1$ and $\dim(I_\ka)_t[D]=2$; hence $(I_\ka)_t[D]/X_t[D]$ is a line
and $Y_t[D]=(I_\ka)_t[D]$ is full.  The new profile is $0$ below $D-2$, of
dimension $s$ at $D-2$, and full from $D$ up: a top truncation with threshold
$D-2$, or, if $s$ is maximal there, with the next colour-$t$ degree below.
\item \emph{$s=0$.}  The new profile is $0$ at every colour-$t$ degree $<D$,
nonzero at $D$, and full above $D$.  Whether $Y_t[D]$ is proper or full, this
is again ``zero below one degree, proper or full at it, full above'', i.e.\ a
top truncation, with threshold $D$ or the next colour-$t$ degree below.
\end{itemize}
In every case $Y$ is a top truncation at $t$, and it is unchanged at the two
leaves.  So $(H)$ is preserved, and by induction $V=X_\ell$ satisfies $(H)$.  In
particular the colour-$t$ multiplicities of $V$ are ``zero below one degree,
proper at it, full above'', i.e.\ the greedy top truncation of the sequence of
Lemma \ref{lem:P} to total $v_t$.
\end{proof}

\begin{remark}
The proof of Proposition \ref{prop:T} uses $(*^\circ)$, whose proof
(\S\S\ref{sec:reps}--\ref{sec:closure}) uses only rigidity and the strand
apparatus, never Proposition \ref{prop:T}.  There is no circularity.
\end{remark}

\begin{remark}
Both $(*^\circ)$ and Claim F genuinely use extremality.  The corresponding
statement for arbitrary graded submodules of $I_\ka$ satisfying $(H)$ is false;
see \S\ref{sec:verification}.
\end{remark}

\section{All characteristics: specialisation from the generic fibre}
\label{sec:charp}

The proof just completed assumed $\operatorname{char}\mathbb{K}\ne2$ at exactly
one place, the strand decomposition of \S\ref{sec:strands}, and used it for
exactly one purpose, the construction in Theorem \ref{thm:B} of a witness
submodule for $(B)$.  We now remove the hypothesis.  Nothing in this section
repairs the strand apparatus at $p=2$; it makes it unnecessary there.

\begin{theorem}\label{thm:charp}
Let $\Gamma=D_n$, $\la=\om_k$ with $2\le k\le n-2$, and let $\mathbb{K}$ be an
arbitrary field.  Then $(B)$, $(*^\circ)$, Proposition \ref{prop:T} and Theorem
\ref{thm:typeD} all hold over $\mathbb{K}$.  More precisely, for every
$\mu\in W\la$ the graded dimensions $\dim_{\mathbb{K}}V(\mu)_i[d]$ are
independent of $\mathbb{K}$.
\end{theorem}

\begin{proof}
Let $p$ be a prime and let $R=\mathbb{Z}_{(p)}$, a discrete valuation ring with
fraction field $\mathbb{Q}$ and residue field $\mathbb{F}_p$.  It suffices to
treat $\mathbb{K}=\mathbb{F}_p$ and $\mathbb{K}=\mathbb{Q}$: a field of
characteristic $0$ or $p$ is a flat --- indeed free --- extension of its prime
field, and all the dimensions below are computed slotwise, so they do not
change under field extension.

\smallskip\noindent
\emph{Step 1: a $\mathbb{Z}$-form, free in every graded slot.}  The sign
convention $\varsigma(i,j)=\operatorname{sign}(j-i)$ of \S\ref{sec:setup} has
coefficients $\pm1$, so $\Pia$ is defined by a presentation over $\mathbb{Z}$:
writing $\Pia_{\mathbb{Z}}$ for the quotient of $\mathbb{Z}\overline{Q}$ by
those relations, we have $\Pia_{\mathbb{Z}}\otimes_{\mathbb{Z}}\mathbb{K}
=\Pia_{\mathbb{K}}$ for every field $\mathbb{K}$, both sides being the same
presented quotient and $-\otimes\mathbb{K}$ being right exact.  Put
$\Lambda:=\Pia_{\mathbb{Z}}e_k$, so that
$\Lambda\otimes\mathbb{K}=\Pia_{\mathbb{K}}e_k=I_\ka(\mathbb{K})$, compatibly
with the path-length grading, which is defined over $\mathbb{Z}$.

\emph{$\Lambda$ is a finitely generated $\mathbb{Z}$-module.}  This needs an
argument, because $\mathbb{Z}\overline{Q}$ is not.  The path-length grading is
defined over $\mathbb{Z}$, so $\Lambda=\bigoplus_{d\ge0}\Lambda[d]$, and
$\Lambda[d]$ is spanned over $\mathbb{Z}$ by the images of the paths of length
$d$ starting at $k$ --- a finite set --- so each $\Lambda[d]$ is finitely
generated.  It remains to see that $\Lambda[d]=0$ for $d>2t$.  Now
$\Lambda[d]\otimes\mathbb{K}=(I_\ka(\mathbb{K}))[d]$, which vanishes for
$d>2t$ over every field $\mathbb{K}$, by Lemma \ref{lem:P} (proved
characteristic-free).  A finitely generated abelian group $A$ with
$A\otimes\mathbb{Q}=0$ is finite, and a nonzero finite abelian group has
$A\otimes\mathbb{F}_\ell\ne0$ for some prime $\ell$; so $\Lambda[d]=0$ for
$d>2t$, and $\Lambda=\bigoplus_{d\le2t}\Lambda[d]$ is finitely generated.

\emph{$\Lambda$ is free.}  Write $\Lambda\cong\mathbb{Z}^N\oplus T$ with $T$
finite.  By Corollary \ref{cor:dimI} --- which is proved over an arbitrary
field, from Lemma \ref{lem:R}, Lemma \ref{lem:soctop} and Proposition
\ref{prop:ext1}, and so in particular holds in characteristic $2$ --- we have
$\dim_{\mathbb{K}}\Lambda\otimes\mathbb{K}=k(2n-1-k)$ for \emph{every}
$\mathbb{K}$.  Taking $\mathbb{K}=\mathbb{Q}$ gives $N=k(2n-1-k)$; taking
$\mathbb{K}=\mathbb{F}_\ell$, \emph{for each prime $\ell$ in turn}, gives
$N+\#\{\text{cyclic summands of }T\text{ of }\ell\text{-power order}\}=N$, so
$T$ has no $\ell$-torsion, for every $\ell$.  Hence $T=0$ and $\Lambda$ is
free.  (Only the absence of $p$-torsion is needed for the fixed $p$ at hand,
and that is what the single prime $p$ would give; we range over all primes
because it costs nothing and makes the conclusion the clean one.)  Each graded
slot $\Lambda_i[d]$ is a direct summand of $\Lambda$ as a $\mathbb{Z}$-module,
hence projective, hence free; and its rank is
$\dim_{\mathbb{K}}(I_\ka(\mathbb{K}))_i[d]$ for every $\mathbb{K}$.
(In particular this re-proves Lemma \ref{lem:P} in every characteristic from
the characteristic-$0$ case, although Lemma \ref{lem:P} was in any case proved
characteristic-free.)

\smallskip\noindent
\emph{Step 2: the witness in the generic fibre.}  Let $\mu$ be one of the
representatives of Proposition \ref{prop:reps} and $v=\la-\mu$.  Since
$\operatorname{char}\mathbb{Q}=0\ne2$, Theorem \ref{thm:B} applies over
$\mathbb{Q}$ and produces a graded submodule
$V_{\mathbb{Q}}\subseteq\Lambda_{\mathbb{Q}}:=\Lambda\otimes\mathbb{Q}$ with
$\dimv V_{\mathbb{Q}}=v$ whose graded profile is the greedy top truncation at
every vertex.  (It is $T(v)$ or $T^{\,l}(v)$, built explicitly inside
$I_\ka(\mathbb{Q})$; no descent from an algebraic closure is needed, and by
Corollary \ref{cor:pol} it is the unique submodule of its dimension vector.)

\smallskip\noindent
\emph{Step 3: saturate over $R$.}  Put $\Lambda_R:=\Lambda\otimes R$ and
\[
   N(v)\;:=\;V_{\mathbb{Q}}\cap\Lambda_R\ \subseteq\ \Lambda_{\mathbb{Q}} .
\]
This is a $\Pia_R$-submodule of $\Lambda_R$, since $\Pia_R$ preserves
$V_{\mathbb{Q}}$ and $\Lambda_R$ separately.  It is graded, because both
$V_{\mathbb{Q}}$ and $\Lambda_R$ are: $\Lambda_R$ carries the path-length
grading, which is defined over $\mathbb{Z}$, and $V_{\mathbb{Q}}$ is graded
\emph{by construction} --- in Step 2 it is $T(v)$ or $T^{\,l}(v)$, which
Definition \ref{def:Tv} defines slotwise as a graded subspace of
$I_\ka(\mathbb{Q})$, and in the run of Step 5 below for an arbitrary $\mu$ it
is $\soc_\sigma(I_\ka(\mathbb{Q}))$, which is graded because a socle step is
degree-homogeneous (see the proof of Theorem \ref{thm:spin}).  No appeal to
Proposition \ref{prop:grading} --- a statement over $\mathbb{C}$ --- is made
or needed.  So $N(v)=\bigoplus_{i,d}N(v)_i[d]$ with
$N(v)_i[d]=V_{\mathbb{Q},i}[d]\cap\Lambda_{R,i}[d]$.

Each $N(v)_i[d]$ is \emph{saturated} in the free $R$-module
$\Lambda_{R,i}[d]$: if $x\in\Lambda_{R,i}[d]$ and $px\in N(v)_i[d]$ then
$x\in V_{\mathbb{Q},i}[d]$, that being a $\mathbb{Q}$-subspace, so
$x\in N(v)_i[d]$.  Over a principal ideal domain a saturated submodule of a
finitely generated free module is a direct summand.  Hence $N(v)_i[d]$ is a
free direct summand of $\Lambda_{R,i}[d]$ of rank
$\dim_{\mathbb{Q}}V_{\mathbb{Q},i}[d]$.

\smallskip\noindent
\emph{Step 4: reduce.}  Being a direct summand in each slot, $N(v)$ stays a
direct summand after $-\otimes_R\mathbb{F}_p$; in particular
\[
   N(v)\otimes_R\mathbb{F}_p\ \longrightarrow\
   \Lambda_R\otimes_R\mathbb{F}_p=I_\ka(\mathbb{F}_p)
\]
is injective, its image is a graded $\Pia_{\mathbb{F}_p}$-submodule, and
\begin{equation}\label{eq:specdims}
   \dim_{\mathbb{F}_p}\bigl(\text{image}\bigr)_i[d]
   =\dim_{\mathbb{Q}}V_{\mathbb{Q},i}[d]\qquad\text{for all } i,d .
\end{equation}
In particular the image has dimension vector $v$ and its graded profile is the
greedy top truncation at every vertex.  This is exactly the assertion $(B)$
over $\mathbb{F}_p$, in the form \eqref{eq:Brig}; by Corollary \ref{cor:pol},
which holds over every field, the image is the unique submodule of
$I_\ka(\mathbb{F}_p)$ of dimension vector $v$, i.e.\ it \emph{is}
$V_{\mathbb{F}_p}(\mu)$.

\smallskip\noindent
\emph{Step 5: conclude.}  With $(B)$ available over $\mathbb{F}_p$,
Corollary \ref{cor:reduce} gives $(*^\circ)$ over $\mathbb{F}_p$, and the
induction of \S\ref{sec:induction} gives Proposition \ref{prop:T} over
$\mathbb{F}_p$.  We list exactly what that induction consumes --- the list is
the transitive closure of the citations in \S\ref{sec:induction}, not only its
immediate ones --- since the point of this theorem is that none of it carries a
characteristic hypothesis.
\begin{itemize}\itemsep1pt
\item $(*^\circ)$, and with it from \S\ref{sec:reps}: Definition
\ref{def:star}, Lemma \ref{lem:Jinv}, Corollary \ref{cor:reduce}, Proposition
\ref{prop:reps}, Definition \ref{def:B} and \eqref{eq:Brig}.  These are
characteristic-free, as the head of \S\ref{sec:reps} records, which is why
quoting them at $p=2$ is legitimate.  (That is everything the induction
consumes from \S\ref{sec:reps}.  The one further result of that section,
Proposition \ref{prop:sharp}, is a consequence of $(*^\circ)$ rather than an
input to it, and is not consumed here.)
\item From \S\ref{sec:profile}: Lemma \ref{lem:P}, Lemma \ref{lem:C}, and Lemma
\ref{lem:nodefree}, the last supplying the node column of a top truncation in
the parameters $v'_t$ and $D_t(X)$ and replacing Lemma \ref{lem:node} --- a
statement of \S\ref{sec:families}, made for the nine representatives and for
$\operatorname{char}\ne2$ --- at every point of use.  Also Definition
\ref{def:upset} (the threshold $D_j(X)$), Theorem \ref{thm:lemA} and Lemma
\ref{lem:Id}.
\item Through the proof of Lemma \ref{lem:P}: the display \eqref{eq:foldedP},
which the proof of Lemma \ref{lem:cumbound} quotes directly, and with it
Proposition \ref{prop:foldsol}\textup{(a)}, Lemma \ref{lem:fold}, Lemma
\ref{lem:halfline} and Lemma \ref{lem:rigid}.  All are proved from Lemma
\ref{lem:R}, Lemma \ref{lem:soctop} and Proposition \ref{prop:ext1} by
integer combinatorics, over an arbitrary field.
\item Remark \ref{rem:repeats}, which is what Case $0$ of the induction in the
proof of Proposition \ref{prop:T} rests on: a colour-$t$ socle step is
idempotent.  It is a consequence of Lemma \ref{lem:cb} alone.
\item From \S\S\ref{sec:setup}--\ref{sec:main}: Convention \ref{conv:D}, the
sign convention and the relation \eqref{eq:relj} of \S\ref{sec:setup},
Definition \ref{def:ctrace}, Lemma \ref{lem:soctop}, Lemma \ref{lem:R}, Lemma
\ref{lem:cbseq} and Proposition \ref{prop:ext1}, Lemma \ref{lem:cb}, Corollary
\ref{cor:mono}, Corollary \ref{cor:pol}, Lemma \ref{lem:extend}, Theorem
\ref{thm:main} with its equation \eqref{eq:lower}, and Corollary
\ref{cor:prefix}.
\item Proposition \ref{prop:dimI} --- equivalently Corollary
\ref{cor:dimI}, which is where it is proved --- and behind it Proposition
\ref{prop:Pifacts}\textup{(b)}.  These enter at four places:
through Lemma \ref{lem:P} and Proposition \ref{prop:foldsol}; through Theorem
\ref{thm:main}; through Corollary \ref{cor:pol} and Lemma \ref{lem:cb} (whose
proofs use the injectivity of $I_\ka$ and $\soc I_\ka=S_\ka$ --- both
definitional, and not part of the citation); and through Lemma \ref{lem:ab}
and Lemma \ref{lem:soctop}.  Consequently the only input of the present
theorem that is quoted rather than proved is \textup{(b)}, a citation to
\cite{BBK} consumed only through the three items \textup{(b1)}--\textup{(b3)}
itemised in Remark \ref{rem:Pifacts}.
\end{itemize}
Nothing from \S\ref{sec:strands}, \S\ref{sec:T012}, \S\ref{sec:families} or
\S\ref{sec:closure} enters --- those four sections reach $p=2$ only through
their consequence $(B)$, supplied by Step 4 --- and neither does Proposition
\ref{prop:ST} or Proposition \ref{prop:grading}.
\S\ref{sec:reps}, which lies between them, \emph{does} enter, and carries no
characteristic hypothesis.

Theorem \ref{thm:typeD} follows from Proposition
\ref{prop:T} exactly as in \S\ref{sec:profile}: the greedy top truncation of
$(1,2^{k-1},1)$ to total $v_t$ has all multiplicities $\le1$ iff $v_t\le2$,
and $v_t=r+2m$ by Lemma \ref{lem:C}.  Finally, running Steps 2--4 for an
arbitrary extremal $\mu$ rather than a representative gives
\eqref{eq:specdims} for every $\mu\in W\la$, which is the last assertion.  For
such a $\mu$, Step 2 is supplied not by Theorem \ref{thm:B} but by Theorem
\ref{thm:main}: over $\mathbb{Q}$ it constructs the submodule
$V_{\mathbb{Q}}(\mu)=\soc_\sigma(I_\ka(\mathbb{Q}))$ of dimension vector
$\la-\mu$, graded, with no hypothesis on the field; Corollary \ref{cor:pol}
then says only that it is the \emph{unique} such submodule, which is all that
is asked of it.  (The existence half of Corollary \ref{cor:pol} is the clause
``over $\mathbb{C}$ there is exactly one'', which is not available over
$\mathbb{Q}$ without the geometry, and is not used here --- exactly as the
paragraph following Corollary \ref{cor:pol} records.)
\end{proof}

\subsection{The \texorpdfstring{$\mathbb{Z}$}{Z}-form of a chain slot}
\label{ss:zslot}

Step 1 of the proof just given produced the $\mathbb{Z}$-form
$\Lambda=\Pia_{\mathbb{Z}}e_k$ and showed that every graded slot
$\Lambda_i[d]$ is a free $\mathbb{Z}$-module whose rank is
$\dim_{\mathbb{K}}(I_\ka(\mathbb{K}))_i[d]$ for \emph{every} field
$\mathbb{K}$.  We now determine $\Lambda_i[d]$ at a chain vertex as a module
over $\mathbb{Z}[\tau]=\mathbb{Z}[C_2]$.  Nothing in the rest of the paper
depends on this subsection; it is here because the answer is exactly what
governs which shadows of \S\ref{sec:strands} survive at $p=2$, and because it
is not visible in the rational data.

The point at issue is the following.  A rank-$2$ $\mathbb{Z}[C_2]$-lattice
whose rational character is $\eps_+\oplus\eps_-$ is either the regular
representation $\mathbb{Z}[C_2]$, which has a $\mathbb{Z}$-basis $\{e,f\}$
with $\tau e=f$, or the split lattice
$\mathbb{Z}_+\oplus\mathbb{Z}_-$; and these are not distinguished by the
character.  They are distinguished mod $2$ --- $\tau$ is a single Jordan block
in the first case and the identity in the second --- and equivalently by the
index of the sublattice spanned by the two primitive rational eigenvectors,
which is $2$ for $\mathbb{Z}[C_2]$ and $1$ for
$\mathbb{Z}_+\oplus\mathbb{Z}_-$.  Only in the regular case do the two
rational eigenlines have the \emph{same} reduction mod $2$.

Two integral facts drive everything, and neither survives passage to
$\mathbb{Q}$.

\begin{remark}[$\tau$ carries no sign twist]\label{rem:tautwist}
With the convention $\varsigma(i,j)=\operatorname{sign}(j-i)$ of
\S\ref{sec:setup} we have
$\varsigma(t,n-1)=\varsigma(t,n)=+1$ and
$\varsigma(n-1,t)=\varsigma(n,t)=-1$, while $\varsigma$ is unchanged on chain
arrows; so $\varsigma(\tau i,\tau j)=\varsigma(i,j)$ for every arrow.  Hence
$\tau$ carries the relation at $i$ to the relation at $\tau i$ on the nose,
and is an automorphism of $\Pia_{\mathbb{Z}}$ over $\mathbb{Z}$ fixing $e_\ka$.
In particular $\tau$ \emph{permutes paths}, and $\tau a_{j\to i}=a_{\tau j\to
\tau i}\tau$ on $\Lambda$ with no intervening sign.
\end{remark}

\begin{lemma}[generation]\label{lem:gen}
For every vertex $i$ and every $d\ge1$,
$\Lambda_i[d]=\sum_{j\sim i}a_{j\to i}\Lambda_j[d-1]$.
\end{lemma}

\begin{proof}
$\Lambda=\Pia_{\mathbb{Z}}e_\ka$ is generated by $e_\ka$ and graded by path
length, so $\Lambda_i[d]$ is spanned by the images of the paths of length $d$
from $\ka$ to $i$, and every such path is an arrow $a_{j\to i}$ applied to a
path of length $d-1$.
\end{proof}

\begin{lemma}[the arrows toward the node are integrally surjective]
\label{lem:chainsurj}
\textup{(i)}  Let $1\le j\le t-1$ and $d\ge1$ with $d\ne j-k$.  Then
\[
    a_{j+1\to j}\bigl(\Lambda_{j+1}[d-1]\bigr)=\Lambda_j[d].
\]
\textup{(ii)}  Let $d\ge1$ with $d\ne t-k$.  Then
\[
    \Lambda_t[d]=a_{n-1\to t}\bigl(\Lambda_{n-1}[d-1]\bigr)
                +a_{n\to t}\bigl(\Lambda_{n}[d-1]\bigr),
\]
that is, the node slot is generated by the two leaves alone.
\end{lemma}

\begin{proof}
For $1\le j\le t-1$ and $d\in\mathbb{Z}$ write $Q(j,d)$ for the assertion of
(i), and for $d\in\mathbb{Z}$ write $N(d)$ for the assertion of (ii).  Only
the preprojective relations and Lemma \ref{lem:gen} are used; no exactness and
no hypothesis on a field.

\emph{Trivial cases.}  $Q(j,d)$ holds whenever $\Lambda_j[d]=0$: in particular
for $d<0$, and for $d=0$ with $j\ne k$, since $\Lambda_j[0]=\mathbb{Z}e_\ka$
or $0$ according as $j=k$ or not.

\emph{Base.}  $Q(1,d)$ holds for every $d\ge1$: the vertex $1$ has $2$ as its
only neighbour, so Lemma \ref{lem:gen} reads
$\Lambda_1[d]=a_{2\to1}\Lambda_2[d-1]$.

\emph{Step.}  Let $2\le j\le t-1$ and $d\ge1$.  Then $Q(j-1,d-1)$ implies
$Q(j,d)$.  Indeed by Lemma \ref{lem:gen}
$\Lambda_j[d]=a_{j-1\to j}\Lambda_{j-1}[d-1]+a_{j+1\to j}\Lambda_{j+1}[d-1]$,
so it suffices to see that the first summand lies in the second.  Let
$w\in\Lambda_{j-1}[d-1]$.  By $Q(j-1,d-1)$ we may write $w=a_{j\to j-1}v$ with
$v\in\Lambda_j[d-2]$.  The preprojective relation at the chain vertex $j$,
applied to $v$, reads
\[
   \varsigma(j,j-1)\,a_{j-1\to j}a_{j\to j-1}v
  +\varsigma(j,j+1)\,a_{j+1\to j}a_{j\to j+1}v=0 ,
\]
whence
$a_{j-1\to j}w=-\varsigma(j,j-1)\varsigma(j,j+1)\,
 a_{j+1\to j}\bigl(a_{j\to j+1}v\bigr)
 \in a_{j+1\to j}\Lambda_{j+1}[d-1]$.

\emph{Step at the node.}  Likewise $Q(t-1,d-1)$ implies $N(d)$ for $d\ge1$: by
Lemma \ref{lem:gen}, $\Lambda_t[d]$ is the sum of the images of
$a_{t-1\to t}$, $a_{n-1\to t}$ and $a_{n\to t}$; writing
$w\in\Lambda_{t-1}[d-1]$ as $a_{t\to t-1}v$ and applying the relation at $t$,
\[
  -a_{t-1\to t}a_{t\to t-1}v+a_{n-1\to t}a_{t\to n-1}v+a_{n\to t}a_{t\to n}v=0,
\]
puts $a_{t-1\to t}w$ in the sum of the two leaf images.

\emph{Descent.}  Iterating the Step from $(j,d)$ gives
$Q(j-m,d-m)\Rightarrow Q(j,d)$ for every $m\le\min(j-1,d)$ --- the $m$-th
application of the Step is at degree $d-m+1$, which must be $\ge1$.  If $d\ge j$ we reach $(1,d-j+1)$ with
$d-j+1\ge1$ and invoke the Base.  If $d<j$ we reach $(j-d,0)$, and
$\Lambda_{j-d}[0]=0$ unless $j-d=k$; so $Q(j,d)$ holds unless $d=j-k$.  This
is (i).  For (ii), $N(d)$ follows from $Q(t-1,d-1)$, which by the same count
holds unless $d-1=(t-1)-k$, i.e.\ unless $d=t-k$.
\end{proof}

The excluded diagonal is harmless: $d=j-k$ is the bottom degree
$\lo^{+}_j=|j-k|$ of the $+$ strand when $j\ge k$, a slot that is fed from
$j-1$ and not from $j+1$, and it never carries multiplicity $2$.

\begin{proposition}[the $\mathbb{Z}$-form at a chain slot]\label{prop:zslot}
Let $1\le j\le t$ and let $d$ be a degree.
\begin{enumerate}
\item If $\operatorname{rk}\Lambda_j[d]=2$ then $\Lambda_j[d]\cong
\mathbb{Z}[C_2]$ as a $\mathbb{Z}[\tau]$-module: it has a $\mathbb{Z}$-basis
$\{x,\tau x\}$.
\item If $\operatorname{rk}\Lambda_j[d]=1$ then $\Lambda_j[d]\cong
\mathbb{Z}_{\eps}$, where $\eps$ is the unique strand carried by the slot.
\end{enumerate}
Consequently $\tau$ reduces mod $2$ to a single Jordan block at every
multiplicity-$2$ chain slot and to the identity at every $1$-dimensional one;
at a multiplicity-$2$ slot the two rational eigenlines
$\langle x+\tau x\rangle$ and $\langle x-\tau x\rangle$ have the same
reduction mod $2$, namely $\ker(\tau-1)=\im(\tau-1)$, and the sublattice they
span has index $2$.
\end{proposition}

\begin{proof}
(2) is immediate: an involution of $\mathbb{Z}$ is $\pm1$, and the sign is the
eigenvalue of $\tau$ on the slot over $\mathbb{Q}$, which is $\eps$ by
Proposition \ref{prop:W2}.

For (1), recall from Proposition \ref{prop:W2} that the multiplicity-$2$
degrees at the chain vertex $j$ are exactly
\begin{equation}\label{eq:window}
   d\in W_j=[\,2t+2-j-k,\ j+k-2\,],\qquad d\equiv j+k\ (2),
\end{equation}
and that by Step 1 of the proof of Theorem \ref{thm:charp} these are also the
$d$ with $\operatorname{rk}\Lambda_j[d]=2$.

\emph{(a) The window moves inward.}  If $d\in W_j$ and $j\le t-1$ then
$d-1\in W_{j+1}$: the lower bound $2t+2-(j+1)-k\le d-1$ is
$d\ge2t+2-j-k$, and the upper bound $d-1\le (j+1)+k-2$ is $d\le j+k$, both
implied by \eqref{eq:window}; the parity is right because $d-1\equiv
(j+1)+k$.  Iterating, $(t,\,d^{\,\circ})$ with $d^{\,\circ}:=d-(t-j)$ is a
multiplicity-$2$ node slot.

\emph{(b) The chain arrows are surjective along the way.}  For
$j\le j'\le t-1$ put $d'=d-(j'-j)$.  Lemma \ref{lem:chainsurj}(i) applies at
$(j',d')$: its degree hypothesis holds because
$d'\ge d^{\,\circ}\ge t+2-k\ge2$, and its exclusion $d'=j'-k$ would read
$2j'=d+j+k\ge2t+2$ by \eqref{eq:window}, forcing $j'\ge t+1$.  Hence
\[
   \Lambda_j[d]=\bigl(a_{j+1\to j}\circ\dots\circ a_{t\to t-1}\bigr)
                \bigl(\Lambda_t[d^{\,\circ}]\bigr),
\]
a surjection of $\mathbb{Z}[\tau]$-modules, $\tau$ fixing every chain vertex.

\emph{(c) The node slot is generated by one $\tau$-orbit.}  Since
$d^{\,\circ}\in W_t=[\,t+2-k,\ t+k-2\,]$ we have $d^{\,\circ}\ne t-k$, so
Lemma \ref{lem:chainsurj}(ii) gives
$\Lambda_t[d^{\,\circ}]=a_{n-1\to t}\Lambda_{n-1}[d^{\,\circ}-1]
+a_{n\to t}\Lambda_{n}[d^{\,\circ}-1]$.  By Lemma \ref{lem:P} the two leaves
are multiplicity-free with support $t-k+1,t-k+3,\dots,t+k-1$, and
$d^{\,\circ}-1$ lies in that range with the right parity; so
$\Lambda_{n-1}[d^{\,\circ}-1]=\mathbb{Z}u$ and
$\Lambda_{n}[d^{\,\circ}-1]=\mathbb{Z}\tau u$ are of rank $1$, and
$J:=\Lambda_{n-1}[d^{\,\circ}-1]\oplus\Lambda_{n}[d^{\,\circ}-1]$ has the
$\mathbb{Z}$-basis $\{u,\tau u\}$ permuted by $\tau$, i.e.\
$J\cong\mathbb{Z}[C_2]$.  By Remark \ref{rem:tautwist} the map
$\varphi:=a_{n-1\to t}+a_{n\to t}:J\to\Lambda_t[d^{\,\circ}]$ is
$\tau$-equivariant, and it is onto.  So
$\Lambda_t[d^{\,\circ}]$ is a quotient of $\mathbb{Z}[C_2]$; both have rank
$2$, so the kernel of $\varphi$ has rank $0$ and, sitting inside the
torsion-free $J$, is zero.  Hence $\Lambda_t[d^{\,\circ}]\cong\mathbb{Z}[C_2]$,
with basis $\{\varphi(u),\tau\varphi(u)\}$.

\emph{(d) Descent along the chain.}  By (b), $\Lambda_j[d]$ is a
$\tau$-equivariant quotient of $\Lambda_t[d^{\,\circ}]\cong\mathbb{Z}[C_2]$ of
rank $2$; by the same argument the quotient map is an isomorphism.  This
proves (1).

The final sentence is the standard dictionary recalled above: in the basis
$\{x,\tau x\}$ the matrix of $\tau$ is $\begin{psmallmatrix}0&1\\1&0
\end{psmallmatrix}$, which mod $2$ is a single Jordan block with
$\ker(\tau-1)=\im(\tau-1)=\langle\overline{x+\tau x}\rangle$; and
$\det\begin{psmallmatrix}1&1\\1&-1\end{psmallmatrix}=-2$.
\end{proof}

\begin{remark}
The proof uses no exactness and no field: only that $\Lambda$ is generated by
$e_\ka$ in degree $0$, that the preprojective relations have coefficients
$\pm1$ with $\varsigma(\tau i,\tau j)=\varsigma(i,j)$, that the slot ranks are
those of Proposition \ref{prop:W2} and Lemma \ref{lem:P}, and that the slots
are $\mathbb{Z}$-free (Step 1 of Theorem \ref{thm:charp}).  The two places
where integrality --- as opposed to the rational character --- enters are
Remark \ref{rem:tautwist} and the surjectivity of Lemma \ref{lem:chainsurj},
which is a statement about $\mathbb{Z}$-modules and not merely about ranks.
Proposition \ref{prop:zslot} was also verified directly, by an independent
integral construction of $\Lambda$ with Smith normal forms in place of row
reduction; see \S\ref{ss:checks}.
\end{remark}

\begin{remark}[what the argument does and does not give]
It gives the theorem at $p=2$; it gives no new information there.  Being a
specialisation argument, it shows only that the $p=2$ answer is the shadow of
the characteristic-$0$ one.  It does \emph{not} make (W1)--(W4), (T0)--(T2) or
the canonical truncation $T(v)$ true at $p=2$ intrinsically.

Some of them do survive as shadows, and we record the mechanism.  By
Proposition \ref{prop:zslot} the $\mathbb{Z}$-form of a multiplicity-$2$ chain
slot is the regular representation $\mathbb{Z}[C_2]=\langle e,f\rangle$ with
$\tau$ swapping $e$ and $f$; hence the two rational eigenlines
$\langle e+f\rangle$ and $\langle e-f\rangle$ have the \emph{same} reduction
mod $2$, namely $\ker(\tau-1)=\im(\tau-1)$.  So $T(v)$ remains definable at
$p=2$ with ``$\tau$-eigenline'' read as $\ker(\tau-1)$, and (T0)--(T2) then
hold verbatim.

This is not formal.  A rank-$2$ $\mathbb{Z}[C_2]$-lattice whose rational
character is $\eps_+\oplus\eps_-$ is either the regular representation
$\mathbb{Z}[C_2]$, whose reduction mod $2$ is a single Jordan block, or the
split lattice $\mathbb{Z}_+\oplus\mathbb{Z}_-$, whose reduction is the
identity, and the rational data do not distinguish them; in the split case the
two eigenlines would have \emph{different} reductions and the previous
paragraph would fail.  Which one occurs is decided in \S\ref{ss:zslot}, by an
integral argument --- the surjectivity of Lemma \ref{lem:chainsurj} and the
absence of a sign twist in $\tau$ (Remark \ref{rem:tautwist}) --- and not by
any rational input.  Nothing in this paper depends on the outcome: Theorem
\ref{thm:charp} goes through the $\mathbb{Z}$-form slotwise and never asks
what $\tau$ does; and, as the next remark shows, even with Proposition
\ref{prop:zslot} in hand one cannot rebuild \S\ref{sec:strands}.
\end{remark}

\begin{remark}[the filtration route does not work]\label{rem:filtfail}
It is natural to try to replace the eigen-splitting at $p=2$ by the filtration
$0\le\ker(\tau-1)\le V$, which is available because
$(\tau-1)^2=2(1-\tau)=0$ there, and to read the associated graded as a
substitute for the two strands.  This fails, and it fails at the level of
(W2), so it is worth recording so that no one repeats the attempt.

The obstruction sits at the $1$-dimensional slots, and it does not depend on
any unproved input.  At a $1$-dimensional slot the only unipotent operator on
a line is the identity, so the filtration grading is $(1,0)$ \emph{always}.
(At a multiplicity-$2$ slot the filtration grading is $(1,1)$, by
Proposition \ref{prop:zslot}; the argument below does not use this.)  Over $\mathbb{Q}$ a
$1$-dimensional chain slot carries one
strand only, with eigen-dimensions $(1,0)$ or $(0,1)$ according to which.  The
filtration grading is therefore wrong at every ``$-$ only'' slot, whatever it
does at the multiplicity-$2$ slots.  These are not
rare: over $D_4$--$D_{14}$ and all $k$ there are $3575$ nonzero chain slots, of
which $1001$ have multiplicity $2$, $1287$ are ``$+$ only'' and $1287$ are
``$-$ only'' --- and the filtration misreads all $1287$ of the last kind.

Consequently Proposition \ref{prop:W2} is false for the filtration grading in
either assignment of sub- and quotient to $+$ and $-$.  Writing $f^0_d(j)$ and
$f^1_d(j)$ for the two graded pieces, $f^0$ is supported on
$S^{+}(j)\cup S^{-}(j)$ and $f^1$ on the window
$W_j=S^{+}(j)\cap S^{-}(j)$; both are symmetric under
$d\mapsto2t-d$, since $S^{+}(j)\cup S^{-}(j)$ and $W_j$ are, whereas
$a^{+}$ is not, so the reversal symmetry $a^{-}_d(j)=a^{+}_{2t-d}(j)$ of (W2)
fails.  The casualty downstream is (T2), whose proof works by
showing that a $1$-dimensional target carries $-$ and not $+$: over
$\mathbb{F}_2$ the filtration cannot see that label at all, since $\tau$ acts
trivially on every line, and (T2) is not even statable in filtration language.
This is a loss of information, not a gap in an argument, and it is why
Theorem \ref{thm:charp} goes through the $\mathbb{Z}$-form instead.
\end{remark}

\section{A word-free complement: the staircase property}\label{sec:staircase}

Proposition \ref{prop:T} says that the \emph{extremal} module $V$ is a top
truncation at $t$.  There is a stronger, word-free statement, which we prove
here because it explains why extremality is the right hypothesis and because it
is the form in which the result is easiest to check.  Nothing in
\S\S\ref{sec:setup}--\ref{sec:charp} depends on this section.

\begin{definition}\label{def:staircase}
A graded submodule $M\subseteq I_\ka$ is a \emph{staircase} if it is a top
truncation at $t$ and an up-set at each of the two leaves.
\end{definition}

\begin{theorem}[staircase property]\label{thm:staircase}
Over an arbitrary field: if $\dimv M$ is $w$-extremal then $M$ is a
staircase.  Equivalently, a graded submodule that is not a staircase satisfies
$\delta(\dimv M)>0$ and is therefore never the unique submodule of its
dimension vector.
\end{theorem}

\begin{proof}[Proof of Theorem \ref{thm:staircase}]
Let $v=\dimv M$ be extremal, say $v=\la-w\la$.  By Corollary \ref{cor:pol} the
submodule of $I_\ka$ with dimension vector $v$ is unique, so $M=V$; and by
Proposition \ref{prop:T} together with Lemma \ref{lem:leafstep} (applied along
the filtration of Theorem \ref{thm:main}, which ends at $V$ and preserves
$(H)$), $V$ is a top truncation at $t$ and an up-set at each leaf.  The
equivalent form follows from Corollary \ref{cor:pol}.  Both inputs are
available over every field --- Corollary \ref{cor:pol} directly and
Proposition \ref{prop:T} by Theorem \ref{thm:charp} --- so no hypothesis on
the characteristic is needed here.  The remaining results of this section do
carry one; see Remark \ref{rem:sec9char}.
\end{proof}

The top-truncation half of Theorem \ref{thm:staircase} has, for leaf-balanced
dimension vectors, a direct proof independent of the whole machinery of
\S\S\ref{sec:reps}--\ref{sec:induction}; that is Proposition \ref{prop:tau}
below, and it is the reason for Lemma \ref{lem:U}.

\begin{remark}\label{rem:sec9char}
Lemma \ref{lem:U0}, Lemma \ref{lem:U} and Proposition \ref{prop:tau} below
are stated for $\operatorname{char}\ne2$ and are \emph{not} transported to
$p=2$ by Theorem \ref{thm:charp}: they are not consequences of $(B)$, and the
determinant \eqref{eq:det} in the proof of Lemma \ref{lem:U} genuinely
degenerates at $p=2$, the binary quadratic form
$\lambda^2a_+m_+-\nu^2a_-m_-$ becoming a square there.  Since this whole
section is a complement on which nothing else depends, we leave it so.
\end{remark}

\subsection{The node double-path lemma}
For $x\in(I_\ka)_t[d]$ set
\[
   A(x)=a_{(t-1)\to t}a_{t\to(t-1)}x,\quad
   B(x)=a_{(n-1)\to t}a_{t\to(n-1)}x,\quad
   C(x)=a_{n\to t}a_{t\to n}x,
\]
and $\Theta(x)=\operatorname{span}\{A(x),B(x),C(x)\}\subseteq(I_\ka)_t[d+2]$.
With the sign convention $\varsigma(i,j)=\operatorname{sign}(j-i)$ the preprojective
relation at $t$ reads
\begin{equation}\label{eq:Rt}
   -A(x)+B(x)+C(x)=0,\qquad\text{i.e.}\qquad P(x):=B(x)+C(x)=A(x),
\end{equation}
so $\Theta(x)=\operatorname{span}\{A(x),B(x)\}$ has dimension at most $2$.  Set
also $Q(x):=B(x)-C(x)$; since $\operatorname{char}\ne2$,
$\operatorname{span}\{B(x),C(x)\}=\operatorname{span}\{P(x),Q(x)\}$.  Call
$x$ \emph{special} if some leaf arrow kills it.

\begin{lemma}[interval facts]\label{lem:intervals}
Let $2\le k\le t$ and recall \eqref{eq:supports}.
\begin{itemize}\itemsep1pt
\item[\textup{(I1$'$)}] If $\eps$ is present at $(t,d)$ and at $(t,d+2)$, then
$\eps$ is present at $(t-1,d+1)$.  \emph{No hypothesis on $\dim(I_\ka)_t[d]$.}
\item[\textup{(I3)}] If $(I_\ka)_t[d]\ne0$ and $\dim(I_\ka)_t[d+2]=1$, then
$d=t+k-2$, the surviving strand at $(d+2,t)$ is $-$, and at $(t-1,d+1)$ the
strand $-$ is present and $+$ is not.
\item[\textup{(I4$'$)}] If $(I_\ka)_t[d]\ne0$ and $(I_\ka)_t[d+2]\ne0$ then
both leaf slots at degree $d+1$ are nonzero, hence lines.
\item[\textup{(I5)}] The node slots with $\dim(I_\ka)_t[d]=1$ and
$(I_\ka)_t[d+2]\ne0$ are exactly $d=t-k$; there the $+$ strand alone is present
and both strands are present at $d+2$.
\end{itemize}
\end{lemma}

\begin{proof}
(I1$'$), plus strand: presence at $(t,d+2)$ gives $d+2\le t+k-2$, hence
$d+1\le(t-1)+k-2$; presence at $(t,d)$ gives $d\ge|t-k|$, and
$|t-1-k|\le|t-k|+1$, hence $d+1\ge|t-1-k|$; the parity is right since
$d\equiv t+k$ gives $d+1\equiv(t-1)+k$.  Minus strand: $-$ is present at $(j,e)$
iff $+$ is present at $(j,2t-e)$, so putting $e=2t-d-2$ the hypothesis becomes
``$+$ at $(t,e)$ and $+$ at $(t,e+2)$'' and the conclusion ``$+$ at $(t-1,e+1)$'',
which is the plus case.  Nothing in this argument inspects the other strand, so
it is valid at a slot of dimension $1$.

(I3): $\dim(I_\ka)_t[d+2]=1$ forces $d+2\in\{t-k,t+k\}$; the value $t-k$ gives
$d=t-k-2$, below the node support, contradicting $(I_\ka)_t[d]\ne0$.  So
$d+2=t+k$ and the survivor is the $-$ strand (as $\hi^+_t=t+k-2$), and
$d=t+k-2$.  Then $d+1=t+k-1>t+k-3=\hi^+_{t-1}$, so $+$ is absent at
$(t-1,d+1)$; and $t-k+3\le t+k-1\le 2t-|t-1-k|$ for $k\ge2$, so $-$ is present.

(I4$'$): the node support is $[t-k,\,t+k]$ and, by Lemma \ref{lem:P}, the leaf
support is $[t-k+1,\,t+k-1]$.  Now $(I_\ka)_t[d]\ne0$ gives $d\ge t-k$, hence
$d+1\ge t-k+1$; and $(I_\ka)_t[d+2]\ne0$ gives $d+2\le t+k$, hence
$d+1\le t+k-1$.  The parity is right: $d\equiv t+k$ gives $d+1\equiv t+k+1$,
the leaf parity.

(I5): $\dim(I_\ka)_t[d]=1$ means $d\in\{t-k,t+k\}$.  At $d=t+k$ the slot $d+2$
lies beyond the node support, so $(I_\ka)_t[d+2]=0$.  At $d=t-k$ only $+$ is
present, and $d+2=t-k+2$ lies in both supports precisely when $k\ge2$.
\end{proof}

\begin{lemma}\label{lem:U0}
Assume $\operatorname{char}\ne2$ and $2\le k\le n-2$.  Let $d$ satisfy
$(I_\ka)_t[d]\ne0$ and $(I_\ka)_{n-1}[d+1]\ne0$ \textup{(}equivalently, by
$\tau$, $(I_\ka)_n[d+1]\ne0$\textup{)}.  Then:
\begin{enumerate}
\item[\textup{(i)}] no $\tau$-eigenline of $(I_\ka)_t[d]$ is special;
\item[\textup{(ii)}] if $\dim(I_\ka)_t[d]=2$, the two kernels
$\ker a_{t\to(n-1)}$ and $\ker a_{t\to n}$ are distinct lines, and the special
vectors are exactly their union --- two lines, never a subspace;
\item[\textup{(iii)}] if $\dim(I_\ka)_t[d]=1$, no nonzero $x\in(I_\ka)_t[d]$ is
special.
\end{enumerate}
\end{lemma}

\begin{proof}
(i) is Proposition \ref{prop:W4n}, applied to each leaf in turn.  (ii): by
Lemma \ref{lem:P} each leaf slot is at most a line, so each kernel of a map from
a two-dimensional space contains a line; if the two kernels shared a line $S$
then, $\tau$ exchanging the two arrows out of $t$, $S$ would be $\tau$-stable,
hence a $\tau$-eigenline that is special, contradicting (i).  Distinctness
forces each kernel to be exactly a line.  (iii): a one-dimensional slot is
$\tau$-stable, hence an eigenline; apply (i).
\end{proof}

\begin{lemma}[node double-path lemma]\label{lem:U}
Assume $\operatorname{char}\ne2$ and $2\le k\le n-2$, and let $d$ satisfy
$(I_\ka)_t[d]\ne0$ and $(I_\ka)_t[d+2]\ne0$.  Exactly one of the following
three shapes occurs, and in each:
\begin{itemize}\itemsep1pt
\item[\textup{(U1)}] $\dim(I_\ka)_t[d]=\dim(I_\ka)_t[d+2]=2$: then
$\Theta(x)=(I_\ka)_t[d+2]$ for every non-special $x$, and $\dim\Theta(x)\le1$
on the two special lines;
\item[\textup{(U2)}] $\dim(I_\ka)_t[d]=2$, $\dim(I_\ka)_t[d+2]=1$: then
$\Theta(x)=(I_\ka)_t[d+2]$ for every $x\ne0$, special or not;
\item[\textup{(U3)}] $\dim(I_\ka)_t[d]=1$ --- necessarily $d=t-k$ and
$\dim(I_\ka)_t[d+2]=2$: then $\Theta(x)=(I_\ka)_t[d+2]$ for every $x\ne0$, and
no $x\ne0$ is special.
\end{itemize}
\end{lemma}

\begin{proof}
That these are the only shapes is (I5).  By (I4$'$) the leaf slots at $d+1$ are
nonzero, so Lemma \ref{lem:U0} is available at $d$ throughout.

\emph{Cases (U1), (U2): $\dim(I_\ka)_t[d]=2$.}  Let $u^{+},u^{-}$ be a
$\tau$-eigenbasis of $(I_\ka)_t[d]$, and $w^{+},w^{-}$ eigenvectors spanning
the nonzero eigen-parts of $(I_\ka)_t[d+2]$.  Both $P=A$ and $A$ commute with
$\tau$, so $P$ preserves the eigen-grading, while $Q$ anticommutes with it and
reverses it.  Write
\[
   A(u^{+})=a_+w^{+},\quad A(u^{-})=a_-w^{-},\quad
   Q(u^{+})=m_+w^{-},\quad Q(u^{-})=m_-w^{+},
\]
with the convention that a coefficient whose target eigenline is absent is
read as $0$.

\emph{Step 1 (the $a$'s).}  Let $\eps$ be present at $(t,d)$ and at $(t,d+2)$.
By (I1$'$), $\eps$ is present at $(t-1,d+1)$, and then Corollary \ref{cor:W4}
makes both $a_{t\to(t-1)}$ and $a_{(t-1)\to t}$ isomorphisms on the
$\eps$-line: $a_\eps\ne0$.

\emph{Step 2 (the $m$'s).}  Put $y:=a_{t\to(n-1)}u^{\eps}$.  By Lemma
\ref{lem:U0}(i) the eigenline $\langle u^\eps\rangle$ is not special, so
$y\ne0$; and by (I4$'$) the folded leaf space $J[d+1]$ is the induced
$\mathbb{Z}/2$-representation on a line, so each of $J[d+1]^{\pm}$ is a line.
From $\tau u^{\eps}=\eps u^{\eps}$ and
$\tau a_{t\to(n-1)}=a_{t\to n}\tau$ one gets
$a_{t\to n}u^{\eps}=\eps\,\tau(y)$, whence
\[
   L(u^{\eps})=\varphi_{\eps}\bigl((y,\eps\,\tau y)\bigr),
   \qquad
   Q(u^{\eps})=\varphi_{-\eps}\bigl((y,-\eps\,\tau y)\bigr),
\]
where $\varphi_{\delta}:J[d+1]^{\delta}\to(I_\ka)_t[d+2]^{\delta}$ is the
folded leaf map.  The vector $(y,-\eps\tau y)$ is a nonzero element of
$J[d+1]^{-\eps}$, hence spans it, so Proposition \ref{prop:W4l} gives
\[
   m_{\eps}\ne0\iff \text{the strand } -\eps \text{ is present at } (d+2,t).
\]

\emph{Step 3 (the determinant).}  For $x=\lambda u^{+}+\nu u^{-}$,
\[
   P(x)=\lambda a_+w^{+}+\nu a_-w^{-},\qquad
   Q(x)=\nu m_-w^{+}+\lambda m_+w^{-},
\]
so in the basis $(w^+,w^-)$,
\begin{equation}\label{eq:det}
   \det\bigl(P(x),Q(x)\bigr)=\lambda^2a_+m_+-\nu^2a_-m_- .
\end{equation}

In case (U1) both strands are present at $(d+2,t)$, so Steps 1--2 make
$a_\pm,m_\pm$ all nonzero, and \eqref{eq:det} is a binary quadratic form with
both coefficients nonzero: its zero locus is a union of at most two lines.
Every special $x$ lies in that locus, since $B(x)=0$ or $C(x)=0$ forces
$\dim\operatorname{span}\{B(x),C(x)\}\le1$; and by Lemma \ref{lem:U0}(ii) there
are two distinct special lines.  Hence the zero locus is exactly those two
lines, and $\dim\Theta(x)=2$ for every non-special $x$.

In case (U2), by (I3) the surviving strand at $d+2$ is $-$ and $d=t+k-2$.  So
$a_+=0$ and $m_-=0$ (their targets are the absent $+$ line), while $a_-\ne0$
(Step 1, with $-$ present at $d$ and $d+2$) and $m_+\ne0$ (Step 2, with
$-\eps=-$ present at $d+2$).  Hence $P(x)=\nu a_-w^-$ and $Q(x)=\lambda m_+w^-$,
so $\Theta(x)=\operatorname{span}\{P(x),Q(x)\}$ is the whole one-dimensional
target unless $\lambda=\nu=0$.

\emph{Case (U3): $\dim(I_\ka)_t[d]=1$.}  By (I5), $d=t-k$, the source carries
the $+$ strand alone (so $x$ spans the $+$ eigenline, $\tau x=x$), and both
strands are present at $d+2$; let $w^{\pm}$ be an eigenbasis there.  Step 1
applies with $\eps=+$, using (I1$'$), which carries no hypothesis on
$\dim(I_\ka)_t[d]$: $A(x)=a_+w^+$ with $a_+\ne0$.  Step 2 applies with
$\eps=+$, using (I4$'$), whose two hypotheses are exactly those of the present
lemma --- at $d=t-k$ the relevant leaf slot is the \emph{bottom} slot
$t-k+1$ of the leaf support --- and Lemma \ref{lem:U0}(iii): $y\ne0$, and since
the $-$ strand is present at $(d+2,t)$, Proposition \ref{prop:W4l} gives
$Q(x)=m_+w^-$ with $m_+\ne0$.  Thus $P(x)=A(x)$ and $Q(x)$ are nonzero vectors
in distinct eigenlines, hence independent, and
$\Theta(x)=(I_\ka)_t[d+2]$.  Finally no $x\ne0$ is special, by Lemma
\ref{lem:U0}(iii).
\end{proof}

\begin{proposition}\label{prop:tau}
Assume $\operatorname{char}\ne2$ and $v_{n-1}=v_n$.  If $M\subseteq I_\ka$ is a
graded submodule which is not a top truncation at $t$, then $\dimv M$ is not
extremal.  In particular the top-truncation half of Theorem \ref{thm:staircase}
holds for leaf-balanced dimension vectors, independently of
\S\S\ref{sec:reps}--\ref{sec:induction}.
\end{proposition}

\begin{proof}
Suppose $M$ is not a top truncation at $t$: say $M_t[d]\ne0$ while $M_t[d+2]$
is a proper subspace of $(I_\ka)_t[d+2]$ (in particular $(I_\ka)_t[d+2]\ne0$).
Let $0\ne x\in M_t[d]$.  Its three double paths lie in $M_t[d+2]$, so by Lemma
\ref{lem:U} $x$ must be special: otherwise $\Theta(x)=(I_\ka)_t[d+2]\subseteq
M_t[d+2]$, contradicting properness.

At every slot Lemma \ref{lem:U} reaches, a special vector forces
$\dim(I_\ka)_t[d]=2$.  Indeed the node slots of dimension $1$ are $d=t-k$ and
$d=t+k$; the second is excluded from Lemma \ref{lem:U} by (I5), because there
$(I_\ka)_t[d+2]=0$; and at the first, the leaf slots at $d+1$ are nonzero by
(I4$'$), so Lemma \ref{lem:U0}(iii) applies and no nonzero $x$ is special.

So $\dim(I_\ka)_t[d]=2$, every nonzero $x\in M_t[d]$ is special, and by Lemma
\ref{lem:U0}(ii) the special vectors are the union of the two leaf kernels ---
two distinct lines, not a subspace.  Hence $M_t[d]$ is one of those two lines.
The automorphism $\tau$ fixes $I_\ka$ and exchanges the two arrows out of $t$
to the leaves, hence exchanges the two special lines; so $\tau(M)\ne M$.  When
$v_{n-1}=v_n$ we have $\dimv\tau(M)=\dimv M$, and two distinct submodules with
one dimension vector contradict Corollary \ref{cor:pol}.  So $\dimv M$ is not
extremal.

\end{proof}

\begin{remark}
The universal statement ``a special vector exists only where
$\dim(I_\ka)_t[d]=2$'' is \emph{false}: at the top node slot $d=t+k$ both leaf
slots at $d+1=t+k+1$ vanish, the leaf top degree being $t+k-1$, so every
nonzero vector there is special, at a slot of dimension $1$.  What is true, and
all that is used, is the relative statement in the proof above: among the slots
that Lemma \ref{lem:U} reaches, a special vector forces dimension $2$.  The work
is done by (I5) and (I4$'$), not by Lemma \ref{lem:U0} alone.
\end{remark}

\section{What is verified, what is proved, and what fails}
\label{sec:verification}

\subsection{Statements asserted on computational evidence}\label{ss:evidence}
Every statement in \S\S\ref{sec:setup}--\ref{sec:staircase} is either proved
here or attributed to a cited reference.  \emph{No} statement of this paper,
and in particular no step in the proof of Theorem \ref{thm:typeD}, rests on
computational evidence.  (The counts reported in \S\ref{ss:checks} below are
assertions of fact resting on computation; they are checks on proved
statements, and nothing in \S\S\ref{sec:setup}--\ref{sec:staircase} depends on
them.)  This list is now empty; three items have left
it, and we record what became of each, because the computations offered for
them are retained below, in \S\ref{ss:checks}, as checks on the now-proved
statements.
\begin{itemize}\itemsep2pt
\item The conclusion of Theorem \ref{thm:typeD} at $\operatorname{char}=2$ is
Theorem \ref{thm:charp}.
\item The Jordan type of $\tau$ at a multiplicity-$2$ chain slot is
Proposition \ref{prop:zslot}, proved in \S\ref{ss:zslot}.
\item \textbf{(V2)}, the \emph{sharpness of the side
condition in $(*^\circ)$}: the assertion that every extremal
$M\subseteq I_\ka$ which fails to be an up-set at $t-1$ has
$\langle\al_t^\vee,\la-\dimv M\rangle=0$.  That is now Proposition
\ref{prop:sharp}, proved in \S\ref{ss:sharp} from $(*^\circ)$ together with
Lemma \ref{lem:Jinv}(b) applied at $j=t$.  It remains true that the statement
is used nowhere in the proof of Theorem \ref{thm:typeD}.
\end{itemize}

\subsection{Ranges of the independent checks}\label{ss:checks}
The following were run as independent confirmations of proved statements.  They
are checks, not steps in any proof: every count recorded below is a report of
what a computation found, and no statement of
\S\S\ref{sec:setup}--\ref{sec:staircase} rests on any of them.  Each range
below is the range the corresponding script actually covers.
\begin{itemize}\itemsep2pt
\item \emph{Proposition \ref{prop:sharp}, the sharpness of the side condition.}
Over the full Weyl orbits of $D_4$--$D_{10}$, all $2\le k\le n-2$ and all
$\mu\in W\om_k$, every extremal module was built from scratch and tested for
the up-set property at $t-1$.  Over $D_4$--$D_7$ there are $2200$ extremal
modules, of which exactly $64$ fail to be up-sets at $t-1$; over $D_4$--$D_8$
the counts are $7464$ and $244$; over the whole of $D_4$--$D_{10}$ they are
$77196$ and $2778$.  Every failure, on the whole range, has $c_t=0$, as
Proposition \ref{prop:sharp} requires.  The check also confirms that the
proposition has content: $c_t$ takes all five values $-2,\dots,2$ on extremal
modules, with the distribution
$(7537,18644,24834,18644,7537)$ over $D_4$--$D_{10}$, so $26181$ of the $77196$
extremal modules have $c_t\le-1$ --- ``$c_t\ge0$ always'' is false --- and
every one of those $26181$ is an up-set at $t-1$.  (Over $D_4$--$D_8$ the
figure is $2565$, the same $2565$ as in (N5) below; that is not a coincidence,
since $s_t$ is a bijection from $\{c_t\le-1\}$ to $\{c_t\ge1\}$ on each orbit.)
The mechanism of the proof was checked separately and slotwise over the same
range, over $\mathbb{Q}$ and over $\mathbb{F}_2$: for each such $\mu$,
$c_t(s_t\mu)=-c_t(\mu)\ge1$ and $V(\mu)$, $V(s_t\mu)$ agree at every vertex
other than $t$, with no exception in $26181$ instances.
\item \emph{Theorem \ref{thm:charp}, the specialisation to $p=2$.}  The
content of that theorem is \eqref{eq:specdims}: the graded dimensions of
$V(\mu)$ do not depend on the field.  This was checked \emph{slotwise} --- not
merely as a simple-spectrum verdict --- by comparing $V(\mu)$ over $\mathbb{Q}$
with $V(\mu)$ over $\mathbb{F}_2$ at every slot, on the \emph{full} Weyl orbits
of $D_4$--$D_{10}$, all $2\le k\le n-2$: $77196$ extremal weights, with $0$
mismatches and $0$ disagreements with the criterion $r+2m\le2$.  The per-rank
counts of simple-spectrum weights at $p=2$ are
$15,53,139,325,719,1545,3267$ for $D_4,\dots,D_{10}$, identical to the counts
over $\mathbb{Q}$.  Step 1 of the proof was checked separately: over
$D_4$--$D_{10}$ and all $k$, the $3080$ graded slot dimensions of $I_\ka$ agree
over $\mathbb{Q}$, $\mathbb{F}_2$ and $\mathbb{F}_3$, with no exception.  The
input to Step 1 that carries $p=2$, namely Corollary \ref{cor:dimI}, was
checked in both of its halves: \emph{algebraically}, by building
$\Pia_{\mathbb{K}}e_\ka$ from scratch over
$\mathbb{K}=\mathbb{Q},\mathbb{F}_2,\mathbb{F}_3,\mathbb{F}_5,\mathbb{F}_7,
\mathbb{F}_{13},\mathbb{F}_{101}$ for $D_4$--$D_{16}$ and \emph{every} vertex
$\ka$ (chain and leaf), and comparing the total dimension, the top degree and
the full folded profile with Proposition \ref{prop:foldsol} --- agreement in
every case; and \emph{combinatorially}, by verifying the rigidity statement
of Lemma \ref{lem:rigid} and the totals of Corollary \ref{cor:dimI} as
integer identities for every $D_n$ with $4\le n\le300$ and every vertex, with
no exception.
\item \emph{The socle of $\Pia e_j$, i.e.\ the part of Proposition
\ref{prop:Pifacts}\textup{(b)} that this paper consumes.}  Each $\Pia e_j$
was rebuilt from scratch --- as the quotient $F/\mathcal{R}$ of the free
module on the paths starting at $j$, taking the cokernel of the map $Y$ of
Lemma \ref{lem:R} as it comes, with no injectivity of $Y$, no self-injectivity
and no socle statement assumed --- for $D_4$ through $D_{10}$, \emph{every}
vertex $j$ including both leaves and both parities of $n$, over $\mathbb{Q}$
and over $\mathbb{F}_2,\mathbb{F}_3,\mathbb{F}_5,\mathbb{F}_7$.  In all $245$
cases $\soc(\Pia e_j)$ came out one dimensional and concentrated in the top
degree, the top degree came out $2t=2(n-2)$ at every vertex, and the total
dimension agreed with Corollary \ref{cor:dimI}.  The socle sits at vertex $j$
for every chain vertex $j\le n-2$ --- which is item \textup{(b3)} of Remark
\ref{rem:Pifacts}, the only part of the Nakayama permutation the paper uses
--- and at the leaf $j$ itself when $n$ is even, at the other leaf when $n$ is
odd, which is the refinement the paper does \emph{not} use.  The
survivals recorded in the first remark of \S\ref{sec:charp} were measured too:
over $D_4$--$D_{11}$ and all $k$, $\tau$ is a single Jordan block at each of the
$330$ multiplicity-$2$ chain slots and the identity at each of the $900$
one-dimensional ones, no chain arrow drops rank from $\mathbb{Q}$ to
$\mathbb{F}_2$, and (T0), (T1), (T2) and Proposition \ref{prop:W4n} hold at all
$1908$, $120$, $120$ and $240$ instances with $\ker(\tau-1)$ read for the
eigenline.
\item \emph{Proposition \ref{prop:zslot}, the $\mathbb{Z}$-form of a chain
slot.}  The lattice $\Lambda=\Pia_{\mathbb{Z}}e_k$ was built from scratch over
$\mathbb{Z}$ --- Smith normal forms in place of row reduction, so that the
object produced is the lattice and not a reduction of it --- together with the
integral matrix of $\tau$ in each slot, over $D_4$--$D_{22}$ and all
$2\le k\le n-2$.  Every slot came out torsion free, of the rank predicted by
Proposition \ref{prop:W2} and Lemma \ref{lem:P}.  At each of the $7315$
multiplicity-$2$ chain slots in that range, $\tau$ is a transposition of a
$\mathbb{Z}$-basis; it is not the identity mod $2$; it is not the identity
mod $4$ either, which is the test $\mathbb{F}_2$ cannot make; and the
sublattice spanned by the two primitive rational eigenvectors has index exactly
$2$.  At each of the $17290$ one-dimensional chain slots $\tau$ acts by $+1$ or
$-1$ according to the strand of Proposition \ref{prop:W2}.  The two lemmas
carrying the proof were checked separately over $D_4$--$D_{18}$ and all $k$:
the integral surjectivity of Lemma \ref{lem:chainsurj} at all $31280$
chain-arrow instances and all $2720$ node instances --- failing, as predicted,
exactly on the excluded diagonals $d=j-k$ and $d=t-k$ and nowhere else --- and
the leaf ranks of Lemma \ref{lem:P} at all $2840$ instances.  As a further
cross-check, the leaf-crossing composite
$a_{t\to n-1}a_{n\to t}\colon\Lambda_n[d-1]\to\Lambda_{n-1}[d+1]$, whose
unimodularity is equivalent to the conclusion, was found to be $\pm1$ at all
$1330$ multiplicity-$2$ node slots of $D_4$--$D_{22}$.  The census behind Remark \ref{rem:filtfail} is in closed form over
$D_4$--$D_{14}$: $3575$ nonzero chain slots, $1001$ of multiplicity $2$, $1287$
``$+$ only'' and $1287$ ``$-$ only''.
\item \emph{The block induction of Theorem \ref{thm:main}}
\textup{(}equation \eqref{eq:lower}\textup{)}.  The two bounds
$\delta_p\ge c(p)$ and $\delta_p\le c(p)$ were checked block by block on the
actual modules, over the \emph{full} Weyl orbits of $D_4$--$D_8$, all
$2\le k\le n-2$, at $p\in\{0,2,3,5\}$: $29856$ weight-runs and $519184$
blocks, every one of which satisfies $\dimv X_p=\la-\mu[p]$, enlarges $X$ only
at $i(p)$, and has $\delta_p=c(p)$ exactly, with $0$ failures.  The identity
$\delta(v+\delta\epsilon_i)-\delta(v)=2\delta\bigl(c-\delta\bigr)$ on which the
upper bound turns was checked as integer arithmetic at every extremal $v$ of
$D_4$--$D_9$, every vertex $i$ and every $-4\le\delta\le6$ ($624140$ cases, no
failure), together with $\delta(v)=0$ at each such $v$.  Recorded separately:
the drop $(Ce)_i=2>0$ that an inexact earlier block would inflict on the
pairing seen by Corollary \ref{cor:mono} --- which is why the induction, and
not a bare appeal to Corollary \ref{cor:mono}, is what establishes
\eqref{eq:lower} --- occurs at all $5074$ instances tested over $D_4$--$D_7$.
\item \emph{The consumption list of Step 5 of Theorem \ref{thm:charp}.}  The
transitive closure of every \verb|\ref| and \verb|\eqref| reachable from
\S\ref{sec:induction}, expanding each statement's own text and proof and
stopping only at the references \S\ref{sec:induction} makes in order to
\emph{disclaim} them (Lemma \ref{lem:node}, Theorem \ref{thm:B} and
Proposition \ref{prop:W2}) and at the two it makes to statements it is
conditional on rather than consumes (Remark \ref{rem:nodefree} and Theorem
\ref{thm:charp}), was computed
from the source of this paper: $42$ labels, $33$ of them external to
\S\ref{sec:induction}, and every one of the $33$ is named in the list of Step 5
--- the single exception being the display \eqref{eq:Iddeg}, which sits inside
the proof of the named Lemma \ref{lem:Id}.
\item \emph{The cosmetic points of \S\ref{sec:profile} and \S\ref{ss:cbproof}.}
The failure of Lemma \ref{lem:nodefree}(a) without its top-truncation
hypothesis (the line in the bottom slot at degree $t-k$ has $v'_t=1$ and
threshold $t+k$, not $t+k-2$) was checked at all $703$ pairs $(n,k)$ with
$n\le40$; the comparison of Remark \ref{rem:nodefree} at all $19684$ greedy
truncations in that range ($9139$ even $v'_t>0$ agreeing, $9842$ odd differing
by $2$, $703$ with $v'_t=0$ where the lowest taken degree does not exist); the
sign of Lemma \ref{lem:cbseq}(ii), $\psi'_a-\psi_a=-\al(\varphi)_a$, on $2000$
random instances over $\mathbb{F}_7$ (the opposite sign fails in $1958$ of
them); the exception in Remark \ref{rem:mu_a} on all $210$ members of the
family $(1,0,0)$ over $D_4$--$D_{12}$ (descent set exactly $\{t-1\}$ except at
the $9$ instances $k=n-2$, $h=k-1$, where $\mu=\la$ has no descent); and the
claim of \S\ref{sec:setup} that $\varsigma(i,j)=\operatorname{sign}(j-i)$ is
induced by exactly one orientation, namely ``every arrow points to the smaller
label'', by enumerating all $2^{\#\text{edges}}$ orientations of
$D_4$--$D_{10}$.
\item \emph{The $\Ext^1$ formula and the four-term sequence}
\textup{(}\S\ref{ss:cbproof}\textup{)}.  The complex $C^0\to C^1\to C^2$ of
Lemma \ref{lem:cbseq} was built literally as written and tested on $1641$
pairs $(M,N)$ of genuine $\Pia$-modules over $D_4$--$D_6$ at
$p\in\{2,3,5\}$ --- all simples $S_i$, all injectives $I_\ka$, and submodules
$V(\mu)\subseteq I_\ka$: $\beta\al=0$ in every case; $\dim\operatorname{coker}
\beta=\dim\Hom_\Pia(N,M)$ in every case; $E(\psi)$ satisfies the preprojective
relations for every $\psi\in\ker\beta$ tested ($117$ cases) and for no $\psi$
with $\beta(\psi)\ne0$ ($108$ cases).  The value
$\dim(\ker\beta/\im\al)$ was compared against $\Ext^1$ computed
\emph{independently of the sequence} in every case where it is known a priori
--- $\Ext^1(S_i,S_j)=[\,i\sim j\,]$, and $\Ext^1(M,I_\ka)=\Ext^1(I_\ka,N)=0$
by injectivity and projectivity of $I_\ka$ --- with agreement at all $1641$
pairs, and Proposition \ref{prop:ext1} agrees with both.
\item \emph{The node column of a top truncation} \textup{(}Lemma
\ref{lem:nodefree}\textup{)}.  Part \textup{(a)} was checked on all $19684$
greedy top truncations of the profile $(1,2^{k-1},1)$ with $v'_t<2k$ over
$D_4$--$D_{40}$, all $2\le k\le n-2$: the threshold $D_t(X)$ of Definition
\ref{def:upset} matches the closed form, and $X_t[D]\ne0$ iff $v'_t$ is even
and $\ge2$, without exception.  The same run confirms the discrepancy recorded
in Remark \ref{rem:nodefree}: the lowest \emph{taken} degree of Lemma
\ref{lem:node} agrees with $D_t(X)$ at all $9139$ even $v'_t>0$ and exceeds it
by $2$ at all $9842$ odd $v'_t$; and at $v'_t=2k$ the column is full, so
$D_t(X)=-\infty$, which is why part \textup{(a)} carries the hypothesis
$v'_t<2k$.  Part \textup{(b)} was checked on the full Weyl orbits of
$D_4$--$D_{10}$, all $k$: of the $77196$ extremal weights, the $26181$ with
$c_t\ge1$ all have $v_t\le2k-1$, and every weight with $v_t=2k$ has
$c_t\le0$, without exception.
\item \emph{The count of Proposition \ref{prop:reps}.}  The
$W_{J'}$-minimal $\mu\in W\om_k$ with $c_t(\mu)\ge1$ were enumerated directly
from the description $W\om_k=\{\sum_{i\in S}\eta_ie_i\}$ over $D_4$--$D_{12}$,
all $2\le k\le n-2$: the count is $9(k-1)$ in every case, \emph{including}
$k=n-2$, and $\la$ itself lies in the list exactly when $k=n-2$, where it
satisfies $c_t(\la)=1$.  In particular $(n,k)=(4,2),(6,4),(7,5)$ give $9$,
$27$, $36$.
\item \emph{Representative-level statements} --- Proposition \ref{prop:reps}
(the classification and the count $9(k-1)$), Lemma \ref{lem:vclosed},
\eqref{eq:sform}, Propositions \ref{prop:incr}, \ref{prop:decr}, Lemma
\ref{lem:node}, Proposition \ref{prop:PS}, Propositions \ref{prop:L1},
\ref{prop:L3}, \ref{prop:L4}, Proposition \ref{prop:leafclasses}, and the
complete obligation census of Theorem \ref{thm:B} --- checked in closed form
over $D_6$--$D_{30}$, every $2\le k\le n-2$ (so up to $k=28$), every family and
every $h$.  In every fatal obligation class (full$\to$empty, line$\to$empty at
a leaf, line$\to$partial at a leaf) the count is $0$ throughout.
\item \emph{The degenerate node column} \textup{(}the last paragraph of the
proof of Proposition \ref{prop:partials}\textup{)}.  The closed form
$v_t=2k-2h-q-x$ of Lemma \ref{lem:vclosed} was checked against $\la-\mu$
computed directly from the Bourbaki coordinates \eqref{eq:reps}, at every one
of the $18\,216$ representatives over $D_4$--$D_{25}$, all $2\le k\le n-2$, all
nine families and every admissible $h$: the chain, node and leaf entries of
\eqref{eq:vclosed}--\eqref{eq:vleaf} agree without exception, no entry of $v$
is negative, and $v_t=0$ at exactly $253$ of them --- precisely the family
$(1,0,0)$ at $h=k-1$, one for each pair $(n,k)$.  At those weights the module
$V(\mu)$ was built from the socle construction over $\mathbb{F}_3$,
$\mathbb{F}_5$, $\mathbb{F}_7$ and $\mathbb{Q}$ for $D_4$--$D_{10}$ and over
$\mathbb{F}_3$ for $D_{11}$--$D_{13}$ ($139$ instances, covering all $66$ pairs
$(n,k)$ with $n\le13$): in every case $\dimv V(\mu)=\la-\mu$, the colour-$t$
column is zero in every degree, $V(\mu)$ is an up-set at \emph{every} vertex
--- in particular at $t-1$, which is $(*^\circ)$ --- and its profile is the
greedy top truncation at every vertex, which is $(B)$.
\item \emph{The obligation census behind Proposition \ref{prop:L5}} --- over
$D_4$--$D_{44}$, every $2\le k\le n-2$, all nine families and every $h$, some
$4.8\times10^{7}$ obligations were classified.  The class ``full$\to$partial
along a chain arrow'' is populated only in the family $(1,0,0)$; in each of the
four pinned families it is empty, and every chain obligation into a partial
slot has a line source.  Independently, Proposition \ref{prop:L5}(a) was
confirmed on all $24682$ instances of $(0,-1,\pm1)$ in that range
($E_{t-1}=E_t+1$, with $(E_t,t)$ the unique partial slot), and
Proposition \ref{prop:L5}(b) on all $22960$ instances of $(1,-1,\pm1)$.  This
is a check; the proof is in \S\ref{ss:L5}.
\item \emph{The interval facts and the rank statements} --- Lemma
\ref{lem:intervals} (I1$'$)--(I5) in closed form up to $D_{200}$; the
statements (T0)--(T2) of Proposition \ref{prop:T012} in closed form up to
$D_{100}$, every $k$, every chain edge and every degree ($2.4\times10^{7}$
arrow/degree pairs, no failure).
\item \emph{The profile of $I_\ka$ and the spin weights} --- the half-line
construction of Proposition \ref{prop:foldsol}(a) (the recursion, the mirror
condition \eqref{eq:mirror}, the symmetry \eqref{eq:Bsym}, non-negativity,
vanishing above $2t$ and the total $k(2n-1-k)$) checked over $D_4$--$D_{80}$
for every $1\le k\le n-2$; the multiplicity-freeness of Corollary
\ref{cor:mfree} in closed form over $D_4$--$D_{150}$ for all three of
$\ka\in\{1,n-1,n\}$, and against the actual modules $I_\ka$ over $D_4$--$D_8$,
where the top degrees $2t$ and $2(n-1)-2$ and the dimensions $2n-2$ and
$n(n-1)/2$ come out as predicted.
\item \emph{Extremal modules built explicitly} --- Theorem \ref{thm:main},
Proposition \ref{prop:T}, $(*^\circ)$ and $(B)$ verified on the actual modules,
over $\mathbb{F}_p$ for $p=3,5,7,11,13$, for all $D_9$--$D_{11}$ with $k\le7$.
\item \emph{All graded submodules enumerated} --- Theorem \ref{thm:staircase},
Lemma \ref{lem:ab} and Corollary \ref{cor:F} checked against every graded
submodule of $I_\ka$, for $D_7$ with $k\le4$ ($1672$, $1783$ and $2011$
submodules over $\mathbb{F}_2$, $\mathbb{F}_3$, $\mathbb{F}_5$ respectively ---
the count is field-dependent), and for $D_7$ and $D_8$ with all $k$.  The
preservation of $(H)$ holds on every \emph{extremal} graded submodule in that
range, and only there: on arbitrary graded submodules it fails, first at $D_6$,
$\ka=4$, where $v(N)=(1,2,3,4,3,2)$ is a staircase whose colour-$4$ socle step
gives $(1,2,3,5,3,2)$, which is not.  This is the same $\ka\ge4$ phenomenon as (N3)
below, and extremality is supplied at every point of use by Corollary
\ref{cor:prefix}.
\item The criterion of Theorem \ref{thm:typeD} itself was compared against a
direct solver on the full Weyl orbit for $D_4$--$D_{11}$ at
$p\in\{2,3,5\}$ ($241008$ weights each) and for $D_4$--$D_{10}$ over
$\mathbb{Q}$ and at $p=7$ ($77196$ each), with agreement at every element;
and beyond the enumerable orbits, on the complete deterministic band $m\le2$ ---
which contains the whole region $r+2m\le2$ together with the adjacent shell
$r+2m\in\{3,4\}$, hence the entire boundary of the criterion --- for $D_{12}$
and $D_{13}$ ($156663$ and $370221$ weights at $p=2$), and the band $m\le1$ for
$D_{14}$ ($257339$).  Agreement is independent of the reduced word (checked on
a second word for $1685$ weights of $D_{11}$).  All of this was produced by a
from-scratch implementation sharing no code with the one used elsewhere in this
project; the two agree on $7020$ graded dimensions and $11880$ arrow ranks over
$D_4$--$D_{12}$.  Separately, Lemma \ref{lem:U} --- which \S\ref{sec:charp}
does \emph{not} transport, and which retains its hypothesis
$\operatorname{char}\ne2$ --- was checked at all $9139$ multiplicity-$2$ node
slots over $D_4$--$D_{40}$, all $k$, at $p\in\{2,3,5,7,11,13\}$.
\end{itemize}

\subsection{Negative results worth recording}\label{ss:negative}
These delimit the argument and explain why several shorter routes do not exist.

\begin{enumerate}
\item[\textbf{(N1)}] \emph{Full commutativity is neither necessary nor
sufficient.}  In $D_4$, $w=s_2s_4s_3s_1s_2$ is fully commutative and does not
have simple spectrum (Example \ref{ex:fcbad}); \cite[Ex.~5.12]{DKK} gives a $w$
with simple spectrum that is not fully commutative.
\item[\textbf{(N2)}] \emph{The greedy top truncation at every vertex is false
for general extremal $\mu$.}  It holds at the representatives of Proposition
\ref{prop:reps} --- that is $(B)$, Theorem \ref{thm:B} --- but not in general.
The smallest failure is $D_5$, $\la=\om_2$, $\mu=(1,-1,0,0,0)$, where $V_2$
occupies degrees $2$ and $6$ while the top truncation would give $6$ and $4$;
here $t-1=2$ is the chain neighbour of the node.  That element has $r+2m=2$ and
so \emph{does} have simple spectrum, so the failure is not confined to
degenerate cases.  It is precisely why the theorem is proved at $t$ only, and
why the reduction to a bounded list of weights is necessary.
\item[\textbf{(N3)}] \emph{Extremality cannot be dropped from $(*^\circ)$.}  The
localised statement --- for every graded $N\subseteq I_\ka$, $(H)$ together with
$\langle\al_t^\vee,\la-\dimv N\rangle\ge1$ implies $N$ is an up-set at $t-1$ ---
holds for every graded submodule when $\ka\le3$ and fails at $\ka=4$, where
$t-1$ first carries multiplicity $2$ in two different degrees.  The smallest
counterexample is $D_6$, $\ka=4$, $t=4$, with graded profile
($\dim N/\dim I_\ka$ by degree)
\[
\begin{array}{ll}
 1:\ 0/1\ \ 1/1 &
 4:\ 0/1\ \ 0/2\ \ 1/2\ \ 2/2\ \ 1/1\\
 2:\ 0/1\ \ 1/2\ \ 1/1 &
 5:\ 0/1\ \ 1/1\ \ 1/1\ \ 1/1\\
 3:\ 0/1\ \ 1/2\ \ 1/2\ \ 1/1 &
 6:\ 0/1\ \ 0/1\ \ 1/1\ \ 1/1
\end{array}
\]
which satisfies $(H)$, has $\langle\al_t^\vee,\la-\dimv N\rangle=1$, and is
partial at two degrees at $t-1$.  Likewise Claim F (Corollary \ref{cor:F}) is
false for arbitrary graded submodules.
\item[\textbf{(N4)}] \emph{The staircase property needs extremality.}  Graded
submodules of $I_\ka$ that are not staircases exist, and in the enumeration
described in \S\ref{sec:verification} every one of them has
$\delta(\dimv M)>0$, in accordance with Theorem \ref{thm:staircase}.  We quote
no count here: the number of graded submodules depends on the ground field
(a two-dimensional slot has $|\mathbb{F}_p|+1$ lines in it), so a bare count
is not a well-posed statistic.
\item[\textbf{(N5)}] \emph{The larger parabolic does not work.}  With
$J=I\setminus\{t-1,t\}$, which contains the two leaves, the side condition
$\langle\al_t^\vee,\mu\rangle\ge1$ is not $W_J$-invariant, and the
$W_J$-minimal representative of the orbit of a $\mu$ satisfying it usually does
not satisfy it: over $D_4$--$D_7$, all $2\le k\le n-2$, there are $764$ weights
$\mu\in W\om_k$ with $c_t(\mu)\ge1$, and $648$ of them have $W_J$-minimal
representative outside the hypothesis set; over $D_4$--$D_8$ the counts are
$2565$ and $2207$.  Since the
$W_J$-minimal element of an orbit is unique, there is no alternative descent
path, and $(*^\circ)$ is false without its side condition --- the $64$
respectively $244$ extremal non-up-sets recorded in \S\ref{ss:checks} are
exactly the counterexamples.  This is why
Proposition \ref{prop:reps} uses $J'=\{1,\dots,t-2\}$ and produces $9(k-1)$
representatives rather than $k$.
\item[\textbf{(N6)}] \emph{The canonical truncation is not always the right
module.}  In the four families with $q+x$ even and $z\ne0$, the leaf pin forces
the line at $(E_t,t)$ to be $\im(a_{l\to t})$, which by Proposition
\ref{prop:W4l} is \emph{not} a $\tau$-eigenvector.  So $T(v)$ fails to be a
submodule there and must be replaced by $T^{\,l}(v)$ (Theorem \ref{thm:B}).
\item[\textbf{(N7)}] \emph{Containment of multiplicity-$2$ windows does not
prove Theorem \ref{thm:lemA}}; see Remark \ref{rem:nowindow}.  Nor is the naive
mechanism for (G1) available: when $X_t[D]$ is a line it is always one of the
two special lines, so one leaf arrow annihilates it and that leaf's fullness at
$D+1$ must come from elsewhere in $X$.
\end{enumerate}

\begin{thebibliography}{99}

\bibitem{BBK} S.~Brenner, M.~C.~R.~Butler and A.~D.~King, \emph{Periodic
algebras which are almost Koszul}, Algebr.\ Represent.\ Theory \textbf{5}
(2002), 331--368.

\bibitem{BK} P.~Baumann and J.~Kamnitzer, \emph{Preprojective algebras and MV
polytopes}, Represent.\ Theory \textbf{16} (2012), 152--188.

\bibitem{BKT} P.~Baumann, J.~Kamnitzer and P.~Tingley, \emph{Affine
Mirkovi\'c--Vilonen polytopes}, Publ.\ Math.\ Inst.\ Hautes \'Etudes Sci.\
\textbf{120} (2014), 113--205.

\bibitem{CB} W.~Crawley-Boevey, \emph{On the exceptional fibres of Kleinian
singularities}, Amer.\ J.\ Math.\ \textbf{122} (2000), 1027--1037.

\bibitem{DKK} H.~Dinkins, I.~Karpov and V.~Krylov, \emph{The quantum Hikita
conjecture via quasimaps}, arXiv:2608.16746.

\bibitem{DEKM} A.~Dranowski, B.~Elek, J.~Kamnitzer and C.~Morton-Ferguson,
\emph{Heaps, crystals, and preprojective algebra modules}, arXiv:2202.02490.

\bibitem{GLS} C.~Geiss, B.~Leclerc and J.~Schr\"oer, \emph{Kac--Moody groups
and cluster algebras}, Adv.\ Math.\ \textbf{228} (2011), 329--433.

\bibitem{Nak} H.~Nakajima, \emph{Quiver varieties and Kac--Moody algebras},
Duke Math.\ J.\ \textbf{91} (1998), 515--560.

\bibitem{ST} A.~Savage and P.~Tingley, \emph{Quiver Grassmannians, quiver
varieties and the preprojective algebra}, Pacific J.\ Math.\ \textbf{251}
(2011), 393--429.

\bibitem{Stem} J.~Stembridge, \emph{Minuscule elements of Weyl groups}, J.\
Algebra \textbf{235} (2001), 722--743.

\end{thebibliography}

\end{document}
